%% =====================================================================
%%  A corrected exposition of Zhuk's simplified CSP dichotomy proof.
%%
%%  GENERATED FILE -- do not edit.  Produced by make_standalone.py from
%%  csp.tex and the per-section sources; edit those and regenerate.
%%  Generated from revision 35c8e30 on 2026-08-02.
%% =====================================================================

%% ---------------------------------------------------------------- csp.tex
\documentclass[11pt]{article}
%% ---------------------------------------------------------------- preamble.tex
% Shared preamble for the CSP dichotomy blueprint.
% House style follows the earlier `zhuk_centers` blueprint: numbered theorem
% environments shared across a single counter, a `longtheorem` upright variant
% for statements with many numbered hypotheses, and explicit `convention`
% environments for standing hypotheses.
\usepackage[a4paper,margin=27mm,headheight=14pt]{geometry}
\usepackage[T1]{fontenc}
\usepackage{lmodern}
\usepackage{amsmath,amssymb,amsthm}
\usepackage{longtable,array}
\usepackage{fancyhdr}
\usepackage[hidelinks]{hyperref}

\newtheoremstyle{mainplain}{9pt}{9pt}{\itshape}{}{\bfseries}{.}{0.55em}{}
\newtheoremstyle{mainroman}{9pt}{9pt}{\normalfont}{}{\bfseries}{.}{0.55em}{}
\theoremstyle{mainplain}
\newtheorem{theorem}{Theorem}[section]
\newtheorem{lemma}[theorem]{Lemma}
\newtheorem{proposition}[theorem]{Proposition}
\newtheorem{corollary}[theorem]{Corollary}
\theoremstyle{mainroman}
\newtheorem{longtheorem}[theorem]{Theorem}
\newtheorem{longlemma}[theorem]{Lemma}
\newtheorem{longproposition}[theorem]{Proposition}
\newtheorem{longcorollary}[theorem]{Corollary}
\theoremstyle{definition}
\newtheorem{definition}[theorem]{Definition}
\newtheorem{remark}[theorem]{Remark}
\newtheorem{convention}[theorem]{Convention}
\newtheorem{hazard}[theorem]{Formalization note}

\setlength{\parindent}{1.25em}\setlength{\parskip}{0pt}\setlength{\emergencystretch}{2em}
\newcommand{\alphlabels}{\renewcommand{\labelenumi}{\textup{(\alph{enumi})}}%
  \renewcommand{\theenumi}{\alph{enumi}}}
\newcommand{\romanlabels}{\renewcommand{\labelenumi}{\textup{(\roman{enumi})}}%
  \renewcommand{\theenumi}{\roman{enumi}}}
\providecommand{\tightlist}{\setlength{\itemsep}{0pt}\setlength{\parskip}{0pt}}
\numberwithin{equation}{section}
\setcounter{tocdepth}{2}\setcounter{secnumdepth}{3}
\pagestyle{fancy}\fancyhf{}
\fancyhead[L]{\small CSP dichotomy}\fancyhead[R]{\small\nouppercase{\leftmark}}
\fancyfoot[C]{\thepage}\renewcommand{\headrulewidth}{0.3pt}

% ---------------------------------------------------------------- operators
\DeclareMathOperator{\Sg}{Sg}
\DeclareMathOperator{\Clo}{Clo}
\DeclareMathOperator{\Pol}{Pol}
\DeclareMathOperator{\Inv}{Inv}
\DeclareMathOperator{\Var}{Var}
\DeclareMathOperator{\Con}{Con}
\DeclareMathOperator{\Sol}{Sol}
\DeclareMathOperator{\Expanded}{Exp}
\DeclareMathOperator{\CSP}{CSP}
\DeclareMathOperator{\GenSAT}{SAT}
\DeclareMathOperator*{\proj}{pr}
\newcommand{\ConOne}{\operatorname{Con}_1}

% ---------------------------------------------------------------- algebras
\newcommand{\bA}{\mathbf A}
\newcommand{\bB}{\mathbf B}
\newcommand{\bC}{\mathbf C}
\newcommand{\bD}{\mathbf D}
\newcommand{\bE}{\mathbf E}
\newcommand{\bR}{\mathbf R}
\newcommand{\bS}{\mathbf S}
\newcommand{\bZ}{\mathbf Z}
\newcommand{\Vn}{\mathcal V_{n}}

% ---------------------------------------------------------- subuniverse types
% The six types of Zhuk's theory.  Typed as upright small capitals so that a
% type never reads as a product of variables.
\newcommand{\TBA}{\textup{\textsc{ba}}}
\newcommand{\TC}{\textup{\textsc{c}}}
\newcommand{\TS}{\textup{\textsc{s}}}
\newcommand{\TD}{\textup{\textsc{d}}}
\newcommand{\TL}{\textup{\textsc{l}}}
\newcommand{\TPC}{\textup{\textsc{pc}}}
\newcommand{\MT}[1]{\mathcal M#1}
\newcommand{\TML}{\MT\TL}
\newcommand{\TMPC}{\MT\TPC}
\newcommand{\TMD}{\MT\TD}

% `C <_{T}^{A} B' and its relatives.
\newcommand{\subt}[2][]{<_{#2}^{#1}}
\newcommand{\subte}[2][]{\le_{#2}^{#1}}
\newcommand{\dsubt}[2][]{\dot<_{#2}^{#1}}
\newcommand{\dsubte}[2][]{\dot\le_{#2}^{#1}}
\renewcommand{\lll}{\mathrel{\lhd\!\lhd\!\lhd}}
\newcommand{\dlll}{\mathrel{\dot{\lhd}\!\dot{\lhd}\!\dot{\lhd}}}
\newcommand{\sdle}{\le_{\mathrm{sd}}}

% congruence cover  sigma^*
\newcommand{\cover}[1]{#1^{*}}
% bridge collapse
\newcommand{\bcol}[1]{\widetilde{#1}}
% absorption
\newcommand{\ab}{\lhd}
\newcommand{\abin}{\lhd_{\mathrm{bin}}}
\newcommand{\abc}{\lhd_{\mathrm{c}}}

\hypersetup{pdftitle={The CSP Dichotomy Theorem: a formalization blueprint},
  pdfsubject={Zhuk's simplified proof, rewritten at formalization granularity},
  pdfauthor={}}

\title{The CSP Dichotomy Theorem\\[2pt]
  \large A formalization blueprint}
\author{}\date{}

\begin{document}
\maketitle\vspace{-2.2em}
\begin{abstract}
The Constraint Satisfaction Problem over a finite constraint language is
solvable in polynomial time if the language admits a weak near-unanimity
polymorphism, and is \textup{NP}-complete otherwise. This document rewrites
Zhuk's simplified proof \cite{ZhukSimplified} --- together with the
prerequisites it imports --- as a blueprint for a formalization in Lean~4 and
Mathlib. It states precisely what a formalization can and cannot claim
(\S\ref{sec:target}); fixes the vocabulary at a granularity a proof assistant
can consume (\S\S\ref{sec:conventions}--\ref{sec:instances}); records the
theory's interface (\S\ref{sec:properties}) and its summit
(\S\ref{sec:main-statements}); states the algorithm and the hard half
(\S\S\ref{sec:algorithm}--\ref{sec:hardness}); and reports what Mathlib is
missing, what has been built, and what the remaining work is
(\S\S\ref{sec:complexity-layer}--\ref{sec:formalization}).

Throughout, \emph{formalization notes} mark the places where the informal proof
has latitude that a formal one does not: a suppressed side condition, a
quantifier whose scope is ambiguous, an object asserted to exist without an
argument, or a step whose well-foundedness is left to the reader.
\end{abstract}

\tableofcontents\clearpage

\section*{Introduction}
\addcontentsline{toc}{section}{Introduction}

The dichotomy conjecture of Feder and Vardi \cite{FederVardi} was settled
independently by Bulatov \cite{Bulatov} and Zhuk \cite{ZhukJACM} in 2017. Both
proofs are long. Zhuk's, in the form published in 2020, is seventy-eight
journal pages, and its central induction alternates between global and local
information in a way its author later described as forcing ``a very complicated
induction to prove most of the claims''.

In 2024 Zhuk published a simplified proof \cite{ZhukSimplified}. The
simplification is a redesign of the theory of strong subalgebras so that
\emph{every reduction is either strong or global}: for each domain $D_{x}$ one
can build a chain
$D_{x}\supsetneq D_{x}^{(1)}\supsetneq\dots\supsetneq\{a\}$ in which every step
is either a strong subset or the intersection with a block of an equivalence
relation with very strong properties. The types of the intermediate sets need
not be tracked --- only that such a chain exists, written $C\lll B$. This is
what makes the argument tractable, and it is what makes it a plausible
formalization target.

\medskip

\noindent\textbf{Why this route.} Among the available proofs,
\cite{ZhukSimplified} is the only complete, recent, rigorously written one whose
author has already removed the obstruction that made its predecessor
unformalizable. Bulatov's proof requires tame congruence theory and centralizer
theory, neither of which is formalized anywhere. Restricting to minimal Taylor
algebras --- which \cite{ZhukSimplified} itself suggests would simplify its
\S2 --- turns out to be a bad trade: it would simplify perhaps a quarter of one
section and none of the CSP argument, at the cost of importing the cyclic term
theorem, which the present route does not otherwise need. Brady's notes
\cite{BradyNotes} are the clearest prose in the subject and are the best source
for several of the imported prerequisites, but they stop short of the
dichotomy. The route taken here is \cite{ZhukSimplified} as the spine, with
\cite{BradyNotes} for the imports, on top of an existing formalization of Zhuk's
centre theorem.

\medskip

\noindent\textbf{What is claimed.} Not that the dichotomy has been formalized.
\S\ref{sec:scope} estimates the transitive closure of what must be proved at
$100$--$130$ printed pages, which extrapolates to tens of thousands of lines of
Lean, none of it attempted in any proof assistant. That is a large number but
not, at current throughput, a long time; what governs the schedule is the
critical path and the defects of \S\ref{sec:defects}, not the line count. What is offered is its design: a route, precise statements, a
vocabulary pinned down where the sources leave it loose, an inventory of the
missing Mathlib layers, and a Lean development whose foundational modules are
complete and whose target theorems are stated.

\medskip

\noindent\textbf{What the formalization notes are for.} A blueprint's value is
mostly in its failures of nerve --- the places where writing the statement down
carefully shows that it does not quite say what it appeared to say. The
predecessor project found four such places in twenty-four pages, and all four
were in the foundations rather than in the mathematics. The notes here record
the same kind of thing: that $\mathbf Z_{p}$ lies in $\Vn$ only when
$p\mid n-1$ (Hazard~\ref{haz:Zp}); that $\cover\sigma$ need not be a congruence
(Hazard~\ref{haz:irreducible}); that images do not commute with intersection,
and that one proof needs them to (Hazard~\ref{haz:quotient-rel}); that ``we
cannot weaken forever'' needs a measure that is not the obvious one
(Hazard~\ref{haz:weakening}); that the main induction quantifies over the
instance \emph{inside} (Hazard~\ref{haz:the-induction}). \S\ref{sec:defects} collects the
harder cases separately: ten places where the source cannot be transcribed at
all, of which one --- \cite{ZhukSimplified}'s restatement of a lemma from
\cite{ZhukStrong}, with two hypotheses dropped --- is false as printed, with a
two-line counterexample.

\clearpage
%% ---------------------------------------------------------------- ch0-target.tex
\section{What is being proved}\label{sec:target}

\subsection{The theorem}

Fix a finite set $A$ and a finite set $\Gamma$ of relations on $A$, the
\emph{constraint language}. An instance of $\CSP(\Gamma)$ is a conjunction
$R_{1}(\bar x_{1})\wedge\dots\wedge R_{s}(\bar x_{s})$ with each
$R_{i}\in\Gamma$; the question is whether it is satisfiable.

\begin{theorem}[Dichotomy; Bulatov, Zhuk]\label{thm:dichotomy}
If some weak near-unanimity operation preserves $\Gamma$, then $\CSP(\Gamma)$ is
solvable in polynomial time. Otherwise $\CSP(\Gamma)$ is \textup{NP}-complete.
\end{theorem}

Both halves are conditional on a model of computation, and this is not a
formality that a formalization may wave through. ``Solvable in polynomial time''
quantifies over an encoding of instances as strings, a machine, and a
polynomial bound on its step count; ``\textup{NP}-complete'' additionally
quantifies over all problems in \textup{NP} and over polynomial-time many-one
reductions. Mathlib supplies Turing machines
(\texttt{Mathlib/Computability/TuringMachine/}) and a notion of computability,
but it has \emph{no} complexity class, no polynomial-time predicate, and no
notion of polynomial-time reduction.

\subsection{Where the work actually is}\label{sec:program-axis}

An earlier draft of this section concluded from that gap that the
complexity-theoretic half should be deferred because building it would dominate
the project. That inference was wrong, and it is worth recording why, because
the corrected version is what organises the rest of the document.

The complexity vocabulary itself --- languages, a cost model, $\mathrm P$,
$\mathrm{NP}$ by verifiers, $\le_{p}$, \textup{NP}-hardness,
\textup{NP}-completeness, and a generic \textup{NP}-hard problem --- has since
been built, in about $850$ lines. That is consistent with the two existing
Cook--Levin formalizations, whose class layers are $702$ lines (Isabelle) and
$992$ lines (Coq) against totals of $56\,514$ and $\approx 16\,500$. Defining
the classes is cheap in every development that has done it.

What is expensive is not complexity theory. It is any statement that requires
one to \emph{exhibit a program}. In a formalization, ``$f$ is computable in
polynomial time'' is not a property one observes of a mathematical function; it
is an existential over programs, discharged only by writing one and proving a
bound on its step count. The dividing line runs along that axis, and it cuts
across the algebra/complexity split rather than agreeing with it:

\begin{center}
\begin{tabular}{@{}lll@{}}
\hline
Target & Exhibit a program? & Cost \\
\hline
\textup{(T0)} the algebraic core & no & the mathematics \\
\textup{(T1)} $\textsc{Solve}$ is correct & no (but see Hazard~\ref{haz:t1-computable}) & moderate \\
\textup{(T2)} $\textsc{Solve}$ is polynomial & \textbf{yes --- the largest one here} & \textbf{large} \\
\textup{(H0)} $\Gamma$ pp-interprets \textsc{Nae} & no & the mathematics \\
\textup{(H1)} pp-interpretation $\Rightarrow\le_{p}$ & yes, a syntactic substitution & small \\
\hline
\end{tabular}
\end{center}

So the document still separates the machine-free statements from the rest, but
for a sharper reason than before, and with a different verdict on \textup{(T2)}.

\subsection{The layered targets}

Write $\mathbf A=(A;w)$ for a finite idempotent algebra whose basic operation
$w$ is a special weak near-unanimity operation, and $\Vn$ for the class of such
algebras with $w$ of arity $n$ (Convention~\ref{conv:standing}).

\paragraph{Tractable half.} Three layers.

\begin{description}\tightlist
\item[\textup{(T0)} \emph{The algebraic core.}] The three theorems of
  \S\ref{sec:main-statements} --- the existence of a strong subuniverse or a
  linear congruence, the safety of a reduction to a strong subuniverse, and the
  codimension-one theorem --- together with everything they rest on. These are
  statements about finite algebras, congruences, and CSP instances regarded as
  finite combinatorial objects. No computation appears in them.
\item[\textup{(T1)} \emph{Correctness of the algorithm.}] Zhuk's algorithm
  written as a total function $\textsc{Solve}$ from instances to Booleans, and
  the proposition
  \[
    \textsc{Solve}(\mathcal I)=\texttt{true}\iff
    \mathcal I\ \text{has a solution}.
  \]
  This is a statement about a recursive function, and it is provable from
  \textup{(T0)} with no cost model. It is the honest formal content of ``the
  algorithm works'' --- \emph{provided} $\textsc{Solve}$ actually computes;
  see Hazard~\ref{haz:t1-computable}.
\item[\textup{(T2)} \emph{The complexity wrapper.}] That $\textsc{Solve}$ runs
  in polynomial time in the size of the instance, for fixed $\Gamma$. This is
  the one place a cost model is needed on this side, and it is a large
  undertaking rather than a wrapper: see Hazard~\ref{haz:t2-is-not-small}.
\end{description}

\paragraph{Hard half.} Two layers.

\begin{description}\tightlist
\item[\textup{(H0)} \emph{The algebraic core.}] If no weak near-unanimity
  operation preserves $\Gamma$ then $\Gamma$ primitively positively interprets
  every finite constraint language --- in particular one whose CSP is
  \textup{NP}-hard. This is again pure algebra: the Galois connection between
  polymorphisms and primitive positive definability, the Maróti--McKenzie
  theorem that a Taylor term yields a weak near-unanimity term, and the
  Bulatov--Jeavons--Krokhin reduction.
\item[\textup{(H1)} \emph{The complexity wrapper.}] That a primitive positive
  interpretation yields a polynomial-time many-one reduction, and hence
  \textup{NP}-hardness. This does need a program, but the program performs a
  syntactic substitution of formulas whose size is linear in the instance, so
  it is genuinely small.
\end{description}

The claim this document makes is that \textup{(T0)}, \textup{(T1)} and
\textup{(H0)} are the mathematics and are what the literature actually proves.
\textup{(H1)} is a small addition on top. \textup{(T2)} is neither small nor
mathematics, but it is not \emph{unrelated} to the dichotomy either --- it is
the assertion that the algorithm one has just proved correct is also fast, and
it is the second-largest component of the project. The blueprint is written for
\textup{(T0)}, \textup{(T1)} and \textup{(H0)};
\S\ref{sec:complexity-layer} states \textup{(T2)} and \textup{(H1)} precisely
and says what building them costs.

\begin{hazard}[\textup{(T1)} is vacuous unless $\textsc{Solve}$ computes]
\label{haz:t1-computable}
``$\textsc{Solve}$ written as a total function'' hides a requirement that no
cost model would catch. If $\textsc{Solve}$ is defined using classical choice
--- for instance by deciding ``does $D_{x}$ have a nonempty proper binary
absorbing subuniverse?'' with \texttt{Classical.dec} --- then
$\textsc{Solve}(\mathcal I)=\texttt{true}\iff\mathcal I$ has a solution is a
true theorem that says nothing algorithmic, because $\textsc{Solve}$ does not
reduce. To mean what it appears to mean, \textup{(T1)} needs every predicate
the algorithm branches on to carry a genuine decision procedure.

That is a real obligation with a locatable cost. Absorption is defined by an
existential over \emph{terms} --- an infinite type --- so
``$C\subt{\TBA}B$'' is not decidable on its face. It becomes decidable through
the bound that a witness of arity at most $|A|^{|A|}$ suffices, which is a
theorem and not a definition. Hazard~\ref{haz:algorithm-formalization}(b)
records that the bound exists; what it does not record, and what matters here,
is that \textup{(T1)} is empty without it. The same applies to the searches
over congruences and over subsets.
\end{hazard}

\begin{hazard}[\textup{(T2)} is not a wrapper]\label{haz:t2-is-not-small}
Two things make \textup{(T2)} much larger than the phrase ``complexity
wrapper'' suggests.

\emph{The program is enormous.} \textup{(T2)} requires $\textsc{Solve}$ ---
\textsc{ForceConsistency}, the search for a strong subuniverse,
\textsc{SolveLinear}, and the mutual recursion between them --- to be expressed
in a cost model, with a proved polynomial bound. Every existing measurement
says this kind of work dominates: it is the $50\,000$ lines of Turing-machine
engineering in the Isabelle Cook--Levin development, against $702$ for the
classes themselves.

\emph{The statement is non-uniform.} Because $\Gamma$ is fixed, the polynomial
may depend on $\Gamma$, and the searches over subsets of $A$, over congruences
on $A$, and over term operations of bounded arity are all constant-time. But
``constant'' here means the search is compiled away into a program whose
\emph{size} depends on $|A|$. So $\textsc{Solve}$ is not one program: it is a
family $\textsc{Solve}_{\Gamma}$ with a polynomial $p_{\Gamma}$ for each
$\Gamma$, and \textup{(T2)} must be stated in that form. Attempting a bound
uniform in $\Gamma$ would be proving something else, and something false.
\end{hazard}

\begin{remark}
The separation is not a dodge. Layer \textup{(T1)} is strictly stronger than
``\,$\CSP(\Gamma)$ is decidable'', which is trivial for a finite template: it
says that one particular, highly non-obvious algorithm --- the one whose
existence is the content of the dichotomy --- is correct. Everything that makes
the dichotomy hard lives in \textup{(T0)}. Conversely, no part of
\textup{(T0)}--\textup{(T1)} becomes easier if a cost model is available.
\end{remark}
\clearpage
%% ---------------------------------------------------------------- ch1-conventions.tex
\section{Standing conventions}\label{sec:conventions}

Zhuk's papers carry their hypotheses in prose: ``in this paper we assume that
every algebra is a finite idempotent algebra having a WNU term operation''.
A formalization cannot carry a hypothesis in prose, and an informal reader
cannot audit which results actually consume which hypothesis. We therefore fix
the hypotheses as a labelled convention and record, for each one, where it is
consumed.

\begin{convention}[Standing hypotheses]\label{conv:standing}
Throughout, $n\ge 3$ is a fixed integer and:
\begin{enumerate}\alphlabels\tightlist
\item\label{conv:finite} every algebra is \emph{finite} and has a
  \emph{nonempty} universe;
\item\label{conv:onebasic} every algebra has exactly one basic operation, of
  arity $n$, written $w$;
\item\label{conv:idem} $w$ is \emph{idempotent}: $w(x,\dots,x)=x$;
\item\label{conv:wnu} $w$ is a \emph{weak near-unanimity} operation:
  $w(y,x,\dots,x)=w(x,y,x,\dots,x)=\dots=w(x,\dots,x,y)$ for all $x,y$;
\item\label{conv:special} $w$ is \emph{special}:
  $w(x,\dots,x,y)=w\bigl(x,\dots,x,w(x,\dots,x,y)\bigr)$ for all $x,y$.
\end{enumerate}
The class of algebras satisfying \ref{conv:finite}--\ref{conv:special} is
written $\Vn$. Since only finite algebras are considered, $\Vn$ is not a
variety.
\end{convention}

\begin{remark}[Why one basic operation]\label{rem:one-basic}
Restricting to a single basic operation is not a loss. The dichotomy hypothesis
is that \emph{some} weak near-unanimity operation preserves $\Gamma$; one may
then discard the original signature and work in the algebra
$(A;w)$ whose only basic operation is that one, because every relation of
$\Gamma$ is a subuniverse of a power of $(A;w)$ and that is all the argument
uses. Passing from a WNU $w$ to a \emph{special} one is
Lemma~\ref{lem:special-wnu}, and costs an increase of arity from $n$ to
$n^{n!}$.

The gain is large for a formalization: a term is a finite tree with $n$-ary
branching over a single symbol, so ``term operation'' needs no signature
bookkeeping, and $\Clo(\mathbf A)$ is the clone generated by one operation.
\end{remark}

\begin{hazard}[Where each hypothesis is consumed]\label{haz:hyp-audit}
\ref{conv:finite} is used in every induction on $|A|$ and in every
minimality argument (a minimal element of a nonempty family of subsets
exists); it cannot be dropped anywhere. \ref{conv:idem} is what makes
singletons subuniverses and makes $\Sg$ of a nonempty set nonempty.
\ref{conv:wnu} is consumed only through the results it implies --- the
existence of absorption-theoretic structure --- and never directly in
\S\S\ref{sec:instances}--\ref{sec:main-statements}. \ref{conv:special} is used
where a unary polynomial must be idempotent under iteration, chiefly in the
theory of $\Sg$-closures of two-element sets. A formalization should carry
\ref{conv:finite}--\ref{conv:special} as typeclass instances on a structure and
not as ambient hypotheses, so that the audit is mechanical.
\end{hazard}

\begin{convention}[The empty set]\label{conv:empty}
A \emph{subuniverse} of $\mathbf A$ is a subset closed under $w$; the empty set
is a subuniverse under this definition, and we allow it. A
\emph{subalgebra} is a nonempty subuniverse; $\mathbf B\le\mathbf A$ always
carries $B\neq\varnothing$.

The relations $\subt{T}$, $\subte{T}$ and $\lll$ of \S\ref{sec:types} are
defined only between \emph{nonempty} sets; in particular
$\varnothing\lll A$ never holds. Where a construction can produce the empty
set --- a projection of an intersection, most often --- the dotted variants
$\dsubt{T}$, $\dsubte{T}$, $\dlll$ are used, and they mean ``the stated relation
holds, or the left-hand side is empty''.
\end{convention}

\begin{hazard}[The dot is not decoration]\label{haz:dot}
Conflating $\lll$ with $\dlll$ is the single most common way to misread
\S\ref{sec:types}. Every propagation statement of \S\ref{sec:properties} that
projects or intersects produces a dotted conclusion, and every statement that
consumes a $\lll$ hypothesis needs the undotted one. In a formalization the two
should be distinct constants, with the dotted one \emph{defined} as
\texttt{C = }$\varnothing$\texttt{ }$\vee$\texttt{ C }$\lll$\texttt{ B}, so that the case split is forced
at every use site rather than being resolved by the reader's judgement.
\end{hazard}

\begin{convention}[Relations, and indices]\label{conv:relations}
A relation is a subset of a product $A_{1}\times\dots\times A_{m}$ of universes,
indexed by $[m]=\{1,\dots,m\}$; $\mathbf R\le\mathbf A_{1}\times\dots\times
\mathbf A_{m}$ means additionally that $R$ is a subuniverse of the product
algebra. We write $\proj_{i_{1},\dots,i_{s}}(R)$ for the projection onto the
listed coordinates and $\mathbf R\sdle\mathbf A_{1}\times\dots\times\mathbf
A_{m}$ when every $\proj_{i}(R)=A_{i}$.

Products are indexed by \emph{arbitrary finite index sets}, not only by initial
segments of $\mathbb N$. Several constructions in \S\ref{sec:main-statements}
form a product over the variable set of an instance, or over a set of
subuniverses, and reindexing those to $[m]$ by hand is exactly the kind of
bookkeeping that a formalization pays for repeatedly and that buys nothing.
\end{convention}

\subsection{Legislated readings}\label{sec:legislation}

Ten places in \cite{ZhukSimplified} admit more than one reading, and in each the
choice is invisible in prose but forced in a formalization. This subsection
fixes all ten. Two of them turn out not to be conventions at all: item
\ref{leg:Zp} is a theorem, proved here, and item \ref{leg:central} is a missing
lemma, proved in \S\ref{sec:properties}. Of the remaining eight, six are
decided by internal evidence --- one reading makes some statement of the source
true or some proof valid and the other does not, so the ``choice'' is really a
reconstruction --- and only items \ref{leg:dimension} and \ref{leg:linked} are
free.

\begin{enumerate}\romanlabels

\item\label{leg:Zp}\textbf{The algebra $\mathbf Z_{p}$.}
\cite{ZhukSimplified} declares that ``every algebra $\mathbf Z_{p}$ belongs to
$\Vn$ for a fixed $n$'', with basic operation $x_{1}+\dots+x_{n}\bmod p$. That
operation is idempotent only when $p\mid n-1$, so the declaration is a
constraint on $p$, and the source never says where the constraint comes from.
It comes from Convention~\ref{conv:standing}\ref{conv:special}; see
Proposition~\ref{prop:p-divides} below. We therefore keep the two roles of
$\mathbf Z_{p}$ apart: the abelian \emph{group} $(\mathbb Z_{p};+,-)$, which
supplies the relation $x_{1}-x_{2}=x_{3}-x_{4}$ in
Definition~\ref{def:linear-congruence} and carries no arity constraint, and the
\emph{algebra} $\mathbf Z_{p}=(\mathbb Z_{p};x_{1}+\dots+x_{n})\in\Vn$, which
does. Only the second occurs in Definition~\ref{def:perfect-linear}, where
$\zeta\le\mathbf A\times\mathbf A\times\mathbf Z_{p}$ is a subuniverse of a
product and therefore needs a common signature.

\item\label{leg:cover}\textbf{$\cover\sigma$ is a tolerance.} It is a
reflexive symmetric subuniverse of $\mathbf A^{2}$ and \emph{not} in general a
congruence (Proposition~\ref{prop:cover}). That $\cover\sigma$ is a congruence
is item (b) of the definition of a linear congruence and conclusion (a) of
Lemma~\ref{lem:nontrivial-bridge}; typing it as a congruence would make the
first vacuous and collapse the linear/PC distinction.

\item\label{leg:empty}\textbf{The empty set.} As in
Convention~\ref{conv:empty}: subuniverses may be empty, subalgebras may not,
and the typed relations hold only between nonempty sets. Two statements of the
source must be restated because they are true only under the opposite reading.
Read syntactically, $\varnothing$ is vacuously absorbing and vacuously central,
so \emph{(a)} the clause defining $\subt{\TS}$ needs its witness $D$ required
nonempty (Definition~\ref{def:ba-center-free}), and \emph{(b)}
Lemma~\ref{lem:central-relation} needs an explicit alternative ``$C=\varnothing$''
--- without it the statement is false, as $\mathbf A=\mathbf B=\mathbf Z_{p}$
with $R=\{(x,x+1)\}$ shows. Likewise Lemma~\ref{lem:linked-implies-abs}
concludes ``there exists a BA or central subuniverse'', which is vacuous
because every algebra is a binary absorbing subuniverse of itself; read
\emph{proper nonempty}, which is how it is used.

\item\label{leg:emptyfamily}\textbf{The empty family is allowed.}
In Definition~\ref{def:irreducible} the family $S_{1},\dots,S_{k}$ may be
empty, the empty intersection being $A^{2}$. This is not a free choice: the
two formulations given there agree \emph{only} under it. Forbid $k=0$ and the
second no longer excludes $\sigma=A^{2}$ while the first still does. With
$k=0$ allowed, an irreducible congruence is automatically proper --- which
recovers the word ``proper'' that \cite{ZhukJACM} carries and
\cite{ZhukSimplified} drops --- and $\cover\sigma$ exists.

\item\label{leg:chain}\textbf{$\lll$ carries its chain.} $C\lll^{A}B$ is
\emph{data}: a list of intermediate sets with their typed steps and witnessing
congruences, not merely the proposition that such a list exists. Several
statements speak of ``the dividing congruences coming from $C\lll^{A}B$'',
which is a function of the chain and not of the pair. See
Hazard~\ref{haz:lll-chain} for the two places where this is load-bearing.

\item\label{leg:dimension}\textbf{``Dimension''.} $\prod_{i}\mathbb Z_{q_{i}}$
with distinct primes is not a vector space over anything and has no dimension.
Use the composition length of the finite abelian group, which agrees with the
source wherever the source is meaningful, and prove its additivity along short
exact sequences. A free choice; nothing in
\S\S\ref{sec:types}--\ref{sec:properties} uses the notion.

\item\label{leg:linked}\textbf{``Linked instance''.} Two inequivalent
definitions coexist in the source: global connectivity of the value graph, and
the per-variable condition. The second does not imply the first, because a
fragmented instance is per-variable linked in each component. We use the
per-variable definition and carry non-fragmentation as a separate hypothesis
wherever the informal reading was doing that work
(Hazard~\ref{haz:linked}).

\item\label{leg:quotient}\textbf{Quotients of subsets.} $B/\delta$ always
means the \emph{image} $\{b/\delta:b\in B\}$ in $\mathbf A/\delta$, for $B$ any
subset of $A$; never the quotient of $B$ by a congruence of $\mathbf B$.
Consequently $(C_{1}\cap\dots\cap C_{t})/\delta\subseteq\bigcap_{i}C_{i}/\delta$
always, with equality false in general
(Hazard~\ref{haz:quotient-rel}). The reverse inclusion is established on the
fly inside the proof of Lemma~\ref{lem:propagation}(fm) and must be extracted
as a standalone lemma with its hypothesis visible.

\item\label{leg:notice}\textbf{Sentences beginning ``Notice that''.} Three of
them are used as lemmas and never proved. They are promoted here:
Lemma~\ref{lem:linear-bridge} (the relation $S$ of
Definition~\ref{def:linear-congruence}(c) is a bridge with
$\bcol S=\proj_{1,2}(S)=\proj_{3,4}(S)=\cover\sigma$),
Lemma~\ref{lem:irred-quotient} ($\sigma$ is irreducible, linear or PC exactly
when $0_{\mathbf A/\sigma}$ is), and
Lemma~\ref{lem:sfree-equiv} (the two formulations of $\TS$-freeness agree).
Lemma~\ref{lem:rect-basics} belongs to the same list.

\item\label{leg:central}\textbf{A citation that drops a case.} The source's
Lemma~\ref{lem:central-relation} cites \cite[Thm.~6.15]{ZhukStrong} but omits
that theorem's third alternative, a nontrivial \emph{projective} subuniverse.
The elimination is legitimate under Convention~\ref{conv:standing}, but it is a
lemma and not a convention, and the source does not state it: see
Lemma~\ref{lem:no-unary-quotient} and Lemma~\ref{lem:central-relation}.

\end{enumerate}

\begin{proposition}[Where $p\mid n-1$ comes from]\label{prop:p-divides}
Let $\mathbf U$ be a finite algebra with a single basic operation $w$ of arity
$n$ that is idempotent, weak near-unanimity and special, and suppose
$\mathbf U$ is affine over $\mathbb Z_{p}$ with $|U|>1$. Then $p\mid n-1$, and
$w$ acts on $\mathbf U$ as $x_{1}+\dots+x_{n}$.
\end{proposition}

\begin{proof}
Being affine, $\mathbf U$ is polynomially equivalent to a
$\mathbb Z_{p}$-module, so $w$ acts as $w(\bar x)=\sum_{i}c_{i}x_{i}+d$ for
scalars $c_{i}\in\mathbb Z_{p}$ and a constant $d$. Idempotence gives
$\bigl(\sum_{i}c_{i}\bigr)x+d=x$ for all $x$, hence $\sum_{i}c_{i}=1$ and
$d=0$.

The weak near-unanimity identities
$w(y,x,\dots,x)=w(x,y,x,\dots,x)=\dots=w(x,\dots,x,y)$ read
$c_{1}y+(1-c_{1})x=c_{2}y+(1-c_{2})x=\dots$, and comparing coefficients of $y$
gives $c_{1}=\dots=c_{n}=:c$, so $nc=1$.

Specialness, $w(x,\dots,x,y)=w\bigl(x,\dots,x,w(x,\dots,x,y)\bigr)$, reads
\[
  (n-1)c\,x+cy \;=\; (n-1)c\,x+c\bigl((n-1)c\,x+cy\bigr),
\]
and comparing coefficients of $y$ gives $c=c^{2}$, so $c\in\{0,1\}$. Now $c=0$
contradicts $nc=1$, so $c=1$; then $n\cdot 1=1$ in $\mathbb Z_{p}$, i.e.\
$p\mid n-1$, and $w$ acts as $x_{1}+\dots+x_{n}$.
\end{proof}

\begin{corollary}\label{cor:Zp-in-Vn}
$\mathbf Z_{p}=(\mathbb Z_{p};x_{1}+\dots+x_{n})$ lies in $\Vn$ if and only if
$p\mid n-1$; and whenever a prime $p$ arises from a linear congruence on some
$\mathbf A\in\Vn$ --- so that some subquotient of $\mathbf A$ is affine over
$\mathbb Z_{p}$ and nontrivial --- one has $p\mid n-1$. Hence the source's
standing identification of $\mathbf Z_{p}$ is legitimate.
\end{corollary}

\begin{hazard}[The proposition is not decoration]\label{haz:p-divides}
Proposition~\ref{prop:p-divides} is needed twice, and in both places its
absence would leave a statement that does not typecheck. It is what makes
$\mathbf Z_{p}$ an object of $\Vn$ at all, hence what makes
Definition~\ref{def:perfect-linear} well formed; and it is what upgrades the
group isomorphism of Lemma~\ref{lem:linear-on-top} to an isomorphism of
algebras. Verified by exhaustion for $p\in\{2,3,5,7\}$ and $3\le n<12$: a
special idempotent affine WNU of arity $n$ on $\mathbb Z_{p}$ exists exactly
when $p\mid n-1$, and is then the sum.
\end{hazard}
%% ---------------------------------------------------------------- ch2-vocabulary.tex
\section{Vocabulary}\label{sec:vocabulary}

This section fixes the universal-algebraic vocabulary. It is Zhuk's
\S2.1, rewritten so that every definition is unambiguous when read by a
machine. The substantive content begins in \S\ref{sec:types}; a reader who
knows the subject should skim to
Definition~\ref{def:irreducible} (irreducible congruences) and
Definition~\ref{def:bridge} (bridges), which are the two places where the
conventions genuinely matter.

\subsection{Algebras, subuniverses, terms}

\begin{definition}[Algebra]\label{def:algebra}
An algebra is a pair $\mathbf A=(A;w^{\mathbf A})$ with $A$ a finite nonempty
set and $w^{\mathbf A}\colon A^{n}\to A$, subject to
Convention~\ref{conv:standing}. We write $A$ for the universe of $\mathbf A$
and drop the superscript on $w$ when no confusion arises.
\end{definition}

\begin{definition}[The algebras $\mathbf Z_{p}$]\label{def:Zp}
For a prime $p$ with $p\mid n-1$, $\mathbf Z_{p}$ is the algebra on
$\{0,1,\dots,p-1\}$ with $w^{\mathbf Z_{p}}(x_{1},\dots,x_{n})
=x_{1}+\dots+x_{n}\bmod p$.
\end{definition}

\begin{hazard}[$\mathbf Z_{p}$ is only sometimes in $\Vn$]\label{haz:Zp}
Zhuk defines $\mathbf Z_{p}$ for every prime $p$ and then asserts that
``every algebra $\mathbf Z_{p}$ belongs to $\Vn$''. That assertion carries a
silent side condition. The operation $x_{1}+\dots+x_{n}\bmod p$ is idempotent
exactly when $n\equiv 1\pmod p$: idempotence says $nx\equiv x$ for all $x$,
i.e.\ $p\mid n-1$. Given that, it is also a WNU (all the identities read
$(n-1)x+y=y$) and special (both sides reduce to $y$).

So $\mathbf Z_{p}\in\Vn$ if and only if $p\mid n-1$, and every statement of the
form ``$\mathbf A/\sigma\cong\mathbf Z_{p}$ for some prime $p$''
(Lemma~\ref{lem:linear-on-top}, Definition~\ref{def:perfect-linear}) silently
produces a $p$ dividing $n-1$. A formalization must decide whether
$\mathbf Z_{p}$ is defined for all $p$ and merely fails to lie in $\Vn$ for bad
$p$, or is defined only for $p\mid n-1$. The first is cleaner; the second makes
the membership free. Either way the side condition must be visible, because
``for some prime $p$'' otherwise reads as an unconstrained existential.
\end{hazard}

\begin{definition}[Subuniverse, subalgebra]\label{def:subuniverse}
$B\subseteq A$ is a \emph{subuniverse} of $\mathbf A$ if
$w(b_{1},\dots,b_{n})\in B$ for all $b_{1},\dots,b_{n}\in B$. If moreover
$B\neq\varnothing$ then $\mathbf B=(B;w\!\restriction_{B^{n}})$ is a
\emph{subalgebra}, written $\mathbf B\le\mathbf A$.
\end{definition}

\begin{definition}[Term operations]\label{def:clone}
$\Clo(\mathbf A)$ is the least set of finitary operations on $A$ containing all
projections and closed under composition with $w$. Its $m$-ary part is written
$\Clo_{m}(\mathbf A)$. Elements of $\Clo(\mathbf A)$ are the \emph{term
operations} of $\mathbf A$.
\end{definition}

\begin{hazard}[Terms versus term operations]\label{haz:term-vs-op}
Zhuk writes $t\in\Clo(\mathbf A)$ and treats $t$ as both a syntactic term and
the operation it induces. These must be distinguished in a formalization: a
term is a finite tree, and distinct terms may induce the same operation. Every
statement below that quantifies over $\Clo(\mathbf A)$ is a statement about
\emph{operations}, so the syntactic layer may be quotiented away --- but it
cannot be skipped, because $\Sg$ is characterised through terms
(Lemma~\ref{lem:sg-terms}) and because the arity bookkeeping in the absorption
arguments is syntactic.
\end{hazard}

\begin{definition}[Generation]\label{def:sg}
For $X\subseteq A^{m}$, $\Sg_{\mathbf A}(X)$ is the least subuniverse of
$\mathbf A^{m}$ containing $X$.
\end{definition}

\begin{lemma}[Generation by terms]\label{lem:sg-terms}
Let $X=\{x_{1},\dots,x_{k}\}\subseteq A^{m}$ be nonempty. Then
\[
  \Sg_{\mathbf A}(X)=\bigl\{\,t^{\mathbf A^{m}}(x_{1},\dots,x_{k})
      \;:\;t\in\Clo_{k}(\mathbf A)\,\bigr\}.
\]
The generator list is \emph{fixed}: the same list $x_{1},\dots,x_{k}$ in the
same order serves every element of the closure, with $t$ varying.
\end{lemma}

\begin{hazard}\label{haz:sg-fixed-list}
The ``fixed list'' phrasing is what makes Lemma~\ref{lem:sg-terms} usable. The
weaker reading --- for each $y$ in the closure there are \emph{some}
generators and \emph{some} term --- is also true but useless, because every
later application applies one term simultaneously to several closure elements
and needs their generator lists to coincide. This was a defect found in review
of the predecessor blueprint and is called out here for the same reason.
\end{hazard}

\subsection{Relations}

\begin{definition}[Relational vocabulary]\label{def:rel-vocab}
Let $R\subseteq A_{1}\times\dots\times A_{m}$.
\begin{enumerate}\alphlabels\tightlist
\item $R$ is \emph{subdirect} if $\proj_{i}(R)=A_{i}$ for every $i$;
  $\mathbf R\sdle\mathbf A_{1}\times\dots\times\mathbf A_{m}$ abbreviates
  ``$R$ is a subdirect subuniverse''.
\item For $R\subseteq A^{m}$: $R$ is \emph{reflexive} if
  $(a,\dots,a)\in R$ for every $a\in A$.
\item For binary $\delta_{1}\subseteq A_{1}\times A_{2}$,
  $\delta_{2}\subseteq A_{2}\times A_{3}$:
  $\delta_{1}\circ\delta_{2}=\{(a,c):\exists b\;(a,b)\in\delta_{1}\wedge
  (b,c)\in\delta_{2}\}$, and $\delta^{-1}=\{(b,a):(a,b)\in\delta\}$.
\item For $B\subseteq A_{1}$ and $\delta\subseteq A_{1}\times A_{2}$:
  $B\circ\delta=\{c:\exists b\in B\;(b,c)\in\delta\}$.
\item A subdirect $\delta\subseteq A\times B$ is \emph{linked} if the bipartite
  graph on $A\sqcup B$ with edge set $\delta$ is connected. It is
  \emph{bijective} if $|\delta|=|A|=|B|$.
\item For binary $\sigma$ and $m\ge 2$,
  $\sigma^{[m]}=\{(a_{1},\dots,a_{m}):(a_{i},a_{j})\in\sigma\ \forall i,j\}$.
\end{enumerate}
\end{definition}

\begin{hazard}[Composition is diagrammatic]\label{haz:comp-order}
$\delta_{1}\circ\delta_{2}$ applies $\delta_{1}$ \emph{first}. This is the
order of \texttt{Rel.comp} and the opposite of \texttt{Function.comp}. Every
composite of bridges and of congruences below is written left to right, and
silently reversing it produces statements that are true-looking and wrong ---
in Corollary~\ref{cor:stable-intersection}(l) the composite
$\sigma_{k}\circ\proj_{k,\ell}(R)\circ\sigma_{\ell}$ is asymmetric in $k$ and
$\ell$ and the order is the content.
\end{hazard}

\begin{hazard}[``Linked'' is used for two different things]\label{haz:linked}
Zhuk uses \emph{linked} both for a binary relation (Definition~\ref{def:rel-vocab}(e))
and for a CSP instance (Definition~\ref{def:instance-linked}). They are related
but not the same predicate, and the coincidence of names causes real confusion
in \S\ref{sec:main-statements}, where both occur in one proof. Use two names.
\end{hazard}

\begin{definition}[Quotients of relations]\label{def:quotient-rel}
For an equivalence $\sigma$ on $A$ and $a\in A$, $a/\sigma$ is the class of
$a$; for $B\subseteq A$, $B/\sigma=\{b/\sigma:b\in B\}$; for
$R\subseteq A^{m}$, $R/\sigma=\{(b_{1}/\sigma,\dots,b_{m}/\sigma):\bar b\in R\}$.
\end{definition}

\begin{hazard}[Images do not commute with intersection]\label{haz:quotient-rel}
$B/\sigma$ and $R/\sigma$ are \emph{images}. Three consequences that the source
navigates silently:
\begin{enumerate}\alphlabels\tightlist
\item $B/\sigma$ is a subset of $A/\sigma$, not the quotient of $B$ by
  $\sigma\cap B^{2}$. For $B$ a subuniverse and $\sigma$ a congruence the two
  are canonically isomorphic as algebras, and the source moves between them
  freely. Fix $B/\sigma$ to be the image and prove the isomorphism once.
\item In general $(C_{1}\cap C_{2})/\sigma\subsetneq C_{1}/\sigma\cap
  C_{2}/\sigma$. Lemma~\ref{lem:propagation}(fm) needs equality for a family
  $C_{1},\dots,C_{t}$, and gets it from the extra hypothesis
  $(C_{i}\circ\sigma)\cap B=C_{i}$, i.e.\ that each $C_{i}$ is
  $\sigma$-saturated inside $B$. That is a content step, not notation.
\item $R/\sigma$ need not be stable under anything, and need not be a
  subuniverse of $(\mathbf A/\sigma)^{m}$ unless $\sigma$ is a congruence.
\end{enumerate}
\end{hazard}

\subsection{Rectangularity}

\begin{definition}[Parallelogram property, rectangularity]\label{def:parallelogram}
Let $R\subseteq A_{1}\times\dots\times A_{m}$.
\begin{enumerate}\alphlabels\tightlist
\item The $i$-th variable of $R$ is \emph{rectangular} if for all
  $\bar a,\bar b$,
  \[
    \bar a[i\mapsto b_{i}]\in R,\quad
    \bar b[i\mapsto a_{i}]\in R,\quad
    \bar b\in R
    \;\Longrightarrow\;\bar a\in R,
  \]
  where $\bar a[i\mapsto c]$ is $\bar a$ with its $i$-th entry replaced by $c$.
  $R$ is \emph{rectangular} if every variable is.
\item $R$ \emph{has the parallelogram property} if for every permutation of the
  coordinates giving $R'$ and every $\ell\in[m-1]$,
  \[
    (a_{1},\dots,a_{\ell},b_{\ell+1},\dots,b_{m})\in R',\quad
    (b_{1},\dots,b_{\ell},a_{\ell+1},\dots,a_{m})\in R',\quad
    \bar b\in R'
    \;\Longrightarrow\;\bar a\in R'.
  \]
\item The \emph{rectangular closure} of $R$ is the least rectangular relation
  containing $R$.
\end{enumerate}
\end{definition}

\begin{lemma}[Two assertions the source makes in passing]\label{lem:rect-basics}
\begin{enumerate}\alphlabels\tightlist
\item A relation with the parallelogram property is rectangular.
\item The rectangular closure exists.
\end{enumerate}
\end{lemma}

\begin{proof}
(a) Apply the parallelogram property to the permutation carrying coordinate
$i$ to position $1$, with $\ell=1$. (b) Rectangularity at coordinate $i$ is a
Horn implication, so the rectangular relations containing $R$ are closed under
arbitrary intersection and the family is nonempty (it contains
$A_{1}\times\dots\times A_{m}$); the intersection of all of them is the least
one.
\end{proof}

\begin{hazard}\label{haz:rect-closure}
Both parts are asserted without proof in \cite{ZhukSimplified} --- (a) as ``as
it follows from the definitions'', (b) by the definite article in ``\emph{the}
minimal rectangular relation''. Both are easy and both must be present. Note
that ``$R$ is rectangular'' unqualified --- the form used in ``rectangular
closure'' --- is never defined in the source; it means every variable is
rectangular.
\end{hazard}

\begin{hazard}[Quantifying over permutations]\label{haz:permutations}
The parallelogram property is stated ``any permutation of its variables gives a
relation $R'$ satisfying \dots''. Formally this is a quantification over
$\mathfrak S_{m}$ combined with a quantification over the split point $\ell$;
equivalently, and more usably, over all ways of splitting $[m]$ into two
nonempty blocks. The second formulation is the one to formalize: it avoids
transporting $R$ along a permutation, which in a dependently typed setting
means an equivalence of index types and a coercion at every use.
\end{hazard}

\subsection{Congruences}

\begin{definition}[Stability]\label{def:stable}
Let $R\subseteq A_{1}\times\dots\times A_{m}$ and let $\sigma$ be a binary
relation on $A_{i}$. The $i$-th variable of $R$ is \emph{stable under $\sigma$}
if $\bar a\in R$ and $(a_{i},b)\in\sigma$ imply $\bar a[i\mapsto b]\in R$.
$R$ is \emph{stable under $\sigma$} if every variable is (each with respect to
$\sigma$ on the corresponding factor, when the factors agree).
\end{definition}

\begin{definition}[Congruence]\label{def:congruence}
A \emph{congruence} on $\mathbf A$ is an equivalence relation on $A$ that is a
subuniverse of $\mathbf A\times\mathbf A$. It is \emph{nontrivial} if it is
neither the equality relation $0_{\mathbf A}$ nor $A^{2}$. The quotient
$\mathbf A/\sigma$ is the algebra on $A/\sigma$ with
$w(\bar a/\sigma)=w(\bar a)/\sigma$.
\end{definition}

\begin{definition}[Irreducible congruence and its cover]\label{def:irreducible}
A congruence $\sigma$ on $\mathbf A$ is \emph{irreducible} if it cannot be
written as an intersection of binary subuniverses of $\mathbf A\times\mathbf A$
each stable under $\sigma$ and each properly containing $\sigma$. Equivalently:
there are no $\mathbf S_{1},\dots,\mathbf S_{k}\le\mathbf A/\sigma\times
\mathbf A/\sigma$ with $0_{\mathbf A/\sigma}=S_{1}\cap\dots\cap S_{k}$ and
$S_{i}\neq 0_{\mathbf A/\sigma}$ for every $i$.

For irreducible $\sigma$, $\cover{\sigma}$ denotes the least
$\delta\le\mathbf A\times\mathbf A$ with $\delta\supsetneq\sigma$ and $\delta$
stable under $\sigma$.
\end{definition}

\begin{hazard}[Three traps in Definition~\ref{def:irreducible}]\label{haz:irreducible}
\begin{enumerate}\alphlabels\tightlist
\item Irreducibility is \emph{not} meet-irreducibility in the congruence
  lattice. The $\mathbf S_{i}$ range over binary \emph{subuniverses} stable
  under $\sigma$, which need not be equivalence relations. The two notions
  differ, and every use below is the stated one.
\item $\cover{\sigma}$ need not be a congruence. It is a subuniverse of
  $\mathbf A^{2}$ stable under $\sigma$, nothing more; that $\cover{\sigma}$
  \emph{is} a congruence is an extra hypothesis in the definition of a linear
  congruence (Definition~\ref{def:linear-congruence}), which would be
  redundant otherwise.
\item Existence and uniqueness of $\cover{\sigma}$ need an argument, given in
  Proposition~\ref{prop:cover}. Uniqueness is exactly what irreducibility buys;
  without it there is no ``the'' cover. A formalization must produce
  $\cover{\sigma}$ as data attached to a proof of irreducibility, not as a
  standalone function.
\item The family $S_{1},\dots,S_{k}$ may be \emph{empty}
  (item \ref{leg:emptyfamily} of
  \S\ref{sec:legislation}). This is not a free choice: the two formulations
  above agree only under it.
\end{enumerate}
\end{hazard}

\begin{proposition}[The cover exists, and is a tolerance]\label{prop:cover}
Let $\sigma$ be an irreducible congruence on $\mathbf A$. Then
\begin{enumerate}\alphlabels\tightlist
\item $\sigma\neq A^{2}$;
\item the family $\mathcal F$ of $\delta\le\mathbf A\times\mathbf A$ with
  $\delta\supsetneq\sigma$ and $\delta$ stable under $\sigma$ is nonempty and
  closed under intersection, so it has a least element $\cover{\sigma}$;
\item $\cover{\sigma}$ is reflexive and symmetric.
\end{enumerate}
\end{proposition}

\begin{proof}
(a) The empty family is allowed, and its intersection is $A^{2}$; so if
$\sigma=A^{2}$ then $\sigma$ is an intersection of a family of relations each
properly containing it, vacuously, and $\sigma$ is reducible.

(b) $\mathcal F\ni A^{2}$ by (a). If $\delta_{1},\delta_{2}\in\mathcal F$ then
$\delta_{1}\cap\delta_{2}$ is again a subuniverse of $\mathbf A^{2}$ stable
under $\sigma$ and contains $\sigma$; it contains $\sigma$ \emph{properly},
since otherwise $\sigma=\delta_{1}\cap\delta_{2}$ exhibits $\sigma$ as an
intersection of two members of $\mathcal F$, contradicting irreducibility. So
$\mathcal F$ is closed under binary, hence finite, intersection, and
$\cover{\sigma}=\bigcap\mathcal F\in\mathcal F$ is its least element.

(c) $\cover{\sigma}\supseteq\sigma\supseteq 0_{\mathbf A}$ is reflexive. For
symmetry, $(\cover{\sigma})^{-1}$ is a subuniverse of $\mathbf A^{2}$, is
stable under $\sigma$ because $\sigma$ is symmetric, and properly contains
$\sigma^{-1}=\sigma$; so $(\cover{\sigma})^{-1}\in\mathcal F$ and therefore
$\cover{\sigma}\subseteq(\cover{\sigma})^{-1}$, whence equality.
\end{proof}

\begin{remark}[Why not a congruence]\label{rem:cover-not-cong}
Nothing in Proposition~\ref{prop:cover} gives transitivity, and the source
supplies the internal evidence that none is available: that $\cover{\sigma}$ is
a congruence is item (b) of Definition~\ref{def:linear-congruence} and
conclusion (a) of Lemma~\ref{lem:nontrivial-bridge}, both of which would be
redundant otherwise. Symmetry of $\cover{\sigma}$, by contrast, is used
silently in \cite{ZhukSimplified} --- in the proof of
Lemma~\ref{lem:pc-bridges-trivial}, where the argument is run again ``switching
$x_{1}$ and $x_{2}$''.
\end{remark}

\begin{lemma}[Passing to the quotient]\label{lem:irred-quotient}
Let $\sigma$ be a congruence on $\mathbf A$. Then $\delta\mapsto\delta/\sigma$
is an inclusion-preserving bijection, commuting with intersections, from the
subuniverses of $\mathbf A^{2}$ that contain $\sigma$ and are stable under
$\sigma$ onto the subuniverses of $(\mathbf A/\sigma)^{2}$, sending $\sigma$ to
$0_{\mathbf A/\sigma}$. Consequently $\sigma$ is irreducible, linear or PC
exactly when $0_{\mathbf A/\sigma}$ is, and then
$\cover{\sigma}/\sigma=\cover{(0_{\mathbf A/\sigma})}$.
\end{lemma}

\begin{proof}
A subuniverse $\delta$ of $\mathbf A^{2}$ containing $\sigma$ and stable under
$\sigma$ is a union of products of $\sigma$-blocks, so it is the preimage of
$\delta/\sigma$; conversely the preimage of a subuniverse of
$(\mathbf A/\sigma)^{2}$ is a subuniverse containing $\sigma$ and stable under
$\sigma$. The two constructions are mutually inverse and preserve inclusion and
intersection, and $\sigma$ is the preimage of $0_{\mathbf A/\sigma}$. The
statement of Definition~\ref{def:irreducible} is a statement about that family
alone, so it transports; and both the least element in
Proposition~\ref{prop:cover}(b) and the data of
Definition~\ref{def:linear-congruence} transport with it.
\end{proof}

\begin{hazard}[This is the licence for ``assume $\sigma=0$'']\label{haz:factor-out}
Three proofs in \S\ref{sec:properties} open with ``we may factor by $\sigma$
and assume $\sigma$ is the equality relation''.
Lemma~\ref{lem:irred-quotient} is what licenses that, and it is one of the
sentences \cite{ZhukSimplified} states without proof
(\S\ref{sec:legislation}, item \ref{leg:notice}).
\end{hazard}

\subsection{Bridges}

\begin{definition}[Bridge]\label{def:bridge}
Let $\sigma_{1},\sigma_{2}$ be congruences on $\mathbf D_{1},\mathbf D_{2}$. A
relation $\delta\le\mathbf D_{1}^{2}\times\mathbf D_{2}^{2}$ is a \emph{bridge
from $\sigma_{1}$ to $\sigma_{2}$} if
\begin{enumerate}\alphlabels\tightlist
\item the first two variables of $\delta$ are stable under $\sigma_{1}$;
\item the last two variables of $\delta$ are stable under $\sigma_{2}$;
\item $\proj_{1,2}(\delta)\supsetneq\sigma_{1}$ and
      $\proj_{3,4}(\delta)\supsetneq\sigma_{2}$;
\item $(a_{1},a_{2},a_{3},a_{4})\in\delta$ implies
      $\bigl((a_{1},a_{2})\in\sigma_{1}\iff(a_{3},a_{4})\in\sigma_{2}\bigr)$.
\end{enumerate}
For a bridge $\delta$ put $\bcol{\delta}=\{(x,y):(x,x,y,y)\in\delta\}$.
A bridge $\delta\subseteq D^{4}$ is \emph{reflexive} if $(a,a,a,a)\in\delta$
for all $a\in D$.
\end{definition}

\begin{remark}[The motivating example]\label{rem:bridge-example}
On $\mathbb Z_{4}$, the relation
$\delta=\{(a_{1},a_{2},a_{3},a_{4}):a_{1}-a_{2}=2a_{3}-2a_{4}\}$ is a bridge
from $0_{\mathbb Z_{4}}$ to congruence mod~$2$. Condition (d) is the content:
$a_{1}=a_{2}$ exactly when $a_{3}\equiv a_{4}\pmod 2$. Bridges are the
formal shadow of a centralizer relation; see \cite{RossSlides}.
\end{remark}

\begin{definition}[Composition of bridges]\label{def:bridge-comp}
For bridges $\delta_{1}$ from $\sigma_{0}$ to $\sigma_{1}$ and $\delta_{2}$
from $\sigma_{1}$ to $\sigma_{2}$, set
\[
  (\delta_{1}\ast\delta_{2})(x_{1},x_{2},z_{1},z_{2})
  =\exists y_{1}\exists y_{2}\;
   \delta_{1}(x_{1},x_{2},y_{1},y_{2})\wedge\delta_{2}(y_{1},y_{2},z_{1},z_{2}).
\]
\end{definition}

\begin{lemma}[Bridge composition; {\cite[Lem.~6.3]{ZhukJACM}}]\label{lem:bridge-comp}
If $\sigma_{0},\sigma_{1},\sigma_{2}$ are irreducible then
$\delta_{1}\ast\delta_{2}$ is a bridge from $\sigma_{0}$ to $\sigma_{2}$, and
$\bcol{\delta_{1}\ast\delta_{2}}=\bcol{\delta_{1}}\circ\bcol{\delta_{2}}$.
\end{lemma}

\begin{hazard}\label{haz:bridge-comp}
Lemma~\ref{lem:bridge-comp} is imported, not proved, in
\cite{ZhukSimplified}; it is \cite[Lemma 6.3]{ZhukJACM}. Both conclusions are
used: the first in every chain of bridges, the second --- the identity
$\bcol{\cdot}$ commutes with composition --- in
Lemma~\ref{lem:perfect-from-linked}, where a bridge with $\bcol{\delta}$ linked
is powered up until $\bcol{\delta}=A^{2}$. The irreducibility hypothesis on
\emph{all three} congruences is needed and is easy to drop by accident.
\end{hazard}

\begin{definition}[Symmetric and pair-reflexive bridges]\label{def:pair-reflexive}
Let $\delta$ be a bridge from $\sigma$ to $\sigma$ on $\mathbf D$. Call $\delta$
\emph{symmetric} if
$\delta(x_{1},x_{2},x_{3},x_{4})\iff\delta(x_{3},x_{4},x_{1},x_{2})$, and
\emph{pair-reflexive} if
\[
  (x_{1},x_{2})\in\proj_{1,2}(\delta)\;\Longrightarrow\;
  (x_{1},x_{2},x_{1},x_{2})\in\delta .
\]
Write $\delta^{-1}(x_{1},x_{2},x_{3},x_{4})=\delta(x_{3},x_{4},x_{1},x_{2})$.
\end{definition}

\begin{lemma}[Symmetrisation]\label{lem:symmetrise}
Let $\delta$ be a reflexive bridge from $\sigma$ to $\sigma$ on $\mathbf D$ and
put $\delta^{\circ}=\delta\ast\delta^{-1}$, that is
\[
  \delta^{\circ}(x_{1},x_{2},x_{3},x_{4})=\exists x_{5}\exists x_{6}\;
  \delta(x_{1},x_{2},x_{5},x_{6})\wedge\delta(x_{3},x_{4},x_{5},x_{6}).
\]
Then $\delta^{\circ}$ is a reflexive, symmetric, pair-reflexive bridge from
$\sigma$ to $\sigma$ with
$\proj_{1,2}(\delta^{\circ})=\proj_{3,4}(\delta^{\circ})=\proj_{1,2}(\delta)$
and $\bcol{\delta^{\circ}}=\bcol{\delta}\circ\bcol{\delta}^{-1}
\supseteq\bcol{\delta}$. Moreover any pair-reflexive bridge $\varepsilon$
satisfies $\varepsilon\subseteq\varepsilon\ast\varepsilon$.
\end{lemma}

\begin{proof}
Symmetry is visible in the formula. For pair-reflexivity, if
$(x_{1},x_{2})\in\proj_{1,2}(\delta^{\circ})$ then
$\delta(x_{1},x_{2},x_{5},x_{6})$ for some $x_{5},x_{6}$, and taking the same
witnesses twice gives $(x_{1},x_{2},x_{1},x_{2})\in\delta^{\circ}$; the same
computation shows $\proj_{1,2}(\delta^{\circ})=\proj_{1,2}(\delta)$, the
inclusion $\subseteq$ being immediate, and symmetry then gives
$\proj_{3,4}(\delta^{\circ})=\proj_{1,2}(\delta^{\circ})$. Clause (c) of
Definition~\ref{def:bridge} follows, and clause (d) because
$(x_{1},x_{2})\in\sigma\iff(x_{5},x_{6})\in\sigma\iff(x_{3},x_{4})\in\sigma$.
Stability is inherited from $\delta$, and reflexivity from
$(a,a,a,a)\in\delta$. The identity
$\bcol{\delta^{\circ}}=\bcol{\delta}\circ\bcol{\delta^{-1}}
=\bcol{\delta}\circ\bcol{\delta}^{-1}$ is Lemma~\ref{lem:bridge-comp};
it contains $\bcol{\delta}$ because $\bcol{\delta}$ is reflexive, so
$(x,y)\in\bcol{\delta}$ gives $(x,y)\in\bcol{\delta}\circ\bcol{\delta}^{-1}$ by
taking $y$ as the intermediate point. Finally, if $\varepsilon$ is
pair-reflexive and $\varepsilon(x_{1},x_{2},x_{3},x_{4})$, then taking
$(x_{1},x_{2})$ itself as the intermediate pair --- legitimate because
$\varepsilon(x_{1},x_{2},x_{1},x_{2})$ --- gives
$(\varepsilon\ast\varepsilon)(x_{1},x_{2},x_{3},x_{4})$.
\end{proof}

\begin{hazard}[Pair-reflexivity is the hypothesis three lemmas need]\label{haz:pair-reflexive}
\cite{ZhukSimplified} states pair-reflexivity only once, as a hypothesis of
Lemma~\ref{lem:pc-inductive-step}, but three further statements require it and
do not say so; without it one of them is false. See
Hazard~\ref{haz:bridge-repairs}. Lemma~\ref{lem:symmetrise} is what makes this
costless: every bridge these statements are ever applied to is a
$\delta^{\circ}$.
\end{hazard}

\begin{definition}[Perfect linear congruence]\label{def:perfect-linear}
A congruence $\sigma$ on $\mathbf A$ is a \emph{perfect linear congruence} if
it is irreducible and there are a prime $p$ and
$\zeta\le\mathbf A\times\mathbf A\times\mathbf Z_{p}$ with
$\proj_{1,2}(\zeta)=\cover{\sigma}$ and
\[
  (a_{1},a_{2},b)\in\zeta\;\Longrightarrow\;
  \bigl((a_{1},a_{2})\in\sigma\iff b=0\bigr).
\]
\end{definition}

Perfect linear congruences are the payoff of the whole bridge machinery: they
let the passage from $\sigma$ to $\cover{\sigma}$ be tracked by a single
element of $\mathbf Z_{p}$, which is what turns a CSP instance into a system of
linear equations in \S\ref{sec:codim-one}.
\clearpage
%% ---------------------------------------------------------------- ch3-types.tex
\section{The six types of strong subuniverse}\label{sec:types}

This section is Zhuk's \S2.2. It defines six binary relations
$C\subt[A]{T}B$ between subsets of a fixed algebra $\mathbf A$, one for each
type $T$, and the transitive-closure relation $C\lll^{A}B$. Everything in
\S\S\ref{sec:properties}--\ref{sec:main-statements} is expressed in this
language, so the definitions repay care.

The informal reading is: $C\subt[A]{T}B$ says that $C$ is a \emph{strong}
subset of $B$ --- strong enough that the CSP algorithm may reduce a variable's
domain from $B$ to $C$ without losing all solutions --- and $T$ records
\emph{why} it is strong. Three of the reasons are absorption-theoretic
($\TBA$, $\TC$, $\TS$) and three are congruence-theoretic
($\TD$, $\TL$, $\TPC$).

\subsection{Absorption}

\begin{definition}[Absorbing subuniverse]\label{def:absorbing}
A subuniverse $B$ of $\mathbf A$ is \emph{absorbing} if there is
$t\in\Clo_{k}(\mathbf A)$, for some $k$, such that for every $i\in[k]$
\[
  t(B,\dots,B,\underbrace{A}_{i},B,\dots,B)\subseteq B .
\]
We then say \emph{$B$ absorbs $\mathbf A$ with the term $t$}. If $t$ may be
chosen binary, $B$ is a \emph{binary absorbing} subuniverse ($\TBA$); if
ternary, a \emph{ternary absorbing} subuniverse.
\end{definition}

\begin{hazard}[The tuple form]\label{haz:absorption-tuples}
Definition~\ref{def:absorbing} must be formalized over \emph{tuples
constrained coordinatewise}, not as ``for all $b_{1},\dots,b_{k}\in B$ and
$a\in A$, $t(b_{1},\dots,b_{i-1},a,b_{i+1},\dots,b_{k})\in B$''. The two agree,
but the second phrasing invites the variant in which $b_{i}$ is quantified and
then discarded, which is vacuously true when $B=\varnothing$ and at $k=1$
where the intended content is not vacuous. This was an actual defect in the
predecessor blueprint. State it as: for every $\bar a\in A^{k}$ such that
$a_{j}\in B$ for all $j\neq i$, one has $t(\bar a)\in B$.
\end{hazard}

\begin{definition}[Central subuniverse]\label{def:central}
A subuniverse $C$ of $\mathbf A$ is \emph{central} ($\TC$) if it is absorbing
and for every $a\in A\setminus C$,
\[
  (a,a)\notin\Sg_{\mathbf A}\bigl((\{a\}\times C)\cup(C\times\{a\})\bigr).
\]
\end{definition}

\begin{lemma}[{\cite[Cor.~6.11.1]{ZhukStrong}}]\label{lem:central-ternary}
A central subuniverse is a ternary absorbing subuniverse.
\end{lemma}

\begin{remark}[Already formalized]\label{rem:center-done}
Lemma~\ref{lem:central-ternary} is the one deep import of this theory that is
already available in Lean: it is the main theorem of the predecessor project
(\texttt{zhuk\_main} in \texttt{ZhukLean}), proved there from the definitions
via the Barto--Kazda relational description of absorption and the Zhuk--Kozik
doubling trick. See \S\ref{sec:reuse}.
\end{remark}

\begin{definition}[BA and centre free; $\TS$]\label{def:ba-center-free}
$\mathbf A$ is \emph{BA and centre free} if it has no proper nonempty binary
absorbing subuniverse and no proper nonempty central subuniverse.

For $\varnothing\neq C\lneq B\le A$ write $C\subt[A]{\TS}B$ if there is a
\emph{nonempty} subuniverse $D\subseteq C$ of $\mathbf B$ that is
simultaneously binary absorbing and central in $\mathbf B$. $\mathbf A$ is
\emph{$\TS$-free} if there is no $C\subt{\TS}\mathbf A$.
\end{definition}

\begin{hazard}[The witness must be required nonempty]\label{haz:S-nonempty}
$\varnothing$ is vacuously binary absorbing and vacuously central
(\S\ref{sec:legislation}, item \ref{leg:empty}), so without the word
``nonempty'' the clause defining $\subt{\TS}$ is universally true and the type
$\TS$ collapses. The global convention that subuniverses are nonempty repairs
it, but the clause does not inherit that convention on its own, because $D$ is
introduced by an existential inside the clause rather than as a subalgebra of
the ambient discourse.
\end{hazard}

\begin{lemma}[The two readings of $\TS$-freeness agree]\label{lem:sfree-equiv}
$\mathbf A$ is $\TS$-free if and only if no proper nonempty $D\le A$ is
simultaneously binary absorbing and central in $\mathbf A$.
\end{lemma}

\begin{proof}
If such a $D$ exists then $D\subt{\TS}\mathbf A$, taking $C=D$. Conversely if
$C\subt{\TS}\mathbf A$ then by definition some nonempty $D\subseteq C\lneq A$
is binary absorbing and central in $\mathbf A$, and $D$ is proper.
\end{proof}

\begin{hazard}[$\TS$ is not the conjunction of $\TBA$ and $\TC$]\label{haz:S-type}
$C\subt{\TS}B$ does not say that $C$ is binary absorbing and central in
$\mathbf B$; it says that $C$ \emph{contains} some $D$ that is. The type
$\TS$ exists precisely so that it is closed under enlarging $C$, which
$\TBA$ and $\TC$ are not. Note also the asymmetry with $\TS$-freeness, which
\emph{is} stated as a conjunction. Getting this backwards makes
Lemma~\ref{lem:propagation}(ft) unprovable.
\end{hazard}

\subsection{Linear and PC congruences}

\begin{definition}[Linear congruence]\label{def:linear-congruence}
A congruence $\sigma$ on $\mathbf A$ is \emph{linear} if
\begin{enumerate}\alphlabels\tightlist
\item $\sigma$ is irreducible;
\item $\cover{\sigma}$ is a congruence;
\item there are a prime $p$ and $S\le(\cover{\sigma})^{[4]}$ such that for
  every block $B$ of $\cover{\sigma}$ there is $m\ge 0$ with
  \[
    \bigl(B/\sigma;\;S\cap(B/\sigma)^{4}\bigr)\;\cong\;
    \bigl(\mathbb Z_{p}^{m};\;x_{1}-x_{2}=x_{3}-x_{4}\bigr).
  \]
\end{enumerate}
An irreducible congruence that is not linear is a \emph{PC congruence}.
\end{definition}

Here $(\cover{\sigma})^{[4]}$ denotes the set of $4$-tuples whose entries lie in
a single block of $\cover{\sigma}$, which is meaningful because (b) makes
$\cover{\sigma}$ a congruence. The source records the following in a sentence
beginning ``Notice that'' and then uses it as a lemma
(\S\ref{sec:legislation}, item \ref{leg:notice}).

\begin{lemma}[The witness of (c) is a bridge]\label{lem:linear-bridge}
Let $\sigma$ be a linear congruence and $S$ as in
Definition~\ref{def:linear-congruence}(c). Then $S$ is a reflexive
pair-reflexive bridge from $\sigma$ to $\sigma$ with
$\bcol{S}=\proj_{1,2}(S)=\proj_{3,4}(S)=\cover{\sigma}$.
\end{lemma}

\begin{proof}
Work inside a block $B$ of $\cover{\sigma}$ and identify $B/\sigma$ with
$\mathbb Z_{p}^{m}$ so that $S\cap(B/\sigma)^{4}$ becomes
$x_{1}-x_{2}=x_{3}-x_{4}$; tuples of $S$ meet no other block. Each variable is
stable under $\sigma$, since the defining equation is an equation between
$\sigma$-classes. Clause (d) of Definition~\ref{def:bridge} holds because
$x_{1}-x_{2}=0$ if and only if $x_{3}-x_{4}=0$, which says
$(x_{1},x_{2})\in\sigma\iff(x_{3},x_{4})\in\sigma$. For the projections, any
$(x_{1},x_{2})$ inside a common block satisfies
$(x_{1},x_{2},x_{1},x_{2})\in S$, so $\proj_{1,2}(S)=\cover{\sigma}$ --- which
is at once clause (c), since $\cover{\sigma}\supsetneq\sigma$, and
pair-reflexivity; symmetry of the defining equation in the two pairs gives
$\proj_{3,4}(S)=\cover{\sigma}$. Finally $(x,x,y,y)\in S$ reads $0=0$, so
$\bcol{S}$ is the whole of $\cover{\sigma}$, and $(a,a,a,a)\in S$ gives
reflexivity.
\end{proof}

This is the useful reformulation:

\begin{lemma}[{\cite[Lem.~\ref*{lem:linear-equiv}]{ZhukSimplified}}]\label{lem:linear-equiv}
For an irreducible congruence $\sigma$ the following are equivalent:
\begin{enumerate}\alphlabels\tightlist
\item $\sigma$ is linear;
\item there is a bridge $\delta$ from $\sigma$ to $\sigma$ with
  $\bcol{\delta}\supsetneq\sigma$.
\end{enumerate}
\end{lemma}

\begin{hazard}[Definition (c) versus criterion (b)]\label{haz:linear-def}
Definition~\ref{def:linear-congruence}(c) is a per-block isomorphism statement
with an existential $m$ \emph{depending on the block}, and it quantifies over
blocks of $\cover{\sigma}$, not of $\sigma$. Lemma~\ref{lem:linear-equiv}
replaces all of that by one bridge. A formalization should take (b) as the
\emph{definition} and prove (c) as a structure theorem, reversing the paper's
order: (b) is a $\Sigma_1$ statement over finite data and is what every
consumer actually uses, while (c) drags in a quotient, a choice of prime, and a
family of isomorphisms.

The dichotomy ``linear or PC'' is then \emph{by definition} exhaustive on
irreducible congruences, which is how the paper uses it. Note that PC-ness is a
negative condition; every lemma about PC congruences is proved by contradiction
against Lemma~\ref{lem:linear-equiv}(b).
\end{hazard}

\begin{lemma}[No bridge crosses the types]\label{lem:no-cross-bridge}
If $\sigma_{1}$ is linear, $\sigma_{2}$ is irreducible, and there is a bridge
from $\sigma_{1}$ to $\sigma_{2}$, then $\sigma_{2}$ is linear.
\end{lemma}

\begin{lemma}[Bridges from PC congruences are trivial]\label{lem:pc-bridges-trivial}
Let $\sigma$ be a PC congruence on $A$ and $\delta$ a reflexive bridge from
$\sigma$ to $\sigma$ with $\proj_{1,2}(\delta)=\proj_{3,4}(\delta)=
\cover{\sigma}$. Then
\[
  \delta(x_{1},x_{2},x_{3},x_{4})=\sigma(x_{1},x_{3})\wedge\sigma(x_{2},x_{4})
  \quad\text{or}\quad
  \delta(x_{1},x_{2},x_{3},x_{4})=\sigma(x_{1},x_{4})\wedge\sigma(x_{2},x_{3}).
\]
\end{lemma}

\begin{lemma}[Bridges into PC congruences]\label{lem:bridge-to-pc}
Let $\delta$ be a bridge from a PC congruence $\sigma_{1}$ on $\mathbf A_{1}$
to an irreducible congruence $\sigma_{2}$ on $\mathbf A_{2}$, with
$\proj_{1,2}(\delta)=\cover{\sigma_{1}}$ and
$\proj_{3,4}(\delta)=\cover{\sigma_{2}}$. Then
\begin{enumerate}\alphlabels\tightlist
\item $\sigma_{2}$ is a PC congruence;
\item $\mathbf A_{1}/\sigma_{1}\cong\mathbf A_{2}/\sigma_{2}$;
\item $\{(a/\sigma_{1},b/\sigma_{2}):(a,b)\in\bcol{\delta}\}$ is bijective;
\item $\delta(x_{1},x_{2},x_{3},x_{4})=
      \bcol{\delta}(x_{1},x_{3})\wedge\bcol{\delta}(x_{2},x_{4})$ or
      $\delta(x_{1},x_{2},x_{3},x_{4})=
      \bcol{\delta}(x_{1},x_{4})\wedge\bcol{\delta}(x_{2},x_{3})$.
\end{enumerate}
\end{lemma}

Two results connect the new vocabulary to the linear and PC subuniverses of the
original proof.

\begin{lemma}\label{lem:linear-on-top}
If $\sigma$ is a linear congruence on $\mathbf A\in\Vn$ with
$\cover{\sigma}=A^{2}$, then $\mathbf A/\sigma\cong\mathbf Z_{p}$ for a prime
$p$.
\end{lemma}

\begin{lemma}\label{lem:pc-on-top}
If $\sigma$ is a PC congruence on $\mathbf A$ with $\cover{\sigma}=A^{2}$, then
$\mathbf A/\sigma$ is polynomially complete.
\end{lemma}

\begin{hazard}[What the two proofs on top actually need]\label{haz:on-top}
Neither statement is proved in \cite{ZhukSimplified} as it stands.

For Lemma~\ref{lem:linear-on-top} the source produces the witness
$\zeta$ of Definition~\ref{def:perfect-linear}, sets
$\xi(x,z)=\zeta(x,a,z)$ for a fixed $a$, and says ``then $\xi$ is a bijective
relation giving an isomorphism'' --- which is the entire content of the lemma.
We prove it instead from Lemma~\ref{lem:nontrivial-bridge}(c): with
$\cover{\sigma}=A^{2}$ that lemma has a single block, so
$\mathbf A/\sigma\cong\mathbb Z_{p}^{m}$ as a group; $m=0$ would give
$\sigma=A^{2}$, excluded because irreducible congruences are proper; and
$m\ge 2$ is excluded by irreducibility, because the $m$ coordinate congruences
$\theta_{i}=\{(x,y):x_{i}=y_{i}\}$ are subuniverses of $(\mathbf A/\sigma)^{2}$
--- linear subspaces are preserved by any coordinatewise affine map --- each
properly above $0_{\mathbf A/\sigma}$, with $\bigcap_{i}\theta_{i}
=0_{\mathbf A/\sigma}$. So $m=1$. That the resulting bijection is an
isomorphism of \emph{algebras}, rather than of groups, is
Proposition~\ref{prop:p-divides}. This route also avoids the import
\cite[Cor.~8.17.1]{ZhukJACM} that the source's route needs.

For Lemma~\ref{lem:pc-on-top} two things are used and not stated: that an
idempotent algebra is polynomially complete exactly when every \emph{reflexive}
subuniverse of a power is a conjunction of equalities --- those being precisely
the invariants of the clone generated by the basic operations together with all
constants --- and that $\mathbf A/\sigma$ is BA and centre free, without which
the step ``$\{a_{i}\}\subt{\TPC}\mathbf A/\sigma$'' is not licensed by
Definition~\ref{def:types}.
\end{hazard}

\begin{hazard}[Polynomial completeness]\label{haz:pc-algebra}
An algebra is \emph{polynomially complete} if the clone generated by its basic
operations together with all constants is the clone of \emph{all} operations on
its universe. Lemma~\ref{lem:pc-on-top} is where the name ``PC'' comes from,
and it is the only place in the development where polynomial completeness is
used as more than a label. Formalizing it needs the Istinger--Kaiser
characterization \cite{IstingerKaiser}; formalizing everything else needs only
the \emph{name}. A staged development should therefore treat ``PC congruence''
as ``irreducible and not linear'' throughout, and prove
Lemma~\ref{lem:pc-on-top} last or not at all --- it is used only to connect
with \cite{ZhukJACM}, not by any argument here.
\end{hazard}

\subsection{The types}

\begin{definition}[The six types]\label{def:types}
Let $\varnothing\neq C\lneq B\le A$. We write:
\begin{itemize}\tightlist
\item $C\subt[A]{\TBA}B$ if $C$ is a binary absorbing subuniverse of
  $\mathbf B$;
\item $C\subt[A]{\TC}B$ if $C$ is a central subuniverse of $\mathbf B$;
\item $C\subt[A]{\TD}B$ if there is an irreducible congruence $\sigma$ on
  $\mathbf A$ with
  \begin{enumerate}\romanlabels\tightlist
  \item $B^{2}\subseteq\cover{\sigma}$,
  \item $C=B\cap E$ for some block $E$ of $\sigma$,
  \item $\mathbf B/\sigma$ is BA and centre free;
  \end{enumerate}
\item $C\subt[A]{\TL}B$ if $C\subt[A]{\TD}B$ with the witnessing $\sigma$
  linear;
\item $C\subt[A]{\TPC}B$ if $C\subt[A]{\TD}B$ with the witnessing $\sigma$ a PC
  congruence;
\item $C\subt[A]{\TS}B$ as in Definition~\ref{def:ba-center-free}.
\end{itemize}
When the witnessing congruence matters we write $C\subt[A]{T(\sigma)}B$; for
$T\in\{\TBA,\TC,\TS\}$, where no congruence is involved, $\sigma$ is taken to
be $A^{2}$.
\end{definition}

\begin{hazard}[The superscript $A$ is not decoration]\label{haz:superscript}
For $T\in\{\TD,\TL,\TPC\}$ the witnessing congruence lives on $\mathbf A$, not
on $\mathbf B$, and condition (i) relates the two: $B^{2}\subseteq
\cover{\sigma}$ says $B$ sits inside a single block of $\cover{\sigma}$. This
is why the relation carries the ambient algebra as a superscript. For
$T\in\{\TBA,\TC,\TS\}$ the superscript is genuinely irrelevant and Zhuk drops
it. A formalization should not overload: define $\subt[A]{T}$ with $A$ always
present, and prove the independence for the three absorption types as a lemma.

Note also that $\mathbf B/\sigma$ in (iii) means $B/\sigma$ with the induced
algebra structure, which requires $B$ to be a subuniverse and $\sigma$
restricted to $B$; since $B^{2}\subseteq\cover{\sigma}$ and not necessarily
$B^{2}\subseteq\sigma$, the quotient $B/\sigma$ is in general \emph{not} a
single point, and its being BA and centre free is a real condition.
\end{hazard}

\begin{definition}[Multi-types]\label{def:multitypes}
For $T\in\{\TL,\TPC,\TD\}$ write $C\subt[A]{\MT T}B$ if $C\neq\varnothing$ and
$C=C_{1}\cap\dots\cap C_{t}$ with $C_{i}\subt[A]{T}B$ for every $i\in[t]$.
\end{definition}

\begin{definition}[The strong-reduction relation]\label{def:lll}
Write $C\lll^{A}B$ if there are $B_{0},\dots,B_{m}\subseteq B$ and
$T_{1},\dots,T_{m}\in\{\TBA,\TC,\TS,\TD\}$ with
\[
  C=B_{m}\subt[A]{T_{m}}B_{m-1}\subt[A]{T_{m-1}}\dots\subt[A]{T_{1}}B_{0}=B .
\]
Here $m=0$ is allowed, so $\lll^{A}$ is reflexive. A congruence \emph{comes
from} $C\lll^{A}B$ if it is one of the witnessing congruences in such a chain.
We write $B\lll A$ for $B\lll^{A}A$, and $C\subte[A]{T(\sigma)}B$ for
``$C=B$ or $C\subt[A]{T(\sigma)}B$''.
\end{definition}

\begin{hazard}[$\lll$ is data, not a proposition]\label{haz:lll-chain}
It is tempting to read $C\lll^{A}B$ as the reflexive transitive closure of
$\bigcup_{T}\subt[A]{T}$ --- in Lean, \texttt{Relation.ReflTransGen} --- and to
forget the chain. That is wrong, and the error is not cosmetic:
Definition~\ref{def:lll} also defines ``the congruences coming from
$C\lll^{A}B$'', which is a function of the chain and not of the pair. Two
places consume it, and in both the \emph{choice} of chain is what makes the
argument work.

\begin{enumerate}\alphlabels\tightlist
\item In the ``moreover'' clause of Lemma~\ref{lem:intersection-good}(d) the
  conclusion is $\delta\supseteq\delta_{1}\cap\dots\cap\delta_{s}$ where the
  $\delta_{i}$ are the congruences of the chains for $B\lll A$ and $C\lll A$
  \emph{fixed at the top of that proof}; the proof establishes exactly
  $\sigma\cap\omega\subseteq\delta$ for those two intersections.
\item In the type-$\TD$ case of Lemma~\ref{lem:two-stable} the clause is
  applied to $C_{2}$ rather than to $B_{2}$, and its output is the intersection
  of $\omega$ with $\sigma_{2}$ --- which is correct only because the chain
  witnessing $C_{2}\lll A$ is chosen to be the chain witnessing
  $B_{2}\lll A$ extended by the single step $C_{2}\subt{\TD(\sigma_{2})}B_{2}$.
  A different chain gives a different family and the step fails.
\end{enumerate}

So $\lll^{A}$ should be an inductive family carrying, at each step, the
intermediate set, its type, and its witnessing congruence; extension by one
step is an operation on that data. The \texttt{ReflTransGen} version is its
propositional truncation and may be used only where no congruence is named,
which is most but not all of the development. What \emph{is} true, and is the
paper's central simplification over \cite{ZhukJACM}, is that the intermediate
\emph{types} may be forgotten: $\TL$ and $\TPC$ are absent from the list in
Definition~\ref{def:lll} because $\TD$ subsumes them.
\end{hazard}

\begin{convention}[Dotted variants]\label{conv:dotted}
$C\dsubt[A]{T}B$, $C\dsubte[A]{T}B$ and $C\dlll^{A}B$ mean that the undotted
relation holds \emph{or} $C=\varnothing$. See Hazard~\ref{haz:dot}.
\end{convention}
\clearpage
%% ---------------------------------------------------------------- ch4-instances.tex
\section{Instances}\label{sec:instances}

This section is Zhuk's \S3.1. It is entirely combinatorial: no algebra is used
beyond the fact that each domain carries a $w$ and each constraint relation is
a subuniverse of a product of domains. It is also where a formalization must
make the most decisions, because the paper's instances are informal syntactic
objects manipulated by operations --- weakening, covering, renaming --- that are
described in words.

\subsection{Instances and solutions}

\begin{definition}[Instance]\label{def:instance}
Fix a finite set $V$ of variables and, for each $x\in V$, an algebra
$\mathbf D_{x}\in\Vn$. A \emph{constraint} is a pair
$C=(\bar x,R)$ where $\bar x=(x_{1},\dots,x_{m})$ is a tuple of variables and
$R\le\mathbf D_{x_{1}}\times\dots\times\mathbf D_{x_{m}}$. An
\emph{instance} $\mathcal I$ is a finite list of constraints; we write
$C\in\mathcal I$, and $\Var(C)$, $\Var(\mathcal I)$ for the variables
occurring. A \emph{subinstance} is a sublist.

A \emph{solution} is a map $s$ with $s(x)\in D_{x}$ for all $x$ such that
$(s(x_{1}),\dots,s(x_{m}))\in R$ for every constraint. The solution set of
$\mathcal I$ is \emph{subdirect} if for every $x$ and every $a\in D_{x}$ there
is a solution with $s(x)=a$.
\end{definition}

\begin{hazard}[Lists, not sets; tuples, not sets of variables]\label{haz:instance-list}
Two representation choices are forced by the way the paper argues, and both
differ from the naive one.

\emph{Constraints form a list, not a set.} Proofs delete ``a constraint
$C\in\mathcal I$'' and reason about $\mathcal I\setminus\{C\}$, and the same
relation may occur twice on different variable tuples --- and, in the
expanded coverings of \S\ref{sec:coverings}, the same relation on the
\emph{same} tuple can legitimately occur twice. Multiset or list semantics is
required; \texttt{List} with membership by index is the safe choice.

\emph{A constraint's scope is a tuple, not a set.} Repeated variables in a
scope are allowed and occur --- $\zeta(x,x',z)$ in
\S\ref{sec:main-statements} arises by splitting one variable into two, and
$R(x,x,\dots,x)$ appears explicitly in Definition~\ref{def:expanded}. So
$\Var(C)$ is the image of the scope, and $|{\Var(C)}|$ may be smaller than the
arity.
\end{hazard}

\begin{definition}[Relations defined by an instance]\label{def:instance-relation}
For $x_{1},\dots,x_{m}\in\Var(\mathcal I)$, $\mathcal I(x_{1},\dots,x_{m})$ is
the set of $(a_{1},\dots,a_{m})$ such that $\mathcal I$ has a solution with
$s(x_{i})=a_{i}$ for all $i$. It is a subuniverse of
$\mathbf D_{x_{1}}\times\dots\times\mathbf D_{x_{m}}$, being defined by a
primitive positive formula over the constraint relations.
\end{definition}

\subsection{Reductions}

\begin{definition}[Reduction]\label{def:reduction}
A \emph{reduction} for $\mathcal I$ is a map $D^{(\top)}$ assigning to each
$x\in\Var(\mathcal I)$ a subuniverse $D^{(\top)}_{x}\subseteq D_{x}$, possibly
empty. It is \emph{nonempty} if every $D^{(\top)}_{x}\neq\varnothing$. The
instance $\mathcal I^{(\top)}$ is $\mathcal I$ with each domain replaced by
$D^{(\top)}_{x}$ and each constraint relation intersected with the
corresponding box. For reductions we write $D^{(\bot)}\lll D^{(\top)}$ and
$D^{(\bot)}\subte{T}D^{(\top)}$ when the corresponding relation holds at every
variable.
\end{definition}

\begin{hazard}[Empty reductions]\label{haz:empty-reduction}
Reductions are allowed to be empty at some or all variables --- the maximal
$1$-consistent reduction below $D^{(B,\top)}$ in the proof of
Theorem~\ref{thm:main-inductive} is constructed exactly so that it may come out
empty, and the case split on whether it does is the branch point of the whole
argument. But $D^{(\bot)}\lll D^{(\top)}$ requires nonemptiness at every
variable (Convention~\ref{conv:empty}). Both notions are needed and they must
be different objects.
\end{hazard}

\subsection{Consistency}

\begin{definition}[Paths]\label{def:path}
A \emph{path} in $\mathcal I$ is an alternating sequence
$z_{1}-C_{1}-z_{2}-\dots-C_{l-1}-z_{l}$ with $z_{i},z_{i+1}\in\Var(C_{i})$.
It \emph{connects $b$ and $c$} if there are $a_{i}\in D_{z_{i}}$ with
$a_{1}=b$, $a_{l}=c$, and $(a_{i},a_{i+1})\in\proj_{z_{i},z_{i+1}}(C_{i})$ for
every $i$.
$\mathcal I$ is a \emph{tree-instance} if it has no path with $l\ge 3$,
$z_{1}=z_{l}$ and $C_{1},\dots,C_{l-1}$ pairwise distinct.
\end{definition}

\begin{definition}[Consistency]\label{def:consistency}
$\mathcal I$ is \emph{$1$-consistent} if $\proj_{z}(C)=D_{z}$ for every
constraint $C$ and every $z\in\Var(C)$. A reduction $D^{(\top)}$ is
$1$-consistent for $\mathcal I$ if $\mathcal I^{(\top)}$ is $1$-consistent.
$\mathcal I$ is \emph{cycle-consistent} if it is $1$-consistent and every path
from $z$ to $z$ connects $a$ to $a$, for every variable $z$ and every
$a\in D_{z}$.
\end{definition}

\begin{definition}[Linked, fragmented, irreducible]\label{def:instance-linked}
$\mathcal I$ is \emph{linked} if for every $z\in\Var(\mathcal I)$ and every
$a,b\in D_{z}$ some path from $z$ to $z$ connects $a$ and $b$.
$\mathcal I$ is \emph{fragmented} if $\Var(\mathcal I)$ splits into two disjoint
nonempty sets $X_{1},X_{2}$ with $\Var(C)\subseteq X_{1}$ or
$\Var(C)\subseteq X_{2}$ for every $C$.
$\mathcal I$ is \emph{irreducible} if there is no instance $\mathcal I'$ with
\begin{enumerate}\romanlabels\tightlist
\item $\Var(\mathcal I')\subseteq\Var(\mathcal I)$;
\item every constraint of $\mathcal I'$ is a projection of a constraint of
  $\mathcal I$ onto some of its variables;
\item $\mathcal I'$ is not fragmented;
\item $\mathcal I'$ is not linked;
\item the solution set of $\mathcal I'$ is not subdirect.
\end{enumerate}
\end{definition}

\begin{hazard}[Irreducibility is a $\Pi_{1}$ condition over a large family]
\label{haz:irreducible-instance}
Conditions (i)--(v) quantify over \emph{all} instances built from projections
of constraints of $\mathcal I$ --- a finite but exponentially large family
(any multiset of projections of any constraints onto any subsets of their
variables). It is finite, so the predicate is decidable, but it is not
something to unfold. Every use in \S\ref{sec:main-statements} goes through
Lemma~\ref{lem:crucial-irreducible} and \emph{never} through the definition.
Formalize it once, prove the two consequences, and never unfold it again.

The typo in the source at this point (``$\Var(C)\subseteq X_{1}$ or
$\Var(C)\subseteq X_{1}$'') is corrected above to $X_{2}$.
\end{hazard}

\subsection{Weakening and cruciality}

\begin{definition}[Weakening]\label{def:weakening}
A constraint $R_{1}(y_{1},\dots,y_{t})$ is \emph{weaker or equivalent to}
$R_{2}(z_{1},\dots,z_{s})$ if $\{y_{i}\}\subseteq\{z_{j}\}$ and
$R_{2}(\bar z)$ implies $R_{1}(\bar y)$; it is \emph{weaker} if in addition
$R_{1}(\bar y)$ does not imply $R_{2}(\bar z)$. \emph{Weakening} a constraint
$C$ in $\mathcal I$ means replacing $C$ by all constraints weaker than it. An
instance $\mathcal I'$ is a \emph{weakening} of $\mathcal I$ if
$\Var(\mathcal I')\subseteq\Var(\mathcal I)$ and every constraint of
$\mathcal I'$ is weaker or equivalent to a constraint of $\mathcal I$.
\end{definition}

\begin{hazard}[``Replace by all weaker constraints'']\label{haz:weakening}
Read literally this is an infinite operation: there are $2^{|D|^{t}}$ relations
weaker than a given one on a given scope, and one must also range over all
sub-scopes. It is finite for a fixed instance, but the resulting instance is
astronomically large, and no algorithm performs this operation --- it appears
only inside proofs, as a device for reaching a crucial instance
(Remark~\ref{rem:get-crucial}).

Two consequences for a formalization. First, weakening must be defined as a
\emph{relation between instances} (``$\mathcal I'$ is a weakening of
$\mathcal I$''), not as a function producing one. Second, the descending chain
argument in Remark~\ref{rem:get-crucial} needs a well-founded measure; the
natural one is the total number of tuples in all constraint relations, which
strictly decreases when a constraint is properly weakened, together with the
fact that only finitely many scopes are available. State that measure
explicitly --- ``we cannot weaken forever'' occurs three times in
\S\ref{sec:main-statements} and is the kind of step that silently fails to be
well-founded if the scope is allowed to grow.
\end{hazard}

\begin{definition}[Crucial]\label{def:crucial}
Let $D^{(\top)}$ be a reduction for $\mathcal I$. A variable $y_{i}$ of a
constraint $R(y_{1},\dots,y_{t})$ is \emph{dummy} if $R$ does not depend on its
$i$-th coordinate. A constraint $C\in\mathcal I$ is \emph{crucial in
$D^{(\top)}$} if it has no dummy variables, $\mathcal I^{(\top)}$ has no
solution, and weakening $C$ in $\mathcal I$ yields an instance with a solution
in $D^{(\top)}$. The instance $\mathcal I$ is \emph{crucial in $D^{(\top)}$} if
it has at least one constraint and all of its constraints are crucial in
$D^{(\top)}$.
\end{definition}

\begin{remark}[Reaching a crucial instance]\label{rem:get-crucial}
If $\mathcal I^{(\top)}$ has no solution, then iteratively replacing each
constraint by all weaker constraints without dummy variables terminates in an
instance crucial in $D^{(\top)}$. Every relation so introduced is still a
subuniverse of the corresponding product of domains.
\end{remark}

\subsection{Coverings}\label{sec:coverings}

\begin{definition}[Expanded covering]\label{def:expanded}
$\mathcal I'$ is an \emph{expanded covering} of $\mathcal I$, written
$\mathcal I'\in\Expanded(\mathcal I)$, if there is
$S\colon\Var(\mathcal I')\to\Var(\mathcal I)$ with
\begin{enumerate}\romanlabels\tightlist
\item $S(x)=x$ for $x\in\Var(\mathcal I)\cap\Var(\mathcal I')$;
\item $D_{x}=D_{S(x)}$ for every $x\in\Var(\mathcal I')$;
\item for every constraint $R(x_{1},\dots,x_{m})$ of $\mathcal I'$, either
  $S(x_{1}),\dots,S(x_{m})$ are pairwise distinct and
  $R(S(x_{1}),\dots,S(x_{m}))$ is weaker or equivalent to a constraint of
  $\mathcal I$, or $S(x_{1})=\dots=S(x_{m})$ and
  $\{(a,\dots,a):a\in D_{x_{1}}\}\subseteq R$.
\end{enumerate}
It is a \emph{covering} if moreover $R(S(x_{1}),\dots,S(x_{m}))\in\mathcal I$
for every constraint; a \emph{tree-covering} if it is a covering and a
tree-instance. $S(x)$ is the \emph{parent} of $x$, and $x$ a \emph{child}.
\end{definition}

\begin{proposition}[Basic properties]\label{prop:covering-basics}
Let $\mathcal I'$ be an expanded covering of $\mathcal I$.
\begin{enumerate}\alphlabels\tightlist
\item Replacing every $x$ by $S(x)$ and deleting the constraints
  $R(x,\dots,x)$ yields a weakening of $\mathcal I$.
\item A weakening is an expanded covering with $S=\mathrm{id}$.
\item Every solution of $\mathcal I$ lifts to a solution of $\mathcal I'$.
\item If $\mathcal I$ is $1$-consistent and $\mathcal I'$ is a tree-covering,
  the solution set of $\mathcal I'$ is subdirect.
\item The union of two expanded coverings is an expanded covering.
\item An expanded covering of an expanded covering is an expanded covering.
\item An expanded covering of a cycle-consistent irreducible instance is
  cycle-consistent and irreducible.
\item Every reduction of $\mathcal I$ extends to $\mathcal I'$, and a
  $1$-consistent one extends to a $1$-consistent one.
\end{enumerate}
\end{proposition}

\begin{hazard}[Variable renaming is the hidden cost]\label{haz:renaming}
Expanded coverings are where a formalization pays. The paper writes ``we
rename the variables so that the only common variable of $\Upsilon_{x}$ and
$\mathcal J$ is $x$'', and forms conjunctions
$\bigwedge_{\mathcal J\in\Omega}\bot(\mathcal J)$ of instances whose variable
sets must first be made disjoint. Doing this with a concrete variable type
means a fresh-name supply and a renaming operation, with a lemma that renaming
preserves each of $1$-consistency, cycle-consistency, cruciality, and being a
covering --- eight lemmas that the paper does not state.

The cheaper design is to make the variable type a parameter and build instances
by \emph{disjoint union} of variable types: $\Upsilon_{x}$ and $\mathcal J$ are
placed over $V_{1}\oplus V_{2}$ quotiented by identifying the single shared
variable. Renaming then disappears into the coproduct, and the eight lemmas
become transport along an equivalence. This is the single most important
representation decision in \S\ref{sec:instances}, and it should be made before
anything downstream is written.
\end{hazard}

\subsection{Induced congruences and connectedness}

\begin{definition}[$\ConOne$]\label{def:conone}
For $R$ of arity $m$ and $i\in[m]$, $\ConOne(R,i)$ is the binary relation
\[
  \{(y,y')\;:\;\exists\,\bar x\ \ R(\bar x[i\mapsto y])\wedge R(\bar x[i\mapsto y'])\}.
\]
For a constraint $C=R(x_{1},\dots,x_{m})$ we write $\ConOne(C,x_{i})$ for
$\ConOne(R,i)$. For an instance, $\Con(\mathcal I,x)=\{\ConOne(C,x):C\in
\mathcal I\}$ and $\Con(\mathcal I)=\bigcup_{x}\Con(\mathcal I,x)$.
\end{definition}

\begin{lemma}\label{lem:conone-basics}
The $i$-th variable of $R$ is rectangular if and only if $R$ is stable under
$\ConOne(R,i)$. If $R$ is subdirect and its $i$-th variable is rectangular,
then $\ConOne(R,i)$ is a congruence.
\end{lemma}

\begin{hazard}\label{haz:conone-not-congruence}
$\ConOne(R,i)$ is always reflexive and symmetric but is \emph{not} transitive
in general, hence not always a congruence; Lemma~\ref{lem:conone-basics} is
what makes it one, and its hypotheses (subdirect \emph{and} rectangular at $i$)
are both needed. The notation $\ConOne$ invites the reader to assume
congruence-hood; a formalization should give the raw relation a neutral name
and reserve $\ConOne$ for the case where the lemma applies. Note also
$\ConOne(C,x_{i})$ is ill-defined if $x$ occurs twice in the scope: it depends
on the \emph{position} $i$, not the variable.
\end{hazard}

\begin{definition}[Types of relations and instances]\label{def:rel-type}
$R$ is \emph{of PC type} (resp.\ \emph{linear type}) if $R$ is rectangular and
every $\ConOne(R,i)$ is a PC (resp.\ linear) congruence. An instance is of PC
or linear type if all its constraints are.
\end{definition}

\begin{definition}[Adjacency, connectedness]\label{def:connected}
Two congruences $\sigma_{1},\sigma_{2}$ on $\mathbf D_{x}$ are \emph{adjacent}
if there is a reflexive bridge from $\sigma_{1}$ to $\sigma_{2}$. Two
rectangular constraints $C_{1},C_{2}$ are \emph{adjacent in a common variable
$x$} if $\ConOne(C_{1},x)$ and $\ConOne(C_{2},x)$ are adjacent. An instance
$\mathcal I$ is \emph{connected} if all its constraints are rectangular, all
congruences in $\Con(\mathcal I)$ are irreducible, and the graph on the
constraints with edges the adjacent pairs is connected.
\end{definition}

Note that every proper congruence is adjacent to itself, via
$\delta(x_{1},x_{2},x_{3},x_{4})=\sigma(x_{1},x_{3})\wedge\sigma(x_{2},x_{4})$.
\clearpage
%% ---------------------------------------------------------------- ch5-properties.tex
\section{Properties of strong subuniverses}\label{sec:properties}

This section is Zhuk's \S2.3. It is the \emph{interface} between the algebraic
theory and the CSP argument: \S\ref{sec:main-statements} uses the results
below and nothing else from \S\S\ref{sec:types}. Proofs are deferred to \S5 of \cite{ZhukSimplified}, which is not
reproduced here; \S\ref{sec:formalization} budgets it.

Its shape is worth naming. Four of the five groups say that the relation
$\lll$ and the types $T$ \emph{propagate}: forward and backward along
surjective homomorphisms, down to quotients, across products of a congruence,
and through subdirect relations. The fifth --- Theorem~\ref{thm:stable-intersection}
--- is the one genuinely hard statement, and it is what converts a failure of
consistency into a bridge. Everything the CSP proof does with bridges traces
back to it.

\begin{convention}\label{conv:type-quantifier}
Where a type is not specified, $T$ ranges over
$\{\TBA,\TC,\TS,\TPC,\TL,\TD\}$. Where a multi-type $\MT T$ appears,
$T$ ranges over $\{\TPC,\TL,\TD\}$.
\end{convention}

\subsection{Ubiquity}

\begin{lemma}[Ubiquity]\label{lem:ubiquity}
If $B\lll A$ and $|B|>1$, then $C\subt[A]{T}B$ for some $C$ and some
$T\in\{\TBA,\TC,\TL,\TPC\}$.
\end{lemma}

Lemma~\ref{lem:ubiquity} is the formal content of Informal
Claim~\ref{ic:existence}: every domain of size at least $2$ that has been
reached by a chain of strong reductions admits another strong reduction. It
is what guarantees the outer loop of the algorithm makes progress, and it is
the reason the algorithm terminates with singleton domains or a linear
instance.

\begin{hazard}\label{haz:ubiquity-hyp}
The hypothesis is $B\lll A$, not $B\le A$. A subalgebra reached by no chain
of strong reductions need admit nothing. In the algorithm this hypothesis is
maintained as an invariant of the reduction loop; in a formalization it should
be carried in the state, not re-derived.
\end{hazard}

\subsection{Propagation}

\begin{longlemma}[Propagation along a homomorphism]\label{lem:propagation}
Let $f\colon\mathbf A\to\mathbf A'$ be a surjective homomorphism. Then
\begin{enumerate}\tightlist
\item[\textup{(f)}] $C\lll^{A}B\Rightarrow f(C)\lll f(B)$;
\item[\textup{(b)}] $C'\lll^{A'}B'\Rightarrow f^{-1}(C')\lll f^{-1}(B')$;
\item[\textup{(ft)}] $C\subt[A]{T(\sigma)}B\lll A\Rightarrow
  \bigl(f(C)=f(B)$ or $f(C)\subt{\TS}f(B)$ or
  $f(C)\subt[A']{T}f(B)\bigr)$;
\item[\textup{(bt)}] $C'\subt[A']{T(\sigma)}B'\Rightarrow
  f^{-1}(C')\subt[A]{T(f^{-1}(\sigma))}f^{-1}(B')$;
\item[\textup{(fs)}] $T\in\{\TBA,\TC,\TS\}$ and $C\subt{T}B$
  $\Rightarrow f(C)\subte{T}f(B)$;
\item[\textup{(fm)}] $C\subte[A]{\MT T}B\lll A$ and $f(B)$ $\TS$-free
  $\Rightarrow f(C)\subte[A']{\MT T}f(B)$;
\item[\textup{(bm)}] $C'\subte[A']{\MT T}B'\lll A'\Rightarrow
  f^{-1}(C')\subte[A]{\MT T}f^{-1}(B')$.
\end{enumerate}
\end{longlemma}

\begin{hazard}[The three-way disjunction in (ft)]\label{haz:ft}
Clause (ft) is the characteristic shape of this theory and the reason it is
usable. Pushing a strong subset forward along a homomorphism can
\begin{itemize}\tightlist
\item collapse it ($f(C)=f(B)$: the reduction becomes vacuous),
\item degrade it ($f(C)\subt{\TS}f(B)$: still strong, but the type is lost),
\item or preserve it ($f(C)\subt[A']{T}f(B)$: same type).
\end{itemize}
The type $\TS$ exists to be the value of the middle case. Every downstream
proof that pushes a reduction through a quotient does a three-way case split
here, and the middle case is the one that is easy to forget; in \cite{ZhukJACM}
it is the source of the ``go forward and backward from global to local''
induction that this paper eliminates.

Note the asymmetry: the backward clauses (b), (bt), (bm) are clean
implications with no disjunction. Preimages behave; images do not.
\end{hazard}

\begin{corollary}[Propagation to a quotient]\label{cor:prop-quotient}
Let $\delta$ be a congruence on $\mathbf A$. Then
\textup{(f)} $C\lll^{A}B\Rightarrow C/\delta\lll^{A/\delta}B/\delta$;
\textup{(t)} $C\subt[A]{T(\sigma)}B\lll A\Rightarrow
 \bigl(C/\delta=B/\delta$ or $C/\delta\subt{\TS}B/\delta$ or
 $C/\delta\subt[A/\delta]{T}B/\delta\bigr)$;
\textup{(s)} for $T\in\{\TBA,\TC,\TS\}$, $C\subt{T}B\Rightarrow
 C/\delta\subte{T}B/\delta$;
\textup{(m)} $C\subte[A]{\MT T}B\lll A$ and $B/\delta$ $\TS$-free
 $\Rightarrow C/\delta\subte[A/\delta]{\MT T}B/\delta$.
\end{corollary}

\begin{corollary}[Reflection from a quotient]\label{cor:prop-from-factor}
Let $\delta$ be a congruence on $\mathbf A$ and $B,C\le\mathbf A$. Then
$C/\delta\lll^{A/\delta}B/\delta\iff C\circ\delta\lll^{A}B\circ\delta$, and
$C/\delta\subt[A/\delta]{T}B/\delta\iff C\circ\delta\subt[A]{T}B\circ\delta$.
\end{corollary}

\begin{corollary}[Multiplying by a congruence]\label{cor:prop-times-cong}
Let $\delta$ be a congruence on $\mathbf A$. Then
\textup{(f)} $C\lll^{A}B\Rightarrow C\circ\delta\lll^{A}B\circ\delta$;
\textup{(t)} $C\subt[A]{T(\sigma)}B\lll^{A}A\Rightarrow
 \bigl(C\circ\delta=B\circ\delta$ or $C\circ\delta\subt[A]{\TS}B\circ\delta$
 or $C\circ\delta\subt[A]{T}B\circ\delta\bigr)$;
\textup{(e)} if $\delta\subseteq\sigma$ and $C\subt[A]{T(\sigma)}B\lll^{A}A$
 then $C\circ\delta\subt[A]{T}B\circ\delta$.
\end{corollary}

\begin{longcorollary}[Propagation to relations]\label{cor:prop-relations}
Let $R\sdle\mathbf A_{1}\times\dots\times\mathbf A_{m}$ and $B_{i}\lll A_{i}$.
Then
\begin{enumerate}\tightlist
\item[\textup{(r)}] $R\cap(B_{1}\times\dots\times B_{m})\dlll R$;
\item[\textup{(r1)}] $\proj_{1}\bigl(R\cap(B_{1}\times\dots\times B_{m})\bigr)
  \dlll A_{1}$;
\item[\textup{(b)}] if $C_{i}\lll^{A_{i}}B_{i}$ for all $i$ then
  $R\cap(C_{1}\times\dots\times C_{m})\dlll^{R}R\cap(B_{1}\times\dots\times B_{m})$;
\item[\textup{(b1)}] under the same hypothesis,
  $\proj_{1}\bigl(R\cap\prod C_{i}\bigr)\dlll^{A_{1}}
   \proj_{1}\bigl(R\cap\prod B_{i}\bigr)$;
\item[\textup{(m)}] if $C_{i}\subte[A_{i}]{\MT T}B_{i}$ for all $i$ then
  $R\cap\prod C_{i}\dsubte[R]{\MT T}R\cap\prod B_{i}$;
\item[\textup{(m1)}] if in addition $\proj_{1}(R\cap\prod B_{i})$ is
  $\TS$-free, then $\proj_{1}(R\cap\prod C_{i})
   \dsubte[A_{1}]{\MT T}\proj_{1}(R\cap\prod B_{i})$.
\end{enumerate}
\end{longcorollary}

For the two absorption types the conclusion is stronger --- no chain, no dot on
the type, and no subdirectness hypothesis:

\begin{lemma}[{\cite[Cor.~6.1.2, 6.9.2]{ZhukStrong}}]\label{lem:ba-center-implies}
Let $R\le\mathbf A_{1}\times\dots\times\mathbf A_{m}$ with
$\proj_{1}(R)=A_{1}$, and let $C_{i}\subte{T}A_{i}$ for every $i$, where
$T\in\{\TBA,\TC\}$ --- and, when $T=\TBA$, with a \emph{single} binary term $t$
witnessing all $m$ absorptions. Then
$\proj_{1}\bigl(R\cap(C_{1}\times\dots\times C_{m})\bigr)\dsubte{T}A_{1}$,
again with the same $t$ when $T=\TBA$.
\end{lemma}

\begin{hazard}[The source's restatement is false; this is the repair]
\label{haz:ba-center-implies}
Lemma~\ref{lem:ba-center-implies} is the workhorse of case (2) of
Theorem~\ref{thm:main-inductive}, and \cite[Lemma~19]{ZhukSimplified} restates
it from \cite{ZhukStrong} \emph{dropping two hypotheses}: the subdirectness
$\proj_{1}(R)=A_{1}$, and --- in the $\TBA$ case --- that the $C_{i}$ absorb
with a common term. Both are present in the cited source
(\cite[Cor.~6.1.2]{ZhukStrong} reads ``Suppose $R\le A_{1}\times\dots\times
A_{n}$, $\proj_{1}(R)=A_{1}$, $B_{i}\le_{\TBA(t)}A_{i}$ for every $i$''), and
without the first the statement is false.

\emph{Counterexample to the printed form.} Take $R=\{(0,0)\}\le\mathbf Z_{p}
\times\mathbf Z_{p}$ --- a subuniverse, by idempotence --- and $C_{i}=A_{i}$,
which is allowed because $\subte{T}$ admits equality. Then
$R\cap(C_{1}\times C_{2})=R$ and $\proj_{1}(R)=\{0\}$, which is neither empty
nor all of $\mathbf Z_{p}$, so the conclusion asserts $\{0\}\subt{\TBA}
\mathbf Z_{p}$ --- contradicting Lemma~\ref{lem:no-abs-in-linear}, the paper's
own Lemma~29. The failure is exactly the missing subdirectness: $\proj_{1}(R)
=\{0\}\neq\mathbf Z_{p}$.

The repair is to restore both hypotheses, as above. A formalization must also
check that they hold at each of the call sites --- in particular in the proof of
Theorem~\ref{thm:main-inductive}(2), where the lemma is applied to the relation
computing $E_{2}$ from $E_{1}$. The common-term condition is the reason
\cite{ZhukStrong} writes $\le_{\TBA(t)}$ with the term in the notation, and
carrying $t$ in the type is the cheapest way to keep it.

This is the single highest-value import to formalize first: small,
self-contained, used everywhere, and adjacent to what the predecessor project
already proves.
\end{hazard}

\subsection{Intersections}

\begin{lemma}[Intersecting with a strong subset]\label{lem:intersect-all}
If $B\lll A$ and $D\lll A$ then \textup{(i)} $B\cap D\dlll A$, and
\textup{(t)} $C\subt[A]{T(\sigma)}B\Rightarrow C\cap D\dsubte[A]{T(\sigma)}B\cap D$.
\end{lemma}

The next theorem is the heart of \S\ref{sec:properties}. Read it as follows:
if a family of strong reductions is \emph{jointly} inconsistent but
\emph{minimally} so --- dropping any one of them restores consistency --- then
the types must agree, and in the linear case one gets bridges between the
witnessing congruences. Bridges are produced nowhere else.

\begin{longtheorem}[Stable intersection]\label{thm:stable-intersection}
Suppose
\begin{enumerate}\romanlabels\tightlist
\item $C_{i}\subt[A]{T_{i}(\sigma_{i})}B_{i}\lll A$ with
  $T_{i}\in\{\TBA,\TC,\TS,\TL,\TPC\}$, for $i\in[m]$, $m\ge 2$;
\item $\bigcap_{i\in[m]}C_{i}=\varnothing$;
\item $B_{j}\cap\bigcap_{i\neq j}C_{i}\neq\varnothing$ for every $j\in[m]$.
\end{enumerate}
Then one of:
\begin{enumerate}\tightlist
\item[\textup{(ba)}] $T_{1}=\dots=T_{m}=\TBA$;
\item[\textup{(l)}] $T_{1}=\dots=T_{m}=\TL$, and for all $k,\ell\in[m]$ there
  is a bridge $\delta$ from $\sigma_{k}$ to $\sigma_{\ell}$ with
  $\bcol{\delta}=\sigma_{k}\circ\sigma_{\ell}$;
\item[\textup{(c)}] $m=2$ and $T_{1}=T_{2}=\TC$;
\item[\textup{(pc)}] $m=2$, $T_{1}=T_{2}=\TPC$ and $\sigma_{1}=\sigma_{2}$.
\end{enumerate}
\end{longtheorem}

\begin{longcorollary}[Stable intersection, relational form]\label{cor:stable-intersection}
Suppose $R\sdle\mathbf A_{1}\times\dots\times\mathbf A_{m}$,
$C_{i}\subt[A_{i}]{T_{i}(\sigma_{i})}B_{i}\lll A_{i}$ with
$T_{i}\in\{\TBA,\TC,\TS,\TL,\TPC\}$ and $m\ge 2$,
$R\cap(C_{1}\times\dots\times C_{m})=\varnothing$, and
$R\cap(C_{1}\times\dots\times B_{j}\times\dots\times C_{m})\neq\varnothing$ for
every $j$. Then one of:
\textup{(ba)} all $T_{i}=\TBA$;
\textup{(l)} all $T_{i}=\TL$ and for all $k,\ell$ a bridge $\delta$ from
$\sigma_{k}$ to $\sigma_{\ell}$ with
$\bcol{\delta}=\sigma_{k}\circ\proj_{k,\ell}(R)\circ\sigma_{\ell}$;
\textup{(c)} $m=2$ and $T_{1}=T_{2}=\TC$;
\textup{(pc)} $m=2$, $T_{1}=T_{2}=\TPC$,
$\mathbf A_{1}/\sigma_{1}\cong\mathbf A_{2}/\sigma_{2}$, and
$\{(a/\sigma_{1},b/\sigma_{2}):(a,b)\in R\}$ is bijective.
\end{longcorollary}

\begin{remark}\label{rem:duplicate-coordinate}
Several applications need more than one restriction on a single coordinate of
$R$. The statement is kept simple; to apply it, duplicate the coordinate (form
$R'=\{(\bar a,a_{k}):\bar a\in R\}$) and restrict the copies separately.
\end{remark}

\begin{hazard}[Minimality is a hypothesis, not a construction]\label{haz:minimality}
Hypotheses (ii)--(iii) of Theorem~\ref{thm:stable-intersection} say the family
is \emph{critically} inconsistent. In applications one starts with an
inconsistent family and shrinks it to a critical subfamily. That shrinking step
is left implicit in the paper (``let $\mathcal M$ be a minimal set of children
of $z$ we need to restrict''). It is easy --- a minimal subset of a finite set
with a monotone property --- but it must be stated as a lemma, because it is
performed at least four times and each time the property being minimised is
slightly different. State once: \emph{if $P$ is a monotone predicate on subsets
of a finite set $X$ and $P(X)$ holds, there is $M\subseteq X$ minimal with
$P(M)$}, and instantiate.
\end{hazard}

\subsection{Multi-types}

\begin{lemma}\label{lem:multitype-chain}
If $C\subte[A]{\MT T}B$ then $C\subt[A]{T}\dots\subt[A]{T}B$ and
$C\lll^{A}B$.
\end{lemma}

That is: an intersection of type-$T$ subsets, though not itself obtained by a
single reduction of type $T$, is reachable by a \emph{chain} of type-$T$
reductions. This is what makes $\MT T$ compatible with $\lll$, and it is used
whenever the algorithm reduces several variables at once.

\begin{lemma}[Preserving linkedness]\label{lem:preserve-linked}
Let $R\sdle\mathbf A_{1}\times\mathbf A_{2}$,
$C_{i}\subte[A_{i}]{\TMD}B_{i}\lll A_{i}$ for $i=1,2$, let $S$ be the
rectangular closure of $R$, and suppose $R\cap(B_{1}\times B_{2})\neq
\varnothing$ and $S\cap(C_{1}\times C_{2})\neq\varnothing$. Then
$R\cap(C_{1}\times C_{2})\neq\varnothing$.
\end{lemma}

\begin{lemma}[Maximal multi-type extension]\label{lem:maximal-mult}
Let $C_{1}\subt[A]{\MT T}B_{1}\lll A$, $B_{2}\lll A$,
$C_{1}\cap B_{2}=\varnothing$, $B_{1}\cap B_{2}\neq\varnothing$, and let
$\sigma$ be a maximal congruence on $\mathbf A$ with
$(C_{1}\circ\sigma)\cap B_{2}=\varnothing$. Then
$\sigma=\omega_{1}\cap\dots\cap\omega_{s}$ where each $\omega_{i}$ is a
congruence of type $T$ on $\mathbf A$ with $\cover{\omega_{i}}\supseteq
B_{1}^{2}$.
\end{lemma}

\subsection{Auxiliary imports}

The following are used but not proved in \cite{ZhukSimplified}.

\begin{lemma}[{\cite[Lem.~4.7]{MarotiMcKenzie}}]\label{lem:special-wnu}
If $w$ is an idempotent WNU operation of arity $n$ on a finite set $A$, then
$\Clo(w)$ contains a special idempotent WNU operation of arity $n^{n!}$.
\end{lemma}

\begin{lemma}[{\cite[Cor.~8.17.1]{ZhukJACM}}]\label{lem:building-perfect}
If $\sigma$ is an irreducible congruence on $\mathbf A\in\Vn$ and $\delta$ is a
bridge from $\sigma$ to $\sigma$ with $\bcol{\delta}=A^{2}$, then $\sigma$ is a
perfect linear congruence.
\end{lemma}

\begin{lemma}\label{lem:perfect-from-linked}
If $\sigma$ is an irreducible congruence on $\mathbf A\in\Vn$ and $\delta$ is a
bridge from $\sigma$ to $\sigma$ with $\bcol{\delta}$ linked, then $\sigma$ is
a perfect linear congruence.
\end{lemma}

\begin{proof}
Let $\delta^{-1}(y_{1},y_{2},x_{1},x_{2})=\delta(x_{1},x_{2},y_{1},y_{2})$,
which is again a bridge from $\sigma$ to $\sigma$. Since $\bcol{\delta}$ is
linked, $(\bcol{\delta}\circ\bcol{\delta^{-1}})^{N}=A^{2}$ for $N$ large.
Composing $2N$ bridges $\delta,\delta^{-1},\dots$ by
Lemma~\ref{lem:bridge-comp} gives a bridge $\delta'$ from $\sigma$ to $\sigma$
with $\bcol{\delta'}=A^{2}$; apply Lemma~\ref{lem:building-perfect}.
\end{proof}

\begin{hazard}[``for sufficiently large $N$'']\label{haz:linked-power}
The step $(\bcol{\delta}\circ\bcol{\delta^{-1}})^{N}=A^{2}$ deserves a lemma of
its own: if a reflexive symmetric relation $\rho$ on a finite set $A$ has
connected graph, then $\rho^{N}=A^{2}$ for $N\ge|A|$. Reflexivity is what makes
the powers increase monotonically; $\bcol{\delta}$ is reflexive because
$(a,a,a,a)\in\delta$ is not assumed --- so the reflexivity must come from
$\delta$ being a bridge from $\sigma$ to $\sigma$ with $\sigma$ reflexive, via
$\proj_{1,2}(\delta)\supsetneq\sigma$. Check this: it is the sort of step where
the informal proof is right for a reason it does not give.
\end{hazard}

\begin{lemma}\label{lem:no-abs-in-linear}
If $B\le\mathbf Z_{p}\times\dots\times\mathbf Z_{p}$ then there is no
$C\subt{T}B$ with $T\in\{\TBA,\TC\}$.
\end{lemma}

\begin{proof}
It suffices to check that for every term $\tau$, if the last variable of
$\tau^{\mathbf Z_{p}}(a,\dots,a,x)$ is not dummy then it takes all values;
hence no absorbing subuniverse exists.
\end{proof}

%% ---------------------------------------------------------------- ch5b-strongproofs.tex
\section{Proofs of the strong-subalgebra theory}\label{sec:strongproofs}

This section proves the statements of \S\ref{sec:properties}. It renders
\S5 of \cite{ZhukSimplified}, the part its author calls the most technical.
Fourteen places there are wrong, incomplete, or state less than their own
proofs give; each repair is marked \textbf{[R]} and carries a note saying what
the source has. Everything else is expansion: about two dozen steps that the
source performs silently and that a formalization cannot.

Throughout, Convention~\ref{conv:standing} is in force --- every algebra is
finite, idempotent, and has a WNU term operation --- and
\S\ref{sec:legislation} fixes the readings.

\subsection{What is imported}\label{sec:imports}

The rendering is self-contained except for the following, which are proved in
the literature and used here as black boxes. Two are substantial; the rest are
routine.

\begin{longtheorem}[Absorption Theorem; {\cite[Thm.~3.11.1]{BradyNotes}}]\label{lem:linked-implies-abs}
Let $R\lneq_{sd}\mathbf A\times\mathbf B$ be a proper subdirect relation
between finite idempotent Taylor algebras. If $R$ is linked then $\mathbf A$ or
$\mathbf B$ has a proper nonempty binary absorbing or central subuniverse.
\end{longtheorem}

\begin{hazard}[Two words the source omits]\label{haz:linked-abs-proper}
\cite{ZhukSimplified} states this as ``there exists a BA or central subuniverse
on $\mathbf A$ or $\mathbf B$'', with no properness. Read literally it is
vacuous, since every algebra is a binary absorbing subuniverse of itself; every
use is the proper nonempty one (\S\ref{sec:legislation}, item
\ref{leg:empty}). This is the single heaviest import in the section.
\end{hazard}

\begin{lemma}[{\cite[Lem.~2.9]{DecidingAbsorption}}, {\cite[Lem.~6.1, Thm.~6.9]{ZhukStrong}}]\label{lem:pp-propagation}
Let $R\le A^{n}$ be defined by a pp-formula $\Phi$ in which a relation $S$
occurs, and let $R'$ be defined by the formula obtained from $\Phi$ by
replacing \emph{every} occurrence of $S$ by some $S'\subt{T}S$, where
$T\in\{\TBA,\TC\}$. Then $R'\subte{T}R$.
\end{lemma}

\begin{lemma}[{\cite[Prop.~2.14]{DecidingAbsorption}}, {\cite[Lem.~3.2]{ZhukStrong}}]\label{lem:barto-kazda}
For $B\le\mathbf A$ and $n\ge2$: $B$ absorbs $\mathbf A$ with an operation of
arity $n$ if and only if there is no $S\le A^{n}$ with $S\cap B^{n}=\varnothing$
and $S\cap(B^{i-1}\times A\times B^{n-i})\neq\varnothing$ for every $i$.
\end{lemma}

\begin{lemma}[{\cite[Lem.~6.8]{ZhukStrong}}]\label{lem:factor-type}
If $B\subte{T}\mathbf A$ with $T\in\{\TBA,\TC,\TS\}$ and $\sigma$ is a
congruence on $\mathbf A$, then $B/\sigma\subte{T}\mathbf A/\sigma$.
\end{lemma}

\begin{hazard}[The $\TBA$ and $\TS$ cases are not imports]\label{haz:factor-type}
\cite{ZhukSimplified} writes ``for $T=\TBA$ it is straightforward, for $T=\TC$
see \cite[Lem.~6.8]{ZhukStrong}, for $T=\TS$ it is just a combination''. The
first is indeed one line --- if $t(B,A)\subseteq B$ and $t(A,B)\subseteq B$ then
the same $t$ works on images --- and the third follows from the other two by
Definition~\ref{def:ba-center-free}. Only $\TC$ is imported.
\end{hazard}

\begin{lemma}[{\cite[Lem.~6.24]{ZhukStrong}}]\label{lem:power-to-base}
If $R$ is a nontrivial strong subuniverse of
$\mathbf A_{1}\times\dots\times\mathbf A_{n}$ of a type $T\neq\TPC$, then some
$\mathbf A_{i}$ has a nontrivial subuniverse of type $T$. In particular
$B\subt{T}\mathbf A^{n}$ with $T\in\{\TBA,\TC,\TS\}$ yields some
$C\subt{T}\mathbf A$.
\end{lemma}

\begin{hazard}[The source undersells its own citation]\label{haz:power-to-base}
\cite{ZhukSimplified} proves the case $T=\TS$ by ``just repeat the same proof
word to word replacing $\TBA$ by $\TS$''. No repetition is needed. The cited
lemma is already stated for every type $T\neq\TPC$, and its proof is
\emph{type-uniform}: it takes $\proj_{1}(R)$ if that is proper and otherwise a
fibre $R'=\{(a_{2},\dots,a_{n}):(a_{1},\dots,a_{n})\in R\}$, and neither
construction mentions the type --- the type enters only through the facts that
projections and fibres preserve $\TBA$ and preserve $\TC$. So a single run of
the proof produces one witness carrying \emph{every} strong type that $R$ has,
and a subuniverse that is simultaneously binary absorbing and central yields one
that is too.
\end{hazard}

\begin{longtheorem}[{\cite[Thm.~6.15]{ZhukStrong}}]\label{thm:central-relation-source}
Let $R\sdle\mathbf A\times\mathbf B$ and
$C=\{c\in A:\{c\}\times B\subseteq R\}$. Then $C$ is a central subuniverse of
$\mathbf A$, or $\mathbf B$ has a nontrivial binary absorbing subuniverse, or
$\mathbf B$ has a nontrivial \emph{projective} subuniverse.
\end{longtheorem}

\begin{hazard}[To do]\label{haz:615-todo}
Theorem~\ref{thm:central-relation-source} is the one import we are not yet
comfortable taking on faith, and it should eventually be proved here. Its
proof in \cite{ZhukStrong} is two pages and uses only the Barto--Kazda
criterion (Lemma~\ref{lem:barto-kazda}) and the theory of projective
subuniverses; none of that theory is needed anywhere else in this document,
which is why absorbing it has been deferred.
\end{hazard}

\begin{lemma}[{\cite{HobbyMcKenzie}}]\label{lem:hobby-mckenzie}
An algebra is Abelian if and only if some congruence of its square has the
diagonal as a block; and a finite algebra with a WNU term operation is Abelian
if and only if it is affine.
\end{lemma}

\begin{lemma}\label{lem:bridge-abelian}
$\mathbf A$ is Abelian if and only if there is a bridge $\delta$ from
$0_{\mathbf A}$ to $0_{\mathbf A}$ with
$\bcol\delta=\proj_{1,2}(\delta)=\proj_{3,4}(\delta)=A^{2}$.
\end{lemma}

\begin{proof}
A congruence of $\mathbf A^{2}$ with the diagonal as a block is such a bridge.
Conversely, composing a bridge with itself enough times yields a reflexive,
symmetric, transitive relation on $\mathbf A^{2}$, whose diagonal class is the
diagonal by clause (d); now apply Lemma~\ref{lem:hobby-mckenzie}.
\end{proof}

\subsection{Absorption preliminaries}\label{sec:abs-prelim}

\begin{lemma}\label{lem:ba-center-intersect}
If $B\subte{T}\mathbf A$ with $T\in\{\TBA,\TC\}$ and $C\le\mathbf A$, then
$B\cap C\dsubte{T}C$.
\end{lemma}

\begin{proof}
$C$ is defined by the pp-formula $A(x)\wedge C(x)$; replacing the occurrence of
the full unary relation $A$ by $B$ yields $B\cap C$, so
Lemma~\ref{lem:pp-propagation} applies.
\end{proof}

\begin{lemma}[No absorbing diagonal]\label{lem:no-absorbing-diagonal}
Let $\mathbf C$ be a finite idempotent algebra with $|C|\ge 2$. Then
$0_{\mathbf C}$ is not an absorbing subuniverse of $\mathbf C^{2}$, of any
arity; in particular $0_{\mathbf C}\not\subte{\TBA}\mathbf C^{2}$ and
$0_{\mathbf C}\not\subte{\TC}\mathbf C^{2}$.
\end{lemma}

\begin{proof}
Suppose $t$ is $k$-ary and absorbs at every position. Fix $i$, take
$a_{j}\in C$ for $j\neq i$ and $b,c\in C$, and apply $t$ coordinatewise to the
tuples $(a_{j},a_{j})\in 0_{\mathbf C}$ for $j\neq i$ and $(b,c)\in C^{2}$ at
position $i$. The result
$\bigl(t(\bar a;b),\,t(\bar a;c)\bigr)$ lies in $0_{\mathbf C}$, so
$t(\bar a;b)=t(\bar a;c)$. As $b,c$ were arbitrary, $t$ does not depend on its
$i$-th argument; and this holds for every $i$, so $t$ is constant. Idempotence
then forces $|C|=1$. A central subuniverse is absorbing by
Definition~\ref{def:central}, so the central case is included.
\end{proof}

\begin{lemma}[Absorption preserves linkedness, subrelation form]\label{lem:absorb-linked-sub}
Let $W\sdle\mathbf P\times\mathbf A$ be linked and let
$\varnothing\neq W'\subte{T}W$ with $T\in\{\TBA,\TC\}$. Put
$P'=\proj_{1}(W')$ and $A'=\proj_{2}(W')$. Then $W'\sdle\mathbf P'\times
\mathbf A'$ is linked.
\end{lemma}

\begin{proof}
Let $\Phi_{k}(x,y)$ be the pp-formula in one binary relation symbol obtained by
iterating $x\sim y\iff\exists z\,(W(x,z)\wedge W(y,z))$ $k$ times. Because $W$
is subdirect, $W\circ W^{-1}$ is reflexive and symmetric on $P$, so for
$k\ge|P|$ the formula $\Phi_{k}$ interpreted in $W$ defines the transitive
closure of $W\circ W^{-1}$, which is $P^{2}$ since $W$ is linked. Interpreted
in $W'$ and for the same $k\ge|P|\ge|P'|$ it defines
$\lambda:=\operatorname{LeftLinked}(W')$. By Lemma~\ref{lem:pp-propagation},
replacing every occurrence of $W$ by $W'$ gives $\lambda\subte{T}P^{2}$. Since
$\lambda\subseteq (P')^{2}\le P^{2}$, Lemma~\ref{lem:ba-center-intersect} gives
$\lambda\subte{T}(P')^{2}$.

Now $\lambda$ is a congruence on $\mathbf P'$: reflexive because $W'$ is
subdirect over $P'$, symmetric by construction, transitive because the powers
stabilise. So $\lambda\times\lambda$ is a congruence on the algebra $(P')^{2}$,
and the image of $\lambda$ under $(P')^{2}\twoheadrightarrow(P'/\lambda)^{2}$ is
exactly $0_{\mathbf P'/\lambda}$; by Lemma~\ref{lem:factor-type},
$0_{\mathbf P'/\lambda}\subte{T}(\mathbf P'/\lambda)^{2}$. Lemma
\ref{lem:no-absorbing-diagonal} forces $|P'/\lambda|=1$, i.e.\
$\lambda=(P')^{2}$, which says $W'$ is linked.
\end{proof}

\begin{hazard}[The box form is not enough]\label{haz:box-vs-sub}
\textbf{[R]} \cite{ZhukSimplified} has only the \emph{box} form of this
principle, \cite[Prop.~2.15(i)]{BartoKozik}: if $W\sdle\mathbf A_{1}\times
\mathbf A_{2}$ is linked, $B_{i}$ absorbs $\mathbf A_{i}$, and
$W\cap(B_{1}\times B_{2})$ is subdirect in $B_{1}\times B_{2}$, then
$W\cap(B_{1}\times B_{2})$ is linked. It cites that form for a restriction that
is not a box, in the proof of Lemma~\ref{lem:pc-inductive-step}, and the two
really do differ: for the meet-semilattice on $\{0,1,2\}$ with $B=\{0,1\}$, the
subalgebra $K=\{(0,0),(1,2)\}$ of $A^{2}$ meets $B^{2}$ and has
$\proj_{1}(K)\cap B=\{0,1\}$ while $\proj_{1}(K\cap B^{2})=\{0\}$.
Lemma~\ref{lem:absorb-linked-sub} is the form the paper actually needs; it also
discharges the subdirectness hypothesis of the box form for free. Checked by
exhaustive search over nine small Taylor algebras and about $10^{5}$ pairs
$(W,W')$: no counterexample.
\end{hazard}

\begin{lemma}[A WNU forbids essentially unary quotients]\label{lem:no-unary-quotient}
Let $\mathbf U$ be a finite algebra with a WNU term operation and $|U|\ge2$.
Then some term operation of $\mathbf U$ depends on more than one variable.
Consequently, if $\mathbf A$ has a WNU term operation then no
$\mathbf U\in\mathrm{HS}(\mathbf A)$ with $|U|\ge2$ is essentially unary.
\end{lemma}

\begin{proof}
Suppose every term operation of $\mathbf U$ has at most one non-dummy variable,
and let $w$ be a WNU term operation of arity $k\ge2$; it is idempotent. If $w$
has no non-dummy coordinate it is constant, and idempotence gives $|U|=1$.
Otherwise $w(\bar x)=g(x_{i})$ for a unary $g$, and $g(x)=w(x,\dots,x)=x$, so
$w$ is the $i$-th projection. The WNU identities equate the $k$ terms in which
a single $y$ sits at each position against a background of $x$'s; two of them
are $w$ with $y$ at position $i$, which evaluates to $y$, and $w$ with $y$ at
some position $\neq i$, which evaluates to $x$. Hence $x=y$ for all $x,y$ and
$|U|=1$. The consequence follows because idempotence and the WNU identities are
identities, hence hold in every member of $\mathrm{HS}(\mathbf A)$.
\end{proof}

\begin{longlemma}[Central relations]\label{lem:central-relation}
\textbf{[R]} Let $\mathbf A,\mathbf B$ satisfy Convention~\ref{conv:standing},
let $R\sdle\mathbf A\times\mathbf B$, and put
$C=\{c\in A:\{c\}\times B\subseteq R\}$. Then
\begin{enumerate}\alphlabels\tightlist
\item $C=\varnothing$, or
\item $C$ is a central subuniverse of $\mathbf A$, or
\item $\mathbf B$ has a proper nonempty binary absorbing subuniverse.
\end{enumerate}
\end{longlemma}

\begin{proof}
Apply Theorem~\ref{thm:central-relation-source}. Its first alternative gives
(a) or (b) according as $C$ is empty or not, its second gives (c). In its third,
$\mathbf B$ has a nontrivial projective subuniverse $P$; by
\cite[Lem.~3.4]{ZhukStrong}, either $P$ is a binary absorbing subuniverse of
$\mathbf B$, which is (c) since $P$ is nontrivial, or there is an essentially
unary $\mathbf U\in\mathrm{HS}(\mathbf B)$ with $|U|\ge2$, which
Lemma~\ref{lem:no-unary-quotient} forbids.
\end{proof}

\begin{hazard}[Both changes are necessary]\label{haz:central-relation}
\cite{ZhukSimplified} states this with two alternatives, (b) and (c), citing
\cite[Thm.~6.15]{ZhukStrong} --- which has three. The elimination of the third
is legitimate but is nowhere stated; it is the proof above, and Zhuk himself
runs exactly this argument on exactly this trichotomy in the proof of
\cite[Thm.~4.15]{ZhukStrong}, step (4)\,$\Rightarrow$\,(5), without ever
connecting it to Theorem~6.15.

Alternative (a) is also not optional. The two-alternative statement is true in
\cite{ZhukStrong} only because $\varnothing$ counts as central there; under
Convention~\ref{conv:empty} it is false, as
$\mathbf A=\mathbf B=\mathbf Z_{p}$ with $R=\{(x,x+1)\}$ shows --- $R$ is
subdirect, $C=\varnothing$, and $\mathbf Z_{p}$ has no nontrivial binary
absorbing subuniverse at all. Every one of the nine call sites discharges (a)
by exhibiting a point of $C$ first.
\end{hazard}

\begin{lemma}\label{lem:bacon-left}
Let $R\sdle\mathbf A\times\mathbf B$ with $\mathbf A$ BA and centre free,
$\operatorname{LeftLinked}(R)=A^{2}$, and
$C=\{c\in B:A\times\{c\}\subseteq R\}$. Then $C\neq\varnothing$ and
$C\subte{\TBA,\TC}\mathbf B$.
\end{lemma}

\begin{proof}
For $k\ge2$ put $W_{k}=\{(a_{1},\dots,a_{k}):\exists b\,\forall i\,
R(a_{i},b)\}$. If $W_{|A|}=A^{|A|}$ then $C\neq\varnothing$, since a witness
$b$ for the tuple listing all of $A$ lies in $C$. Otherwise choose $k$ minimal
with $W_{k}\neq A^{k}$. Since $\operatorname{LeftLinked}(R)=A^{2}$ we get
$\operatorname{LeftLinked}(W_{k})=A^{2}$, and viewing $W_{k}$ as a binary
relation in $\mathbf A\times\mathbf A^{k-1}$ --- subdirect by minimality of $k$
--- Theorem~\ref{lem:linked-implies-abs} and Lemma~\ref{lem:power-to-base}
produce a proper nonempty binary absorbing or central subuniverse of
$\mathbf A$, contrary to hypothesis. So $C\neq\varnothing$.

By Lemma~\ref{lem:central-relation} applied to $R^{-1}$ --- alternative (a)
being excluded --- either $C$ is central in $\mathbf B$ or $\mathbf A$ has a
proper nonempty binary absorbing subuniverse; the latter is excluded, so $C$ is
central. For binary absorption, suppose not; by Lemma~\ref{lem:barto-kazda}
there is $S\le B\times B$ with $S\cap(C\times C)=\varnothing$,
$S\cap(B\times C)\neq\varnothing$ and $S\cap(C\times B)\neq\varnothing$. Put
\[
  W=\{(a_{1},\dots,a_{|A|}):\exists(b,c)\in S\;\bigl(c\in C\wedge
     \forall i\,R(a_{i},b)\bigr)\}.
\]
By Lemma~\ref{lem:pp-propagation} $W\subt{\TC}A^{|A|}$, so by
Lemma~\ref{lem:power-to-base} $\mathbf A$ has a proper nonempty central
subuniverse --- again a contradiction.
\end{proof}

\begin{lemma}\label{lem:strong-nonempty}
If $C\lll^{A}B\lll A$ and $D\subt{\TBA,\TC}B$ then $C\cap D\neq\varnothing$.
\end{lemma}

\begin{proof}
Suppose not, and choose $C''$ minimal with
$C\lll^{A}C'\subt[A]{T(\sigma)}C''\lll^{A}B$ and $C''\cap D\neq\varnothing$.
By Lemma~\ref{lem:ba-center-intersect} $C''\cap D\subt{\TBA,\TC}C''$, so by
Lemma~\ref{lem:possible-intersections} $T=\TD$; and then
Lemma~\ref{lem:factor-type} gives
$(C''\cap D)/\sigma\subt{\TBA}C''/\sigma$, contradicting the requirement in
Definition~\ref{def:types} that $C''/\sigma$ be BA and centre free.
\end{proof}

\begin{lemma}[{\cite[Lem.~6.25]{ZhukStrong}}]\label{lem:possible-intersections}
If $B\subt{T_{1}}\mathbf A$, $C\subt{T_{2}}\mathbf A$, $B\cap C=\varnothing$
and $T_{1},T_{2}\in\{\TBA,\TC,\TS\}$, then $T_{1}=T_{2}\in\{\TBA,\TC\}$.
\end{lemma}

\subsection{The intersection property}\label{sec:intersection-property}

The goal is Lemma~\ref{lem:intersection-good}: if $B\lll A$, $C\lll A$ meet and
$\delta$ is a dividing congruence for $B$, then $(B\cap C)/\delta$ is a single
point or all of $B/\delta$. Everything in \S\ref{sec:strongproofs} that
produces a bridge goes through it.

\begin{lemma}[The shift manoeuvre]\label{lem:multiset-shift}
Let $R\le A^{k}$ be totally symmetric and closed under
$(a_{1},\dots,a_{k})\mapsto(a_{1},a_{1},a_{2},\dots,a_{k-1})$. If
$\bar a\in R$ has entry multiset $M$, then $R$ contains a tuple with entry
multiset $M+\{u\}-\{v\}$ for any $u,v\in M$ occupying distinct positions. In
particular, if $|\{$distinct entries of $\bar a\}|=r$ then $R$ contains a tuple
whose entries are any prescribed $k$-element multiset drawn from those $r$
values, provided that multiset uses at least one of them.
\end{lemma}

\begin{proof}
Total symmetry lets us place $u$ first and $v$ last; the shift then duplicates
$u$ and deletes $v$. Iterating and permuting gives every reachable multiset.
\end{proof}

\begin{hazard}[The shift deletes the last entry]\label{haz:shift}
\textbf{[R]} \cite{ZhukSimplified} uses this twice with the words ``by the
conditions of the lemma we have\dots''. A single application of the shift to
$(a,b_{2},\dots,b_{k-1},c)$ destroys $c$, so the claimed conclusion
$(a,\dots,a,c)\in R$ does not follow from one step; it needs $k-2$ steps
interleaved with permutations, which is what Lemma~\ref{lem:multiset-shift}
packages.
\end{hazard}

\begin{lemma}\label{lem:tot-sym-full}
Let $\mathbf A$ be BA and centre free and let $R\le A^{k}$ be totally
symmetric, closed under the shift, with $\proj_{1,2}(R)=A^{2}$. Then $R=A^{k}$.
\end{lemma}

\begin{proof}
Induct on $k$; for $k=2$ the hypothesis is the conclusion. Let $k>2$. By total
symmetry $\proj_{1,k}(R)=A^{2}$, so for any $a,c$ there is a tuple of $R$ with
first entry $a$ and last entry $c$; by Lemma~\ref{lem:multiset-shift} we may
replace all other entries by $a$, giving $(a,\dots,a,c)\in R$, and deleting $c$
gives $(a,\dots,a)\in R$. Hence, viewing $R\le A^{k-1}\times A$, every left
vertex $(a,\dots,a)$ is adjacent to every $c$, so $R$ is subdirect and
$\operatorname{RightLinked}(R)=A^{2}$.

The relation $R'=\proj_{1,\dots,k-1}(R)$ inherits total symmetry, has
$\proj_{1,2}(R')=A^{2}$, and is shift-closed: from $\bar a\in R'$ extend to
$R$, shift there, and project. By induction $R'=A^{k-1}$, so
$R\sdle A^{k-1}\times A$. If $R\neq A^{k}$ then
Theorem~\ref{lem:linked-implies-abs} gives a proper nonempty binary absorbing
or central subuniverse on $\mathbf A^{k-1}$ or on $\mathbf A$, and
Lemma~\ref{lem:power-to-base} moves the first case to $\mathbf A$ ---
contradicting the hypothesis either way.
\end{proof}

\begin{lemma}\label{lem:tot-sym-irred}
Let $\sigma$ be a dividing congruence for $B\le A$ and let
$R\le(\mathbf A/\sigma)^{k}$ be reflexive, totally symmetric and shift-closed.
Then $(B/\sigma)^{k}\subseteq R$, or $R$ is the set of constant tuples.
\end{lemma}

\begin{proof}
If $\proj_{1,2}(R)$ is the equality relation then, by total symmetry, all
coordinates of every tuple of $R$ agree, and reflexivity supplies all constant
tuples.

Otherwise $\proj_{1,2}(R)$ is a subuniverse of $(\mathbf A/\sigma)^{2}$
containing $0_{\mathbf A/\sigma}$ properly and stable under it, so by
minimality of the cover (Proposition~\ref{prop:cover}, transported by
Lemma~\ref{lem:irred-quotient}) it contains $\cover{\sigma}/\sigma$, hence
$(B/\sigma)^{2}$ --- the last step by Definition~\ref{def:types}(i). Now apply
Lemma~\ref{lem:tot-sym-full} to $R\cap(B/\sigma)^{k}$ over the algebra
$B/\sigma$, which is BA and centre free by Definition~\ref{def:types}(iii). Its
hypothesis $\proj_{1,2}\bigl(R\cap(B/\sigma)^{k}\bigr)=(B/\sigma)^{2}$ is not
immediate from $\proj_{1,2}(R)\supseteq(B/\sigma)^{2}$, since a witnessing
tuple may leave $B/\sigma$; Lemma~\ref{lem:multiset-shift} repairs it, pulling
$(x_{1},x_{2},y_{3},\dots,y_{k})$ back to $(x_{1},x_{2},x_{1},\dots,x_{1})$.
\end{proof}

\begin{lemma}[Full fibre]\label{lem:full-fibre}
Let $F$ be a finite set with $|F|\le N$ and let $W\subseteq F^{N}\times E$
satisfy $\proj_{1..N}(W)=F^{N}$ and be closed under coordinate substitution: if
$(x_{1},\dots,x_{N},e)\in W$ and $g:[N]\to[N]$ then
$(x_{g(1)},\dots,x_{g(N)},e)\in W$. Then $F^{N}\times\{d\}\subseteq W$ for some
$d\in E$.
\end{lemma}

\begin{proof}
Write $F=\{F_{1},\dots,F_{t}\}$ with $t\le N$ and put
$\bar x=(F_{1},\dots,F_{t},F_{t},\dots,F_{t})\in F^{N}$. By surjectivity pick
$d$ with $(\bar x,d)\in W$. Every $\bar y\in F^{N}$ is $\bar x\circ g$ for some
$g:[N]\to[N]$, since every entry of $\bar y$ occurs among the entries of
$\bar x$; closure gives $(\bar y,d)\in W$.
\end{proof}

\begin{hazard}[Why the arity is $|A|$]\label{haz:full-fibre}
\textbf{[R]} Three proofs in \cite{ZhukSimplified} contain, verbatim modulo
names, ``since $\proj_{1,\dots,|A|}(R')=(B/\delta)^{|A|}$, there exists $d$ with
$(B/\delta)^{|A|}\times\{d\}\subseteq R'$''. Surjectivity of a projection does
not produce a full fibre, and as stated this is a non sequitur. It is true
because each such $R'$ is cut out by a conjunction of unary and binary
conditions on its first $|A|$ entries, hence closed under substituting entries
for one another, and because $|B/\delta|\le|A|$. \textbf{The exponent $|A|$ is
load-bearing}: it must be at least the number of blocks, so that one tuple can
enumerate them. A rendering that normalises the arity away breaks all three
proofs.
\end{hazard}

\begin{longlemma}\label{lem:self-intersection-pc}
Let $B_{k}\subt[A]{T_{k}(\sigma_{k})}\dots\subt[A]{T_{1}(\sigma_{1})}B_{0}=A$
with $T_{i}\in\{\TBA,\TC,\TS,\TD\}$, let $\delta$ be a congruence on
$\mathbf A$, and let $m\in[k]$ with $T_{m}=\TD$. Then
$\bigl((B_{k}\circ\delta)\cap B_{m-1}\bigr)/\sigma_{m}$ is either
$B_{m-1}/\sigma_{m}$ or $B_{m}/\sigma_{m}$; and in the second case
$\sigma_{m}\supseteq\delta\cap\sigma_{1}\cap\dots\cap\sigma_{m-1}$.
\end{longlemma}

\begin{proof}[Proof sketch]
Induction on $k$, splitting on $T_{k}$ and on whether $k>m$. The absorption
cases run through Lemmas~\ref{lem:pp-propagation} and \ref{lem:factor-type} to
produce a strong subuniverse of $B_{k-1}/\sigma_{k}$ and contradict
Definition~\ref{def:types}(iii); the $\TD$ cases build the auxiliary relation
$R'$ on $(B_{k-1}/\sigma_{k})^{|A|}\times B_{n-1}/\sigma_{n}$, obtain a full
fibre from Lemma~\ref{lem:full-fibre}, and finish with
Lemma~\ref{lem:central-relation}. Two steps must be supplied at each such
finish: the centre is \emph{proper}, the witness being the displayed
non-containment two lines earlier, and the binary absorbing subuniverse
produced lives on a \emph{power} and is moved down by
Lemma~\ref{lem:power-to-base}, which the source cites at one of the three sites
and omits at the other two.
\end{proof}

\begin{longlemma}[The intersection property]\label{lem:intersection-good}
Let $B\lll A$, $C\lll A$ with $B\cap C\neq\varnothing$. Then
\begin{enumerate}\alphlabels\tightlist
\item[\textup{(d)}] if $\delta$ is a dividing congruence for $B\le A$ then
  $|(B\cap C)/\delta|=1$ or $(B\cap C)/\delta=B/\delta$; and if
  $|(B\cap C)/\delta|=1$ then $\delta\supseteq\delta_{1}\cap\dots\cap
  \delta_{s}$, where $\delta_{1},\dots,\delta_{s}$ are the congruences of the
  \emph{chosen} chains witnessing $B\lll A$ and $C\lll A$;
\item[\textup{(s)}] if $G\subt{\TBA,\TC}B$ then $G\cap C\neq\varnothing$.
\end{enumerate}
\end{longlemma}

\begin{proof}[Proof outline, with the repaired steps in full]
Fix chains $B=B_{k}<\dots<B_{0}=A$ and $C=C_{\ell}<\dots<C_{0}=A$, put
$\sigma=\bigcap_{i}\sigma_{i}$ and $\omega=\bigcap_{j}\omega_{j}$, and induct on
$k+\ell$. Every appeal to the inductive hypothesis lands at strictly smaller
$k+\ell$.

\smallskip\noindent\textbf{(s).} Split on the type $\mathcal T_{\ell}$ of the
last step of $C$'s chain. In the case $\mathcal T_{\ell}=\TD$ the source
asserts $(G\cap C_{\ell-1})/\omega_{\ell}\subt{\TBA,\TC}C_{\ell-1}/
\omega_{\ell}$ with a \emph{proper} containment, where the cited lemmas supply
only $\subte{}$. \textbf{[R]} Split the case instead: if the containment is an
equality then every $\omega_{\ell}$-block meeting $C_{\ell-1}$ meets
$G\cap C_{\ell-1}$, in particular the block $E$ with
$C_{\ell}=C_{\ell-1}\cap E$, so $G\cap C_{\ell}\neq\varnothing$, which is what
was to be proved.

\smallskip\noindent\textbf{(d).} Put
$S=\{(a_{1}/\delta,\dots,a_{|A|}/\delta): \forall i,j\;(a_{i},a_{j})\in
\sigma\cap\omega\}$, a reflexive, totally symmetric, shift-closed subuniverse of
$(\mathbf A/\delta)^{|A|}$, and apply Lemma~\ref{lem:tot-sym-irred}.

\emph{Case 1: $S$ is the diagonal.} Then $\sigma\cap\omega\subseteq\delta$,
which is the ``moreover'' clause; and any $a,b\in B\cap C$ satisfy
$(a,b)\in\sigma\cap\omega$, so $|(B\cap C)/\delta|=1$.

\emph{Case 2: $(B/\delta)^{|A|}\subseteq S$.} Refine $S$ to $S_{m,n}$ by
requiring $a_{i}\in B_{m}\cap C_{n}$, and take the least $m$ (subcase 2A) or
least $n$ (subcase 2B) at which containment fails, splitting on the type of the
corresponding step. \textbf{[R]} Three steps are supplied here that the source
does not state. \emph{Properness}: in 2A1 and 2B1 the containments
$S_{m,0}\cap(B/\delta)^{|A|}\subt{T_{m}}(B/\delta)^{|A|}$ are proper by the
minimal choice of $m$, and nonempty because the constant tuples
$(b/\delta,\dots,b/\delta)$ for $b\in B$ lie in $S_{m,0}$, since $B\subseteq
B_{m}$ and $(b,b)\in\sigma\cap\omega$. \emph{Which block}: in 2B3 the inductive
hypothesis offers ``size $1$ or all of $C_{\ell-1}/\omega_{n}$'', and the
stronger reading ``$=C_{n}/\omega_{n}$'' used there is correct because
$B\cap C\neq\varnothing$ and $C\subseteq C_{n}$ force $B_{k}\cap C_{n}\neq
\varnothing$, so the single block is the one defining $C_{n}$. \emph{Full
fibres}: the appeals in 2A3 and 2B3 to a $d$ with
$(B/\delta)^{|A|}\times\{d\}\subseteq R'$ are Lemma~\ref{lem:full-fibre}.
Subcase 2C survives and gives $(B\cap C)/\delta=B/\delta$.
\end{proof}

\begin{hazard}[What remains to be expanded]\label{haz:intersection-draft}
The proof above is written at the level of its case structure, with every
repaired step in full. Expanding the twelve subcases of (d) to the level of the
rest of this section is the largest single piece of writing left; the source
occupies 292 lines here and nothing in it is wrong. One loose end: the proof
fixes \emph{minimal} chains for $B$ and $C$ and I could find no step that
consumes minimality --- a shorter chain only makes the induction measure
smaller. Either identify the use or drop the hypothesis.
\end{hazard}

\begin{corollary}\label{cor:intersection-good}
Let $R\sdle\mathbf A_{1}\times\mathbf A_{2}$, $B_{1}\lll A_{1}$,
$B_{2}\lll A_{2}$, and let $\sigma$ be a dividing congruence for
$B_{1}\lll A_{1}$. Then $\proj_{1}\bigl(R\cap(B_{1}\times B_{2})\bigr)/\sigma$
is empty, of size $1$, or equal to $B_{1}/\sigma$.
\end{corollary}

\begin{proof}
Let $f_{i}:\mathbf R\to\mathbf A_{i}$ be the projections. By
Lemma~\ref{lem:propagation}(b), $f_{i}^{-1}(B_{i})\lll\mathbf R$; apply
Lemma~\ref{lem:intersection-good}(d) inside $\mathbf R$ to these two, with the
congruence $f_{1}^{-1}(\sigma)$, and read the conclusion back through $f_{1}$.
\end{proof}

\begin{hazard}\label{haz:cor-intersection}
\cite{ZhukSimplified} writes $\delta=f^{-1}(\sigma)$ with $f$ undefined; read
$f_{1}$. The alternative ``empty'' is present here and absent from
Lemma~\ref{lem:intersection-good}, which hypothesises $B\cap C\neq\varnothing$;
the empty case needs its own line at every call site.
\end{hazard}

\begin{lemma}[Parallelogram property for $\TD$]\label{lem:parallelogram-D}
Let $B_{1},B_{2}\lll A$, $C_{1},C_{1}'\subt[A]{\TD(\sigma_{1})}B_{1}$ and
$C_{2},C_{2}'\subt[A]{\TD(\sigma_{2})}B_{2}$ with $C_{1}'\cap C_{2}$,
$C_{1}\cap C_{2}'$ and $C_{1}'\cap C_{2}'$ all nonempty. Then
$C_{1}\cap C_{2}\neq\varnothing$.
\end{lemma}

\begin{proof}
By Lemma~\ref{lem:intersection-good}(d), $(B_{1}\cap B_{2})/\sigma_{1}$ is of
size $1$ or all of $B_{1}/\sigma_{1}$. In the first case the single block is
the one defining $C_{1}$, because $C_{1}\cap C_{2}'\neq\varnothing$ puts a
point of $C_{1}$ in $B_{1}\cap B_{2}$; hence $B_{1}\cap B_{2}\subseteq C_{1}$
and $C_{1}\cap C_{2}=B_{1}\cap C_{2}\supseteq C_{1}'\cap C_{2}\neq\varnothing$.
The symmetric argument applies if $(B_{1}\cap B_{2})/\sigma_{2}$ has size $1$.

So assume $(B_{1}\cap B_{2})/\sigma_{1}=B_{1}/\sigma_{1}$ and
$(B_{1}\cap B_{2})/\sigma_{2}=B_{2}/\sigma_{2}$, and put
$S=\{(a/\sigma_{1},a/\sigma_{2}):a\in B_{1}\cap B_{2}\}$, which is then
subdirect in $(B_{1}/\sigma_{1})\times(B_{2}/\sigma_{2})$. Apply
Lemma~\ref{lem:intersection-good}(d) to each $\sigma_{1}$-class $C$ of $B_{1}$
against $B_{2}$ with the congruence $\sigma_{2}$; the hypothesis
$C\cap B_{2}\neq\varnothing$ holds because $S$ is subdirect. Each fibre of $S$
is therefore a single point or the whole of $B_{2}/\sigma_{2}$.

If some fibre is full, then either $S$ is full --- and then the fibre of
$C_{1}/\sigma_{1}$ contains $C_{2}/\sigma_{2}$, giving $C_{1}\cap C_{2}\neq
\varnothing$ --- or the centre of $S$ is proper and nonempty, and
Lemma~\ref{lem:central-relation} produces a proper nonempty binary absorbing or
central subuniverse of $B_{1}/\sigma_{1}$ or $B_{2}/\sigma_{2}$, contradicting
Definition~\ref{def:types}(iii). If every fibre is a single point, then the
fibre of $C_{1}'/\sigma_{1}$ is both $C_{2}/\sigma_{2}$ (witnessed by
$C_{1}'\cap C_{2}$) and $C_{2}'/\sigma_{2}$ (witnessed by $C_{1}'\cap C_{2}'$),
so $C_{2}=C_{2}'$ and $C_{1}\cap C_{2}=C_{1}\cap C_{2}'\neq\varnothing$.
\end{proof}

\subsection{Bridges}\label{sec:bridge-proofs}

This subsection is where the source is least reliable: of its five proofs, one
rests on a formula that defines the wrong relation, one proves a statement that
is false, and one asserts a hypothesis it does not have. All three are cured by
a single missing property, and Lemma~\ref{lem:symmetrise} supplies it free.

\begin{hazard}[Three defects, one cause]\label{haz:bridge-repairs}
\textbf{[R]} In \cite{ZhukSimplified}:
\begin{enumerate}\alphlabels\tightlist
\item the symmetrisation opening the proof of
  Lemma~\ref{lem:block-good-bridge} is written
  $\sigma(x_{1},x_{2},x_{3},x_{4})=\exists x_{5}\exists x_{6}\,
  \delta(x_{1},x_{2},x_{5},x_{5})\wedge\delta(x_{3},x_{4},x_{5},x_{6})$.
  Clause (d) of Definition~\ref{def:bridge} turns the repeated $x_{5}$ into
  $x_{1}=x_{2}$, so $\proj_{1,2}(\sigma)\subseteq 0_{A}$ and the relation
  defined is \emph{never} a bridge, for any $\sigma_{1}$. It is a
  one-character slip: the intended relation is $\delta^{\circ}$ of
  Lemma~\ref{lem:symmetrise}, which the source writes correctly $210$ lines
  later, in the proof of Lemma~\ref{lem:linear-equiv}.
\item Lemma~\ref{lem:nice-bridge-abelian} is stated there without
  pair-reflexivity, and is then \emph{false}: on $\mathbf Z_{3}\in\mathcal
  V_{4}$ the relation
  $\delta=\{(x_{1},x_{2},x_{3},x_{4}):x_{1}-x_{2}+x_{3}-x_{4}=0\}$ satisfies
  every hypothesis, and none of the six bijections $\mathbb Z_{3}\to\mathbb
  Z_{3}$ carries it to $\{x_{1}-x_{2}=x_{3}-x_{4}\}$ (verified by exhaustion).
  The step that breaks is ``composing $\delta$ with itself we get
  $\delta_{0}\supseteq\delta$''.
\item Lemma~\ref{lem:nontrivial-bridge} asserts that
  $\proj_{1,2}(\delta\cap B^{4}\cap(\cover\sigma\times\cover\sigma))$ is linked
  ``since $\bcol\delta\supseteq\cover\sigma$''. That inclusion bears on a
  different hypothesis; linkedness does not follow, and without
  pair-reflexivity the projection can collapse to $0_{B}$.
\end{enumerate}
Pair-reflexivity (Definition~\ref{def:pair-reflexive}) repairs all three, and
costs nothing: Lemma~\ref{lem:nontrivial-bridge} is invoked exactly once in
\cite{ZhukSimplified}, on a $\delta^{\circ}$, and
Lemma~\ref{lem:nice-bridge-abelian} exactly once from inside that proof.
\end{hazard}

\begin{lemma}\label{lem:nice-bridge-abelian}
Let $\sigma$ be a congruence on $\mathbf A$ and $\delta$ a symmetric,
\emph{pair-reflexive} bridge from $\sigma$ to $\sigma$ with
$\proj_{1,2}(\delta)=\bcol\delta=A^{2}$. Then there is an abelian group
$(G;+,-)$ with $(A/\sigma;\delta/\sigma)\cong(G;x_{1}-x_{2}=x_{3}-x_{4})$.
\end{lemma}

\begin{proof}
By Lemma~\ref{lem:bridge-abelian} and Hobby--McKenzie, $\mathbf A/\sigma$ is
affine, hence polynomially equivalent to a module $\mathbf G$. Composing
$\delta$ with itself gives $\delta_{0}=\delta\ast\delta$, an equivalence
relation on $\mathbf A^{2}$, and $\delta\subseteq\delta_{0}$ by
Lemma~\ref{lem:symmetrise} --- this is the step that needs pair-reflexivity.
Then $\delta_{0}/\sigma$ is a congruence of $\mathbf G^{2}$ having the diagonal
as a block, so it is $\{((x_{1},y_{1}),(x_{2},y_{2})):x_{1}-x_{2}=y_{1}-y_{2}\}$.
Finally $\delta/\sigma$ is preserved by the Maltsev operation $x-y+z$ and has
$\bcol\delta=\proj_{1,2}(\delta)=A^{2}$, so $\proj_{1,2,3}(\delta)=A^{3}$; since
the last entry of $\delta_{0}/\sigma$ is determined by the first three, counting
gives $\delta=\delta_{0}$.
\end{proof}

\begin{remark}[What the untwisted form really is]\label{rem:phi-twist}
Without pair-reflexivity the correct conclusion is
$\delta=\{(x_{1},x_{2},x_{3},x_{4}):\varphi(x_{1}-x_{2})=x_{3}-x_{4}\}$ for an
automorphism $\varphi$ of $G$: factoring out the diagonal exhibits $\delta$ as
the preimage of a subgroup $K\le G\times G$ that meets $G\times 0$ and
$0\times G$ trivially and is subdirect, hence a graph. Symmetry says
$\varphi^{2}=\mathrm{id}$; pair-reflexivity says $\varphi=\mathrm{id}$. The
counterexample of Hazard~\ref{haz:bridge-repairs}(b) is $\varphi=-\mathrm{id}$,
which is why it needs $p$ odd.
\end{remark}

\begin{lemma}\label{lem:block-good-bridge}
Let $\delta$ be a bridge from $0_{\mathbf A}$ to $0_{\mathbf A}$ with
$\proj_{1,2}(\delta)\subseteq\bcol\delta$ and $\proj_{1,2}(\delta)$ linked.
Then $\mathbf A$ is BA and centre free.
\end{lemma}

\begin{proof}[Proof outline, with the supplied steps]
Replace $\delta$ by $\sigma:=\delta^{\circ}$ of Lemma~\ref{lem:symmetrise} and
put $\omega=\proj_{1,2}(\sigma)=\proj_{1,2}(\delta)$. Note first that
$\bcol\delta$ is \emph{reflexive}, because $0_{A}\subseteq\proj_{1,2}(\delta)
\subseteq\bcol\delta$ by clause (c); hence $\omega\subseteq\bcol\delta
\subseteq\bcol\sigma$, and $\bcol\sigma$ is symmetric, which is what legitimises
the ``without loss of generality'' below. Induct on $|A|$.

Suppose $B\subt{T}\mathbf A$ with $T\in\{\TBA,\TC\}$. Since $\omega$ is linked,
$\omega$ meets $(A\setminus B)\times B$ or $B\times(A\setminus B)$; say the
first, and pick $(a,b)$ there.

\emph{Case 1: some subalgebra $D\lneq A$ contains $a$ and $b$.} Let $E$ be the
block containing $a,b$ of the transitive closure of
$(\omega\cup\omega^{-1})\cap D^{2}$; that closure is a congruence on $\mathbf D$,
because the transitive closure of a subalgebra of $D^{2}$ is a subalgebra, so
$E$ is a subuniverse, and $\omega\cap E^{2}$ is linked by construction.
\emph{This is the existence of $E$, which the source asserts.} Moreover
$\proj_{1,2}(\sigma\cap E^{4})=\omega\cap E^{2}$: the inclusion $\subseteq$ is
trivial and $\supseteq$ holds because pair-reflexivity puts
$(x_{1},x_{2},x_{1},x_{2})$ in $\sigma$ for every $(x_{1},x_{2})\in\omega$.
So the inductive hypothesis applies to $\sigma\cap E^{4}$ and makes $E$ BA and
centre free, while Lemma~\ref{lem:ba-center-intersect} makes $B\cap E$ a proper
nonempty subuniverse of $E$ of type $T$ --- proper because $a\in E\setminus B$,
nonempty because $b\in B\cap E$. Contradiction.

\emph{Case 2: no such $D$.} Then $\{a\}\circ\omega=A$, because $\{a\}$ is a
subuniverse by idempotence, hence so is $\{a\}\circ\omega$, and it contains
$a$ and $b$. Put $\xi(x_{1},x_{2})=\sigma(a,x_{1},x_{2},b)$; from
$(a,b,a,b),(a,a,b,b)\in\sigma$ both projections of $\xi$ contain $a$ and $b$,
hence equal $A$. Put $C=B\circ\xi$, so $C\subte{T}\mathbf A$ by
Lemma~\ref{lem:pp-propagation}, with $a\in C$ and $b\notin C$ --- the latter
because $b\in C$ would give $(a,c,b,b)\in\sigma$ with $c\in B$, and clause (d)
would force $a=c\in B$.

If $T=\TBA$, applying the absorbing term to $(a,b,a,b)$ and $(a,a,b,b)$ gives
$(a,b_{1},b_{2},b)\in\sigma$ with $b_{1},b_{2}\in B$, so $C\cap B\neq
\varnothing$. If $T=\TC$, the three tuples $(a,b,a,b)$, $(a,a,a,a)$ and
$(a,a,b,b)$ witness
\[
  \sigma\cap(\{a\}\times A\times C\times B),\quad
  \sigma\cap(\{a\}\times C\times C\times A),\quad
  \sigma\cap(\{a\}\times C\times A\times B)\;\neq\;\varnothing ,
\]
which is hypothesis (4) of Corollary~\ref{cor:stable-intersection} for the
ternary relation $\{(x_{2},x_{3},x_{4}):\sigma(a,x_{2},x_{3},x_{4})\}$ with
$n=3$, all types $\TC$ and all ambients $A$ --- legitimate because $\lll$ is
reflexive. \emph{The source displays the three tight boxes
$\{a\}\times B\times C\times B$, $\{a\}\times C\times C\times C$,
$\{a\}\times C\times B\times B$ instead, which are subsets of these and so do
not match the corollary's shape; widening them is the missing step.} Since
$n=3$ and all types are $\TC$, every alternative of the corollary fails, so
hypothesis (3) must fail:
$\sigma\cap(\{a\}\times C\times C\times B)\neq\varnothing$. Put
$C'=\proj_{2}$ of that set. Then $C'\subte{\TC}A$ by
Lemma~\ref{lem:pp-propagation}, using $\{a\}\circ\omega=A$, and hence
$C'\subte{\TC}C$ by Lemma~\ref{lem:ba-center-intersect}; and by clause (d)
either $a\notin C'$ or $C\cap B\neq\varnothing$.

So in every case we have $\varnothing\neq C'\subt{T}C\subt{T}A$ --- taking
$C'=B\cap C$ when $C\cap B\neq\varnothing$. Then $|C|>1$, and applying the
inductive hypothesis to $\sigma\cap C^{4}$ contradicts $C'\subt{T}C$: its
hypotheses hold because $\{a\}\circ\omega=A$ makes $\omega\cap C^{2}$ a star
centred at $a$, hence linked.
\end{proof}

\begin{longlemma}\label{lem:nontrivial-bridge}
Let $\sigma$ be an irreducible congruence on $\mathbf A$ and $\delta$ a
\emph{pair-reflexive} bridge from $\sigma$ to $\sigma$ with
$\bcol\delta\supsetneq\sigma$. Then
\begin{enumerate}\alphlabels\tightlist
\item $\cover\sigma$ is a congruence;
\item $B/\sigma$ is BA and centre free for every block $B$ of $\cover\sigma$;
\item if $\delta$ is also symmetric then there is a prime $p$ such that every
  block $B$ of $\cover\sigma$ satisfies
  $(B/\sigma;(\delta\cap B^{4})/\sigma)\cong(\mathbb Z_{p}^{n_{B}};
   x_{1}-x_{2}=x_{3}-x_{4})$ for some $n_{B}\ge 0$.
\end{enumerate}
\end{longlemma}

\begin{proof}
By Lemma~\ref{lem:irred-quotient} we may assume $\sigma=0_{\mathbf A}$. Since
$\bcol\delta$ is a subuniverse of $\mathbf A^{2}$ properly containing $\sigma$
and stable under it, minimality gives $\cover\sigma\subseteq\bcol\delta$; the
same argument applied to $\proj_{1,2}(\delta)$, which properly contains $\sigma$
by clause (c), gives $\cover\sigma\subseteq\proj_{1,2}(\delta)$.

Let $B$ be a block of $\operatorname{LeftLinked}(\cover\sigma)$ with $|B|>1$ and
put $\delta'=\delta\cap B^{4}\cap(\cover\sigma\times\cover\sigma)$. Then
$\proj_{1,2}(\delta')=\cover\sigma\cap B^{2}$: for
$(x_{1},x_{2})\in\cover\sigma\cap B^{2}$ pair-reflexivity puts
$(x_{1},x_{2},x_{1},x_{2})$ in $\delta$ --- legitimate because
$\cover\sigma\subseteq\proj_{1,2}(\delta)$ --- and that tuple lies in $B^{4}$
and in $\cover\sigma\times\cover\sigma$. \emph{This is the step the source
asserts.} The resulting relation is linked precisely because $B$ is a block,
and is contained in $\bcol{\delta'}$ because $\cover\sigma\subseteq\bcol\delta$;
so Lemma~\ref{lem:block-good-bridge} makes $B$ BA and centre free. Then
$\cover\sigma\cap B^{2}$, being linked and subdirect on $B$, must be all of
$B^{2}$ by Theorem~\ref{lem:linked-implies-abs}. Blocks of size $1$ satisfy
$B^{2}\subseteq\cover\sigma$ trivially, so
$\cover\sigma=\operatorname{LeftLinked}(\cover\sigma)$ is an equivalence
relation and hence a congruence, giving (a); and (b) follows, one-element
algebras being vacuously BA and centre free.

For (c), fix a block $B$ of $\cover\sigma$ with $|B|\ge2$ and apply
Lemma~\ref{lem:nice-bridge-abelian} to $\delta\cap B^{4}$, whose hypotheses now
hold: $\proj_{1,2}(\delta\cap B^{4})=B^{2}$ by the same pair-reflexivity
computation, and $\bcol{\delta\cap B^{4}}=B^{2}$ because
$B^{2}\subseteq\cover\sigma\subseteq\bcol\delta$. This gives an abelian group
$G_{B}$ with $(B;\delta\cap B^{4})\cong(G_{B};x_{1}-x_{2}=x_{3}-x_{4})$.

It remains to make every $G_{B}$ elementary abelian for one common prime. Put
$\omega=\delta\cap(\cover\sigma\times\cover\sigma)$, which has the same
properties. From $x_{1}-x_{2}=x_{3}-x_{4}$ one pp-defines
$(k+1)x_{1}=k\,x_{2}+x_{3}$ for every $k$, by
\[
  \bigl((k+1)x_{1}=k\,x_{2}+x_{3}\bigr)=\exists x_{4}\;
  \bigl(k\,x_{1}=(k-1)x_{2}+x_{4}\bigr)\wedge(x_{1}-x_{2}=x_{3}-x_{4}),
\]
and hence, identifying $x_{3}$ with $x_{1}$, the relation
$q\,x_{1}=q\,x_{2}$ for every $q$. Let $S$ be what that pp-definition defines
when $x_{1}-x_{2}=x_{3}-x_{4}$ is replaced by $\omega$. If some $G_{B}$ had an
element of order $p_{1}^{2}$, or elements of coprime orders, or if two blocks
had elements of different prime orders, then for a suitable $q$ the relation $S$
would be a subuniverse of $\mathbf A^{2}$ with $S\supsetneq\sigma$ and
$\cover\sigma\not\subseteq S$, contradicting the minimality of $\cover\sigma$.
Hence every element of every $G_{B}$ has one and the same prime order $p$, and
$G_{B}\cong\mathbb Z_{p}^{n_{B}}$.
\end{proof}

\begin{hazard}[It is minimality, not irreducibility]\label{haz:minimality-not-irred}
The contradiction at the end is with the \emph{minimality of $\cover\sigma$},
which is a consequence of irreducibility; \cite{ZhukSimplified} says
``contradicts the irreducibility of $\sigma$''. The distinction matters in a
formalization, where $\cover\sigma$ is data attached to a proof of
irreducibility (Hazard~\ref{haz:irreducible}).
\end{hazard}

\begin{proof}[Proof of Lemma~\ref{lem:linear-equiv}]
$(a)\Rightarrow(b)$: take $\delta=S$ of Definition~\ref{def:linear-congruence},
which by Lemma~\ref{lem:linear-bridge} is a bridge with
$\bcol S=\cover\sigma\supsetneq\sigma$.

$(b)\Rightarrow(a)$: given a bridge $\delta$ with $\bcol\delta\supsetneq\sigma$,
which we may take reflexive, form $\delta^{\circ}$ of
Lemma~\ref{lem:symmetrise}; it is symmetric and pair-reflexive, and
$\bcol{\delta^{\circ}}\supseteq\bcol\delta\supsetneq\sigma$. Now
Lemma~\ref{lem:nontrivial-bridge} gives items (b) and (c) of
Definition~\ref{def:linear-congruence}, with the witness
$S=\delta^{\circ}\cap(\cover\sigma\times\cover\sigma)$.
\end{proof}

\begin{longlemma}\label{lem:pc-inductive-step}
Let $\sigma$ be a congruence on $\mathbf A$ and $\delta$ a reflexive bridge
from $\sigma$ to $\sigma$ with
\begin{enumerate}\romanlabels\tightlist
\item $\delta(x_{1},x_{2},x_{3},x_{4})=\delta(x_{3},x_{4},x_{1},x_{2})$;
\item $(a,b,a,b),(b,a,b,a)\in\delta$ for every $(a,b)\in\proj_{1,2}(\delta)$;
\item $\operatorname{RightLinked}(\proj_{1,2,3}(\delta))=A^{2}$.
\end{enumerate}
Then $(a,a,b,b)\in\delta$ for some $a\neq b$.
\end{longlemma}

\begin{proof}
Factor by $\sigma$ and assume $\sigma=0_{\mathbf A}$; clause (c) of
Definition~\ref{def:bridge} then forces $|A|\ge2$. Write
$P=\proj_{1,2}(\delta)$, which is symmetric by (ii), and
$W=\proj_{1,2,3}(\delta)\sdle\mathbf P\times\mathbf A$. Induct on $|A|$.

\emph{Case 1: some $B\subt{T}\mathbf A$ with $|B|>1$, $T\in\{\TBA,\TC\}$.}
\textbf{[R]} Put $W'=\proj_{1,2,3}(\delta\cap B^{4})$ and write
$W(x_{1},x_{2},x_{3})=\exists x_{4}\,\delta(x_{1},x_{2},x_{3},x_{4})\wedge
A(x_{1})\wedge A(x_{2})\wedge A(x_{3})\wedge A(x_{4})$ with $A$ the full unary
relation. Replacing $A$ by $B$ --- a \emph{single} application of
Lemma~\ref{lem:pp-propagation}, which replaces every occurrence, so no
transitivity of $\subte{\TC}$ is needed --- gives $W'\subte{T}W$. It is
nonempty, containing $(b,b,b)$ for every $b\in B$, and its projections are
$\proj_{1,2}(W')=P\cap B^{2}$, by (ii), and $\proj_{3}(W')=B$. So
Lemma~\ref{lem:absorb-linked-sub} makes $W'$ linked, which is (iii) for
$\delta\cap B^{4}$. Conditions (i) and (ii) are inherited and $\delta\cap B^{4}$
is reflexive. If it is a bridge, the inductive hypothesis applies and produces
the required pair inside $B$. If it is not, then $\proj_{1,2}(\delta\cap
B^{4})=0_{B}$, so by clause (d) every tuple of $\delta\cap B^{4}$ has the form
$(x,x,y,y)$; linkedness of $W'$ on $|B|\ge2$ points then forbids the perfect
matching, so some $(a,a,b,b)\in\delta$ with $a\neq b$.

\emph{Case 2: some $\{a\}\subt{T}\mathbf A$, $T\in\{\TBA,\TC\}$, and no such $B$
of size $>1$.} By (i) the relation $\proj_{1,3}(\delta)$ is symmetric, and by
(iii) it is linked, so $\{a\}\circ\proj_{1,3}(\delta)\neq\{a\}$; pick
$(a,b,c,d)\in\delta$ with $c\neq a$. If $a=b$ then $c=d$ by clause (d) and we
are done. If $\{a\}\circ\proj_{1,3}(\delta)\neq A$ then that set is a strong
subuniverse containing $a$ and $c$, so of size $>1$, which is Case 1.
Otherwise choose $(a,e,b,f)\in\delta$. By (ii), $(b,a,b,a)\in\delta$. Let $g$
be a ternary absorbing operation for $\{a\}$ --- available for $T=\TC$ by
Lemma~\ref{lem:central-ternary}, which the source does not cite. Applying $g$
to $(a,e,b,f)$, $(b,a,b,a)$ and $(a,a,a,a)$ gives $(a,a,g(b,b,a),a)\in\delta$,
whence $g(b,b,a)=a$ by clause (d); applying $g$ to $(b,a,b,a)$, $(b,a,b,a)$ and
$(a,e,b,f)$ then gives $(a,a,b,a)\in\delta$, whence $b=a$, a contradiction.

\emph{Case 3: $\mathbf A$ is BA and centre free.} Put
$C=\{(a,b)\in P:\{(a,b)\}\times A\subseteq W\}$. By
Lemma~\ref{lem:bacon-left}, applied to the transpose of $W$,
$C\subte{\TBA}\mathbf P$. If $C=P$ then in particular $(a,a)\in C$ for any $a$;
otherwise $C\subt{\TBA}P$ and Lemma~\ref{lem:absorbing-equality} gives
$C\cap 0_{\mathbf A}\neq\varnothing$. Either way $(a,a)\in C$ for some $a$, so
for every $c$ there is $d$ with $\delta(a,a,c,d)$, and clause (d) forces $d=c$.
Taking $c\neq a$ finishes the proof.
\end{proof}

\begin{lemma}[{\cite[Lem.~7.2]{ZhukJACM}}]\label{lem:absorbing-equality}
If $0_{\mathbf A}\subseteq\sigma\le\mathbf A^{2}$ and
$\omega\subt{\TBA}\sigma$, then $\omega\cap 0_{\mathbf A}\neq\varnothing$.
\end{lemma}

\begin{proof}[Proof of Lemma~\ref{lem:pc-bridges-trivial}]
Put $\xi=\delta\ast\delta^{-1}$, which is symmetric and pair-reflexive. Each of
$\operatorname{RightLinked}(\proj_{1,2,3}(\xi))$ and
$\operatorname{RightLinked}(\proj_{1,2,4}(\xi))$ is a congruence containing
$\sigma$, so by minimality of $\cover\sigma$ each is $\sigma$ or contains
$\cover\sigma$; this gives three cases.

If both equal $\sigma$, then for $(a,b,c,d)\in\xi$ the classes $c/\sigma$ and
$d/\sigma$ are determined by $a/\sigma,b/\sigma$; since $(a,b,a,b)\in\xi$ we get
$(a,c),(b,d)\in\sigma$, so $\xi$ is trivial. \emph{The source then asserts that
the same determinacy holds for $\delta$, ``since
$\proj_{1,2}(\delta)=\proj_{3,4}(\delta)$''. The bridge between the two is a
counting argument, and this is where that hypothesis is spent:} on
$\cover\sigma/\sigma$, which is finite, $\delta$ induces a bipartite relation
whose left neighbourhoods are nonempty, pairwise disjoint --- that is what the
case gives --- and cover the right side; equal finite cardinalities force every
neighbourhood to be a singleton, so the induced relation is a bijection and
determinacy holds in both directions. With that, the two auxiliary relations
$\zeta_{1},\zeta_{2}$ of the source are bridges when they need to be, and
$\bcol{\zeta_{1}}=\bcol{\zeta_{2}}=\sigma$ --- forced, since $\sigma$ is PC ---
yields $\proj_{1,3}(\delta)=\sigma$ or $\proj_{1,4}(\delta)=\sigma$, and by
the same argument with $x_{1},x_{2}$ interchanged, $\proj_{2,4}(\delta)=\sigma$
or $\proj_{2,3}(\delta)=\sigma$. \emph{Two of the four combinations must now be
excluded, which the source does not do:} $\proj_{1,3}=\proj_{2,3}=\sigma$ would
give $\proj_{1,2}(\delta)\subseteq\sigma$, against
$\proj_{1,2}(\delta)=\cover\sigma\supsetneq\sigma$, and likewise for
$\proj_{1,4}=\proj_{2,4}$. The two surviving combinations are the conclusion,
the reverse inclusion following from stability under $\sigma$.

If $\operatorname{RightLinked}(\proj_{1,2,3}(\xi))\supseteq\cover\sigma$, choose
a block $B$ of it that is not a $\sigma$-block. Then $a,b,d\in B$ whenever
$c\in B$ and $(a,b,c,d)\in\xi$, so $\xi\cap(A^{2}\times B\times A)\subseteq
B^{4}$ and therefore $\proj_{1,2,3}(\xi\cap B^{4})=\proj_{1,2,3}(\xi)\cap
(A^{2}\times B)$, which is linked. \emph{If $\cover\sigma\cap B^{2}\supsetneq
\sigma\cap B^{2}$ then $\xi\cap B^{4}$ is a bridge and
Lemma~\ref{lem:pc-inductive-step} applies, contradicting PC-ness via
Lemma~\ref{lem:linear-equiv}. If not, the lemma cannot be cited --- it is not a
bridge --- and the case must be inlined: clause (d) then puts both pairs of
every tuple in $\sigma$, and linkedness on a $B$ that is not a $\sigma$-block
produces $(x_{1},x_{1},x_{3},x_{3})\in\xi$ with $(x_{1},x_{3})\notin\sigma$
directly.} The third case is the mirror image, using $\proj_{1,2,4}$.
\end{proof}

\begin{hazard}[Two more hedges]\label{haz:pc-bridges-hedges}
The third case above is ``this case can be considered in the same way as
Case 2'' in \cite{ZhukSimplified}. It is genuinely the mirror: $(a,b,a,b)\in\xi$
places $b$ in the block, and $\cover\sigma\subseteq\operatorname{RightLinked}$
places $a$ and $c$. Writing it out is mechanical.
\end{hazard}

\subsection{Types interaction}\label{sec:types-interaction}

\begin{lemma}[{\cite[Lem.~8.19]{ZhukJACM}}]\label{lem:bridge-between-congruences}
Let $\omega,\sigma_{1},\sigma_{2}$ be congruences on $\mathbf A$ with
$\omega\cap\sigma_{1}=\omega\cap\sigma_{2}$ and
$\omega\setminus\sigma_{1}\neq\varnothing$. Then there is a bridge $\delta$
from $\sigma_{1}$ to $\sigma_{2}$ with $\bcol\delta=\sigma_{1}\circ\sigma_{2}$.
\end{lemma}

\begin{proof}
Put
$\delta(x_{1},x_{2},y_{1},y_{2})=\exists z_{1}\exists z_{2}\;
\sigma_{1}(x_{1},z_{1})\wedge\sigma_{1}(x_{2},z_{2})\wedge\omega(z_{1},z_{2})
\wedge\sigma_{2}(z_{1},y_{1})\wedge\sigma_{2}(z_{2},y_{2})$.
Clauses (a) and (b) are visible. For (d): if $(x_{1},x_{2})\in\sigma_{1}$ then
$(z_{1},z_{2})\in\sigma_{1}\cap\omega=\sigma_{2}\cap\omega$, so
$(y_{1},y_{2})\in\sigma_{2}$; \emph{and conversely}, if
$(y_{1},y_{2})\in\sigma_{2}$ then $(z_{1},z_{2})\in\omega\cap\sigma_{2}=
\omega\cap\sigma_{1}$, so $(x_{1},x_{2})\in\sigma_{1}$ --- the direction the
source omits, and the one that spends the other half of the hypothesis. For
(c): $\sigma_{1}\subseteq\proj_{1,2}(\delta)$ by taking $z_{i}=y_{i}=x_{1}$,
and choosing $(a,b)\in\omega\setminus\sigma_{1}$ gives $(a,b,a,b)\in\delta$
with $(a,b)\notin\sigma_{1}$ and, since
$\omega\setminus\sigma_{1}=\omega\setminus\sigma_{2}$, also
$(a,b)\notin\sigma_{2}$. Finally $\bcol\delta=\sigma_{1}\circ\sigma_{2}$: the
inclusion $\supseteq$ by taking $z_{1}=z_{2}$, using reflexivity of $\omega$,
and $\subseteq$ by reading off $(x,z_{1})\in\sigma_{1}$, $(z_{1},y)\in\sigma_{2}$.
\end{proof}

\begin{longlemma}\label{lem:two-stable}
Let $C_{1}\subt[A]{T_{1}(\sigma_{1})}B_{1}\lll A$ and
$C_{2}\subt[A]{T_{2}(\sigma_{2})}B_{2}\lll A$ with
$T_{1},T_{2}\in\{\TBA,\TC,\TS,\TD\}$, $C_{1}\cap B_{2}\neq\varnothing$,
$B_{1}\cap C_{2}\neq\varnothing$ and $C_{1}\cap C_{2}=\varnothing$. Then
$T_{1}=T_{2}\in\{\TBA,\TC,\TD\}$; and if $T_{1}=T_{2}=\TD$ there is a bridge
$\delta$ from $\sigma_{1}$ to $\sigma_{2}$ with
$\bcol\delta=\sigma_{1}\circ\sigma_{2}$.
\end{longlemma}

\begin{proof}
\emph{No $\TS$.} If $T_{1}=\TS$, pick $S_{1}\subseteq C_{1}$ with
$S_{1}\subt{\TBA,\TC}B_{1}$. Since $C_{2}\lll A$ --- the chain for $B_{2}$
extended by one $\TD$ step --- and $B_{1}\cap C_{2}\neq\varnothing$,
Lemma~\ref{lem:intersection-good}(s) gives $S_{1}\cap C_{2}\neq\varnothing$,
hence $C_{1}\cap C_{2}\neq\varnothing$. Symmetrically for $T_{2}$.

\emph{Both absorbing.} If $T_{1},T_{2}\in\{\TBA,\TC\}$ then by
Lemma~\ref{lem:ba-center-intersect} $C_{1}\cap B_{2}\subte{T_{1}}B_{1}\cap
B_{2}$ and $B_{1}\cap C_{2}\subte{T_{2}}B_{1}\cap B_{2}$; both are proper,
since $C_{1}\cap B_{2}=B_{1}\cap B_{2}$ would put $B_{1}\cap C_{2}$ inside
$C_{1}\cap C_{2}=\varnothing$. Their intersection is empty, so
Lemma~\ref{lem:possible-intersections} gives $T_{1}=T_{2}\in\{\TBA,\TC\}$.

\emph{Mixed.} Let $T_{1}\in\{\TBA,\TC\}$ and $T_{2}=\TD$ (the mirror case,
which \cite{ZhukSimplified} does not mention, is identical). By
Lemma~\ref{lem:intersection-good}(d), $(B_{1}\cap B_{2})/\sigma_{2}$ has size
$1$ or is all of $B_{2}/\sigma_{2}$. In the first case the single block is the
one defining $C_{2}$, because $B_{1}\cap C_{2}\neq\varnothing$; then
$B_{1}\cap B_{2}=B_{1}\cap C_{2}$, so $C_{1}\cap B_{2}\subseteq C_{1}\cap
C_{2}=\varnothing$, a contradiction. In the second,
Lemmas~\ref{lem:ba-center-intersect} and \ref{lem:factor-type} give
$(C_{1}\cap B_{2})/\sigma_{2}\subte{T_{1}}B_{2}/\sigma_{2}$, and this is
\emph{proper}: equality would make the $\sigma_{2}$-block defining $C_{2}$ meet
$C_{1}$, i.e.\ $C_{1}\cap C_{2}\neq\varnothing$. So $B_{2}/\sigma_{2}$ has a
proper nonempty strong subuniverse, contradicting
Definition~\ref{def:types}(iii).

\emph{Both $\TD$.} Let $\omega$ be the intersection of the dividing congruences
of the chosen chains for $B_{1}\lll A$ and $B_{2}\lll A$. Apply
Lemma~\ref{lem:intersection-good}(d) to $B_{1}$ and $C_{2}$ with the congruence
$\sigma_{1}$ --- \emph{to $C_{2}$, not to $B_{2}$; against $B_{2}$ the step does
not follow.} Its second alternative, $(B_{1}\cap C_{2})/\sigma_{1}=B_{1}/
\sigma_{1}$, would put a point of $C_{2}$ in the $\sigma_{1}$-block defining
$C_{1}$; so $|(B_{1}\cap C_{2})/\sigma_{1}|=1$ and the ``moreover'' clause
gives $\sigma_{1}\supseteq\omega\cap\sigma_{2}$, because the chain witnessing
$C_{2}\lll A$ is the chain for $B_{2}$ extended by the step
$C_{2}\subt{\TD(\sigma_{2})}B_{2}$ (Hazard~\ref{haz:lll-chain}). Symmetrically
$\sigma_{2}\supseteq\omega\cap\sigma_{1}$, so
$\omega\cap\sigma_{1}=\omega\cap\sigma_{2}$. Finally, for $c_{1}\in C_{1}\cap
B_{2}$ and $c_{2}\in B_{1}\cap C_{2}$ we have $(c_{1},c_{2})\in\omega$, since
both lie in $B_{1}\cap B_{2}$ and each $B_{i}$ sits inside a single block of
every dividing congruence of its own chain; and $(c_{1},c_{2})\notin\sigma_{1}$,
since $c_{2}\in B_{1}$ and $c_{2}\,\sigma_{1}\,c_{1}\in C_{1}$ would give
$c_{2}\in C_{1}\cap C_{2}$. Now Lemma~\ref{lem:bridge-between-congruences}
applies.
\end{proof}

\subsection{Factorisation}\label{sec:factorisation}

\begin{lemma}\label{lem:center-pushed-in}
Let $R\sdle\mathbf A_{1}\times\mathbf A_{2}$ with
$R\cap(B_{1}\times B_{2})\neq\varnothing$, $B_{i}\lll A_{i}$, $\sigma$ a
congruence on $\mathbf A_{1}$ with $B_{1}/\sigma$ BA and centre free. If some
$c\in A_{2}$ has $(E\times\{c\})\cap R\neq\varnothing$ for every $E\in
B_{1}/\sigma$, then such a $c$ may be chosen in $B_{2}$.
\end{lemma}

\begin{proof}[Proof outline]
If $c$ may already be chosen in $B_{2}$ there is nothing to do; otherwise walk
down a chain for $B_{2}\lll A_{2}$ and let $B_{2}''\subt{T(\delta)}B_{2}'$ be
the first step at which the choice fails, so $B_{2}\subseteq B_{2}''$. Put
$S'$ and $S''$ for the sets of $\sigma$-class tuples of length $|A_{1}|$
witnessed by some $c$ in $B_{2}'$, resp.\ $B_{2}''$. A single good $c\in B_{2}'$
makes $S'=(B_{1}/\sigma)^{|A_{1}|}$, and $S''\neq\varnothing$ because
$R\cap(B_{1}\times B_{2})\neq\varnothing$ and $B_{2}\subseteq B_{2}''$ ---
\emph{this is where that hypothesis is spent}. For $T\in\{\TBA,\TC\}$,
Lemmas~\ref{lem:pp-propagation} and \ref{lem:power-to-base} then put a strong
subuniverse on $B_{1}/\sigma$, a contradiction; for $T=\TS$ the same follows
once $R\cap(B_{1}\times D_{2})\neq\varnothing$, which
Lemma~\ref{lem:intersection-good}(s) supplies inside $\mathbf R$.

For $T=\TD$ one builds
$S\le(B_{1}/\sigma)^{|A_{1}|}\times B_{2}'/\delta$ and finds a full column over
$d=c/\delta$, while $(B_{1}/\sigma)^{|A_{1}|}\times B_{2}''/\delta\not\subseteq
S$ by minimality. Corollary~\ref{cor:intersection-good} then computes
$\proj_{|A_{1}|+1}(S)$. \textbf{[R]} It offers three alternatives and the source
takes the third without comment. ``Empty'' is excluded by
$R\cap(B_{1}\times B_{2})\neq\varnothing$. ``Size one'' is excluded thus: the
single $\delta$-class would have to be $[c]_{\delta}$, since $c$ has an
$R$-partner in $B_{1}$; and $B_{2}''=B_{2}'\cap E$ for a $\delta$-block $E$, so
either $E=[c]_{\delta}$, putting $c$ in $B_{2}''$ against the minimal choice of
$B_{2}'$, or $E\cap[c]_{\delta}=\varnothing$, making
$R\cap(B_{1}\times B_{2}'')=\varnothing$ against $B_{2}\subseteq B_{2}''$. So
$S$ is a central relation, and Lemmas~\ref{lem:central-relation} and
\ref{lem:power-to-base} contradict the hypotheses.
\end{proof}

\begin{lemma}\label{lem:leftlinked-stay-full}
Under the hypotheses of Lemma~\ref{lem:center-pushed-in}, if
$(\operatorname{LeftLinked}(R)\cap B_{1}^{2})/\sigma=(B_{1}/\sigma)^{2}$ then
$\operatorname{LeftLinked}\bigl(R\cap(B_{1}\times B_{2})\bigr)/\sigma=
(B_{1}/\sigma)^{2}$.
\end{lemma}

\begin{proof}[Proof outline]
Write $R_{n}=(R\circ R^{-1})^{n}$, which is reflexive and symmetric because $R$
is subdirect, so $R_{2n}=R_{n}\circ R_{n}^{-1}$ --- \emph{this is why the
source restricts to $n=2^{k}$}. If $(R_{1}\cap B_{1}^{2})/\sigma$ is already
full, Lemma~\ref{lem:bacon-left} applied to
$S=\{(a/\sigma,b):a\in B_{1},(a,b)\in R\}$, subdirect over
$(B_{1}/\sigma)\times\proj_{2}(S)$ rather than over $(B_{1}/\sigma)\times A_{2}$,
produces the required $c$, which Lemma~\ref{lem:center-pushed-in} moves into
$B_{2}$; two classes are then linked through $c$. Otherwise take the largest
$n=2^{k}$ at which fullness fails, apply Lemma~\ref{lem:bacon-left} to $R_{n}$,
push $c$ into $B_{1}$, and conclude by Lemma~\ref{lem:central-relation} that
$(R_{n}\cap B_{1}^{2})/\sigma$ is full after all.
\end{proof}

\begin{longlemma}\label{lem:cut-or-nothing}
Let $\sigma$ be a dividing congruence for $B\lll A$, let $\delta$ be a
congruence on $\mathbf A$ with $(\delta\cap B^{2})/\sigma\neq B^{2}/\sigma$, and
let $\omega$ be the intersection of the dividing congruences of the chosen chain
for $B\lll A$. Then \textup{(1)} $\delta\cap B^{2}\subseteq\sigma\cap B^{2}$;
\textup{(2)} $\sigma\supseteq\delta\cap\omega$; \textup{(3)}
$(\delta\vee(\sigma\cap\omega))\cap B^{2}=\sigma\cap B^{2}$; \textup{(4)}
$(\delta\vee(\sigma\cap\omega))\cap\omega=\sigma\cap\omega$.
\end{longlemma}

\begin{proof}
Throughout we use $B^{2}\subseteq\omega$, which holds because each $B_{i}$ in
the chain lies inside a single block of its own dividing congruence, and $B$ is
contained in every $B_{i}$. \emph{The source uses this at least four times and
never states it.}

(2). The hypothesis gives $\sigma$-classes $C_{1},C_{2}$ meeting $B$ with
$\bigl((C_{1}\cap B)\times(C_{2}\cap B)\bigr)\cap\delta=\varnothing$. Then
$C_{1}\cap B\subt{\TD(\sigma)}B$, and Lemma~\ref{lem:self-intersection-pc} with
$m=k$ and $B_{k}=C_{1}\cap B$ gives
$\bigl((C_{1}\cap B)\circ\delta\cap B\bigr)/\sigma=\{C_{1}\}$, since the value
$B/\sigma$ is excluded by the choice of $C_{2}$; its second clause is
$\sigma\supseteq\delta\cap\omega$.

(1) follows, since $\delta\cap B^{2}\subseteq\delta\cap\omega\subseteq\sigma$.

(4). Put $\delta'=\delta\vee(\sigma\cap\omega)$, which is
$\operatorname{LeftLinked}(\delta\circ(\sigma\cap\omega))$. If
$(\delta'\cap B^{2})/\sigma$ were full then Lemma~\ref{lem:leftlinked-stay-full}
would make $\operatorname{LeftLinked}$ of the restricted relation full, so some
single linking step would cross $\sigma$-classes: there would be
$a_{1},a_{2},c\in B$ and $b_{1},b_{2}\in A$ with $(a_{i},b_{i})\in\delta$,
$(b_{i},c)\in\sigma\cap\omega$ and $(a_{1},a_{2})\notin\sigma$. But
$a_{1},a_{2},c\in B$ gives $(a_{i},c)\in\omega$, hence $(a_{i},b_{i})\in\omega$,
hence $(a_{i},b_{i})\in\sigma$ by (2), hence $(a_{1},a_{2})\in\sigma$ --- a
contradiction. So $(\delta'\cap B^{2})/\sigma\neq(B/\sigma)^{2}$, and applying
(2) to $\delta'$ gives $\sigma\supseteq\delta'\cap\omega$; with
$\delta'\supseteq\sigma\cap\omega$ this is (4).

(3) follows from (4) by intersecting with $B^{2}\subseteq\omega$.
\end{proof}

\begin{lemma}\label{lem:main-existence}
\textbf{[R]} Let $B\lll A$ and let $\sigma$ be a congruence on $\mathbf A$ with
$|B/\sigma|>1$ and $B/\sigma$ BA and centre free. Then there is a dividing
congruence $\delta\supseteq\sigma$ for $B\le A$.
\end{lemma}

\begin{proof}
Choose $\delta\supseteq\sigma$ maximal with $|B/\delta|>1$; it exists because
$\sigma$ qualifies. Then $B/\delta$ is BA and centre free, since the preimage
under $B/\sigma\twoheadrightarrow B/\delta$ of a proper nonempty strong
subuniverse would be one of $B/\sigma$.

It remains to see that $\delta$ is irreducible. Suppose
$\delta=S_{1}\cap\dots\cap S_{k}$ with each $S_{i}\supsetneq\delta$ a
subuniverse of $\mathbf A^{2}$ \emph{stable under $\delta$}. Some $S_{i}$ has
$B^{2}\not\subseteq S_{i}$, else $B^{2}\subseteq\delta$. Now
$\operatorname{LeftLinked}(S_{i})$ is a congruence containing
$S_{i}\supsetneq\delta\supseteq\sigma$, so maximality of $\delta$ forces
$B^{2}\subseteq\operatorname{LeftLinked}(S_{i})$, i.e.\
$(\operatorname{LeftLinked}(S_{i})\cap B^{2})/\delta=(B/\delta)^{2}$. By
Lemma~\ref{lem:leftlinked-stay-full}, applied \emph{with the congruence
$\delta$}, $\operatorname{LeftLinked}(S_{i}\cap B^{2})/\delta=(B/\delta)^{2}$,
so $(S_{i}\cap B^{2})/\delta$ is linked. It is also \emph{proper}: stability
under $\delta$ makes $S_{i}$ a union of products of $\delta$-blocks, so
$(S_{i}\cap B^{2})/\delta=(B/\delta)^{2}$ would give $B^{2}\subseteq S_{i}$.
Theorem~\ref{lem:linked-implies-abs} then produces a proper nonempty strong
subuniverse of $B/\delta$, a contradiction.
\end{proof}

\begin{hazard}[Two slips, one cause]\label{haz:main-existence}
\cite{ZhukSimplified} factors by $\sigma$ in the last four lines and states the
conclusion for $B/\delta$; those do not match, and with $\sigma$ the properness
step has no support. Reading $\delta$ throughout is plainly what was meant ---
$B/\delta$ was shown BA and centre free two lines earlier for exactly this use.
Properness then needs the stability clause, which the source drops when it
introduces the $S_{i}$; stability is the only part of
Definition~\ref{def:irreducible} doing work here.
\end{hazard}

\begin{lemma}\label{lem:factor-by-delta}
Let $\delta$ be a congruence on $\mathbf A$ and
$C\subt[A]{T(\sigma)}B\lll A$ with $T\in\{\TPC,\TL\}$. Then
$C/\delta=B/\delta$, or $C/\delta\subt{\TS}B/\delta$, or
$C/\delta\subt[A/\delta]{T}B/\delta$; and in the last case
$\delta\cap B^{2}\subseteq\sigma\cap B^{2}$.
\end{lemma}

\begin{proof}[Proof outline]
Put $R=\{(a/\sigma,a/\delta):a\in B\}$ and note
$R\circ R^{-1}=(\delta\cap B^{2})/\sigma$, so
$\operatorname{LeftLinked}(R)=(B/\sigma)^{2}$ exactly when
$(\delta\cap B^{2})/\sigma=B^{2}/\sigma$. In that case
Lemma~\ref{lem:bacon-left} yields $U\subte{\TBA,\TC}B/\delta$ with
$(B/\sigma)\times U\subseteq R$, and $U\subseteq C/\delta$, giving the first two
alternatives.

Otherwise Lemma~\ref{lem:cut-or-nothing}(4) gives
$\sigma\cap\omega=\delta'\cap\omega$ for $\delta'=\delta\vee(\sigma\cap\omega)$,
and $B/\delta'\cong B/\sigma$ by (3), so $B/\delta'$ is BA and centre free and
Lemma~\ref{lem:main-existence} produces a dividing congruence
$\delta''\supseteq\delta'$ for $B$. Applying
Lemma~\ref{lem:cut-or-nothing}(2) to $\delta''$ --- legitimate because
$\sigma\cap B^{2}\subseteq\delta''$ and $|B/\delta''|>1$ give
$(\delta''\cap B^{2})/\sigma\neq B^{2}/\sigma$ --- yields
$\delta''\cap\omega=\sigma\cap\omega$, hence $\delta''\cap B^{2}=\sigma\cap
B^{2}$, so the $\delta''$-classes and $\sigma$-classes agree on $B$ and
$C/\delta\subt[A/\delta]{\mathcal T(\delta''/\delta)}B/\delta$. Finally
$\mathcal T=T$: $\omega\setminus\delta''\neq\varnothing$, since
$\omega\supseteq B^{2}$ and $|B/\sigma|>1$, so
Lemma~\ref{lem:bridge-between-congruences} gives a bridge from $\delta''$ to
$\sigma$, and Lemma~\ref{lem:no-cross-bridge} in both directions makes the two
congruences of the same type. The last clause holds because
$\delta\cap B^{2}\subseteq\delta'\cap B^{2}=\sigma\cap B^{2}$.
\end{proof}

\subsection{The headline properties}\label{sec:headline}

\begin{proof}[Proof of Lemma~\ref{lem:ubiquity}]
If $\mathbf B$ has a proper nonempty binary absorbing or central subuniverse,
take it. Otherwise $\mathbf B$ is BA and centre free and $|B|>1$, so
Lemma~\ref{lem:main-existence} with $\sigma=0_{\mathbf A}$ gives a dividing
congruence $\delta$ for $B$; take $C=B\cap E$ for any block $E$ \emph{of
$\delta$} with $B\cap E\neq B$, which exists because $|B/\delta|>1$. Its type is
$\TL$ or $\TPC$ according to $\delta$.
\end{proof}

\begin{proof}[Proof of Lemma~\ref{lem:intersect-all}]
(t). For $T\in\{\TBA,\TC\}$ this is Lemma~\ref{lem:ba-center-intersect}. For
$T=\TS$, take $E\subseteq C$ with $E\subt{\TBA,\TC}B$; if $C\cap D=\varnothing$
there is nothing to prove, and otherwise $B\cap D\neq\varnothing$, so
Lemma~\ref{lem:intersection-good}(s) gives $E\cap D\neq\varnothing$ and
Lemma~\ref{lem:ba-center-intersect} gives $E\cap D\subt{\TBA,\TC}B\cap D$. For
$T\in\{\TPC,\TL\}$, Lemma~\ref{lem:intersection-good}(d) makes $(B\cap D)/\sigma$
empty, a point, or all of $B/\sigma$; the three cases give $C\cap D=\varnothing$,
$C\cap D\in\{\varnothing,B\cap D\}$, and $C\cap D\subt{T(\sigma)}B\cap D$
respectively.

(i). Apply (t) to each step of a chain for $B\lll^{A}A$, obtaining
$B\cap D\dlll^{A}D$; with $D\lll A$ this gives $B\cap D\dlll A$.
\end{proof}

\begin{hazard}[The statement is weaker than the proof]\label{haz:intersect-all-strength}
\textbf{[R]} The proof of (i) establishes $B\cap D\dlll^{A}D$ and then weakens
it. Three later arguments use the stronger form while citing (i):
Corollary~\ref{cor:prop-relations}(b) and (m), and
Lemma~\ref{lem:multitype-chain}. State (i) as $B\cap D\dlll^{A}D$ and derive
$B\cap D\dlll A$ as a corollary.
\end{hazard}

\begin{proof}[Proof of Lemma~\ref{lem:propagation}]
(bt) is the reverse-homomorphism lemma and (b) follows by iteration. (fs) is
Lemma~\ref{lem:factor-type}. (ft) is Lemma~\ref{lem:factor-by-delta} for
$T\in\{\TPC,\TL,\TD\}$ and follows from (fs) for $T\in\{\TBA,\TC,\TS\}$ ---
\emph{the source cites only the former, which does not cover the absorption
types}; (f) then follows by iteration, all three outcomes of (ft) staying inside
the type list of Definition~\ref{def:lll}.

(fm). Let $C=C_{1}\cap\dots\cap C_{t}$ with $C_{i}\subt[A]{T(\sigma_{i})}B$ and
induct on $t$. For $t=1$ this is (ft), and \emph{the hypothesis that $f(B)$ is
$\TS$-free is exactly what kills its middle alternative}, which $\MT T$ cannot
absorb. If $C_{i}/\delta=B/\delta$ for some $i$, say $i=t$, put
$D_{j}=C_{j}\cap C_{t}$; by Lemma~\ref{lem:intersect-all}(t)
$D_{j}\subte[A]{T(\sigma_{j})}C_{t}$, and the inductive hypothesis applied to
$C\subte{\MT T}C_{t}$ gives the claim. Otherwise every
$C_{i}/\delta\subt[A/\delta]{T}B/\delta$, and
$\delta\cap B^{2}\subseteq\sigma_{i}\cap B^{2}$ for every $i$ by the last clause
of Lemma~\ref{lem:factor-by-delta}. That is what makes the images meet
correctly: $(C_{i}\circ\delta)\cap B=C_{i}$, so a $\delta$-class $E$ lying in
every $C_{i}/\delta$ has $E\cap B\subseteq C_{1}\cap\dots\cap C_{t}$, whence
\[
  (C_{1}\cap\dots\cap C_{t})/\delta=C_{1}/\delta\cap\dots\cap C_{t}/\delta ,
\]
the inclusion $\subseteq$ being trivial. This is the one place where the failure
of images to commute with intersections (\S\ref{sec:legislation}, item
\ref{leg:quotient}) is repaired, and it should be extracted as a lemma in its
own right. (bm) follows from (bt) by intersecting.
\end{proof}

\begin{longtheorem}[Stable intersection]\label{thm:stable-intersection}
Let $C_{i}\subt[A]{T_{i}(\sigma_{i})}B_{i}\lll A$ with
$T_{i}\in\{\TBA,\TC,\TS,\TL,\TPC\}$ for $i\in[n]$, $n\ge2$, and suppose
$\bigcap_{i}C_{i}=\varnothing$ while $B_{j}\cap\bigcap_{i\neq j}C_{i}\neq
\varnothing$ for every $j$. Then all $T_{i}$ are $\TBA$; or all are $\TL$ and
for all $k,\ell$ there is a bridge $\delta$ from $\sigma_{k}$ to $\sigma_{\ell}$
with $\bcol\delta=\sigma_{k}\circ\sigma_{\ell}$; or $n=2$ and
$T_{1}=T_{2}=\TC$; or $n=2$, $T_{1}=T_{2}=\TPC$ and $\sigma_{1}=\sigma_{2}$.
\end{longtheorem}

\begin{proof}
For $n=2$ combine Lemmas~\ref{lem:two-stable} and \ref{lem:no-cross-bridge}.
For general $n$ put $B_{2}'=B_{2}\cap C_{3}\cap\dots\cap C_{n}$ and
$C_{2}'=C_{2}\cap B_{2}'$; then $C_{2}'\subt{T_{2}(\sigma_{2})}B_{2}'\lll A$ by
Lemma~\ref{lem:intersect-all}, properly because $C_{2}'=B_{2}'$ would make
$C_{1}\cap B_{2}'$ empty. The three hypotheses of the pair case hold, so
$T_{1}=T_{2}$ and, for the dividing types, the bridge exists.

If $T_{1}=\TPC$ then $\sigma_{1}\circ\sigma_{2}=\sigma_{1}=\sigma_{2}$, since
otherwise composing the bridge with its transpose gives a bridge from
$\sigma_{1}$ to itself with $\bcol{}\supsetneq\sigma_{1}$, making $\sigma_{1}$
linear. With $\sigma_{1}=\sigma_{2}=:\sigma$, the blocks $E_{1},E_{2}$ defining
$C_{1},C_{2}$ are \emph{distinct} --- equal blocks would put the point of
$B_{1}\cap C_{2}'$ in $C_{1}\cap C_{2}'$ --- so $C_{1}\cap C_{2}=\varnothing$;
applying this to every pair and then the hypothesis at $j=3$ forces $n=2$.

\textbf{[R]} For $T_{1}=\TC$ the source offers nothing, and the argument above
has no analogue: there is no congruence and no pair of blocks. Here is the
missing proof. Suppose $n\ge3$, put $D=\bigcap_{i}B_{i}$ and $E_{i}=C_{i}\cap D$.
Then $D$ is a nonempty subuniverse, containing $B_{1}\cap C_{2}\cap\dots\cap
C_{n}$; each $E_{i}\subte{\TC}D$ by Lemma~\ref{lem:ba-center-intersect}, and
properly, since $E_{j}=D$ would give $\bigcap_{i}E_{i}=\bigcap_{i\neq j}E_{i}
\neq\varnothing$. Moreover $\bigcap_{i}E_{i}=\varnothing$ and, because
$C_{i}\subseteq B_{i}$,
\[
  \textstyle\bigcap_{i\neq j}E_{i}=D\cap\bigcap_{i\neq j}C_{i}
  =B_{j}\cap\bigcap_{i\neq j}C_{i}\neq\varnothing .
\]
Now the diagonal $R=\{(a,\dots,a):a\in D\}\le D^{n}$ is a subuniverse by
idempotence, and the two displays say exactly that $R$ is
$(E_{1},\dots,E_{n})$-essential. By \cite[Lem.~6.11]{ZhukStrong} no such
relation exists for $n\ge3$ when the $E_{i}$ are central. Hence $n=2$.
\end{proof}

\begin{hazard}[What the repair costs]\label{haz:d4}
\cite{ZhukSimplified} argues $n=2$ only for $\TPC$, then asserts case (c). The
repair needs \cite[Lem.~6.11]{ZhukStrong}, whose printed proof cites
\emph{itself} three times where it means the preceding lemma; a self-contained
rendering must supply both. Every ingredient is already formalised in the
predecessor Lean development: essentiality, the Barto--Kazda criterion, and
``central implies ternary absorbing''.
\end{hazard}

\begin{proof}[Proof of Corollary~\ref{cor:stable-intersection}]
Pull back along the projections $f_{i}:\mathbf R\to\mathbf A_{i}$ using
Lemma~\ref{lem:propagation}(b),(bt), apply
Theorem~\ref{thm:stable-intersection} inside $\mathbf R$, and translate
$\bcol\delta=\sigma_{k}'\circ\sigma_{\ell}'$ back as
$\sigma_{k}\circ\proj_{k,\ell}(R)\circ\sigma_{\ell}$.
\end{proof}

\begin{proof}[Proof of Lemma~\ref{lem:multitype-chain}]
With $C=\bigcap_{i\le n}C_{i}$ and $D_{j}=\bigcap_{i\le j}C_{i}$,
Lemma~\ref{lem:intersect-all}(t) gives $D_{j+1}\subte[A]{T}D_{j}$, using
$D_{j}\lll A$ from part (i); collapsing the equalities gives the chain.
\end{proof}

\begin{proof}[Proof of Lemma~\ref{lem:linear-on-top}]
\textbf{[R]} By Lemma~\ref{lem:linear-equiv} there is a bridge $\delta$ from
$\sigma$ to $\sigma$ with $\bcol\delta\supsetneq\sigma$, which we may take
reflexive; replace it by $\delta^{\circ}$, symmetric and pair-reflexive.
Lemma~\ref{lem:nontrivial-bridge}(c) applies, and $\cover\sigma=A^{2}$ is a
congruence with the single block $A$, so $A/\sigma\cong\mathbb Z_{p}^{m}$ as a
group. Now $m=0$ gives $\sigma=A^{2}$, impossible for an irreducible congruence
(\S\ref{sec:legislation}, item \ref{leg:emptyfamily}). And $m\ge2$ contradicts
irreducibility: $A/\sigma$ is affine, so the basic operation acts as
$\sum_{i}c_{i}x_{i}$ with $\sum_{i}c_{i}=1$, and the coordinate congruences
$\theta_{i}=\{(x,y):x_{i}=y_{i}\}$ are linear subspaces of
$(\mathbb Z_{p}^{m})^{2}$, hence subuniverses of $(\mathbf A/\sigma)^{2}$, each
properly containing $0_{\mathbf A/\sigma}$, with
$\bigcap_{i}\theta_{i}=0_{\mathbf A/\sigma}$ --- so
$0_{\mathbf A/\sigma}$ is reducible, and by Lemma~\ref{lem:irred-quotient} so is
$\sigma$. Hence $m=1$, and the group isomorphism is an isomorphism of algebras
by Proposition~\ref{prop:p-divides}.
\end{proof}

\begin{lemma}\label{lem:multiply-all-linear}
Let $C_{1}\subt[A]{\MT T}B_{1}\lll A$ with $T\in\{\TPC,\TL,\TD\}$, let
$B_{2}\lll A$, $C_{1}\cap B_{2}=\varnothing$ and $B_{1}\cap B_{2}\neq
\varnothing$. Then $\bigl(C_{1}\circ(\omega_{1}\cap\dots\cap\omega_{s})\bigr)
\cap B_{2}=\varnothing$, where $\omega_{1},\dots,\omega_{s}$ are all congruences
of type $T$ on $\mathbf A$ with $\cover{\omega_{i}}\supseteq B_{1}^{2}$.
\end{lemma}

\begin{proof}
Downward induction on $|B_{1}|$, starting at $B_{1}=A$. Write
$C_{1}=\bigcap_{i}C_{1}^{i}$ with $C_{1}^{i}\subt[A]{\TD(\sigma_{i})}B_{1}$, so
each $\sigma_{i}$ is among the $\omega_{j}$ and
$C_{1}^{i}=(C_{1}^{i}\circ\sigma_{i})\cap B_{1}$; by
Corollary~\ref{cor:prop-times-cong}(e),(f),
$C_{1}^{i}\circ\sigma_{i}\subt[A]{\TD(\sigma_{i})}B_{1}\circ\sigma_{i}\lll A$.
Since $\bigcap_{i}(C_{1}^{i}\circ\sigma_{i})\cap B_{1}\cap B_{2}=C_{1}\cap
B_{2}=\varnothing$, walk down a chain from $A$ to $B_{1}$: either the
intersection with $B_{2}$ is already empty at $A$, and then
$C_{1}\circ\bigcap_{j}\omega_{j}\subseteq\bigcap_{i}(C_{1}^{i}\circ\sigma_{i})$
finishes; or there is a first step $B_{1}'\subt{\TD}B_{1}''$ at which it becomes
empty, and the inductive hypothesis at $B_{1}''$, together with
$C_{1}\subseteq B_{1}'$ and $\bigcap_{j}\omega_{j}\subseteq\sigma_{i}$, finishes.
\end{proof}

\begin{hazard}[A citation that does no work]\label{haz:multiply-citation}
\cite{ZhukSimplified} invokes Theorem~\ref{thm:stable-intersection} for the
dichotomy above. No theorem is consumed: the intersection is empty at $B_{1}$
by hypothesis, so either it is empty at $A$ or there is a first step where it
becomes so.
\end{hazard}

\begin{proof}[Proof of Lemma~\ref{lem:maximal-mult}]
\textbf{[R]} Write $\bar X$ for $X/\sigma$ and $C_{1}=\bigcap_{j\le t}C_{1}^{j}$
as above. By Lemma~\ref{lem:propagation}(ft) each $\bar C_{1}^{j}$ is
$\bar B_{1}$, or $\subt{\TS}\bar B_{1}$, or $\subt{T}\bar B_{1}$. Choose
$F\subseteq\{\bar C_{1}^{1},\dots,\bar C_{1}^{t}\}$ \emph{minimal} with
$\bigcap F\cap\bar B_{2}=\varnothing$; it exists because the whole family works,
it is nonempty because $\bigcap\varnothing=\bar B_{1}$ meets $\bar B_{2}$, and no
member of it equals $\bar B_{1}$.

Walk down a chain for $\bar B_{2}\lll\bar A$ to the first step
$\bar B_{2}'\subt{T'}\bar B_{2}''$ with $\bigcap F\cap\bar B_{2}'=\varnothing$
and $\bigcap F\cap\bar B_{2}''\neq\varnothing$. Apply
Theorem~\ref{thm:stable-intersection} to the family $F\cup\{\bar B_{2}'\}$, with
ambient $\bar B_{1}$ for the members of $F$ and $\bar B_{2}''$ for
$\bar B_{2}'$. Its leave-one-out hypothesis holds in both shapes: omitting a
member of $F$ is minimality of $F$ together with $\bar B_{2}\subseteq
\bar B_{2}'$, and omitting $\bar B_{2}'$ is the choice of the step. So all the
types coincide, and the theorem's four conclusions admit only $\TBA$, $\TC$,
$\TL$ and $\TPC$: \emph{no member of $F$ has type $\TS$}.

Hence every member of $F$ is $\subt{T}\bar B_{1}$, so
$\bigcap F\subt{\MT T}\bar B_{1}$ --- properly, since $\bigcap F\cap\bar B_{2}
=\varnothing$ while $\bar B_{1}\cap\bar B_{2}\neq\varnothing$ --- with the
ambient still $\bar B_{1}$. Lemma~\ref{lem:multiply-all-linear} inside
$\mathbf A/\sigma$ now returns congruences $\delta_{i}$ of type $T$ with
$\cover{\delta_{i}}\supseteq\bar B_{1}^{2}$ and
$\bigl(\bigcap F\circ\bigcap_{i}\delta_{i}\bigr)\cap\bar B_{2}=\varnothing$.
Their preimages $\omega_{i}\supseteq\sigma$ satisfy
$\cover{\omega_{i}}\supseteq B_{1}^{2}$ and, since $C_{1}$ is contained in the
preimage of $\bigcap F$, also
$\bigl(C_{1}\circ\bigcap_{i}\omega_{i}\bigr)\cap B_{2}=\varnothing$. Maximality
of $\sigma$ gives $\sigma=\bigcap_{i}\omega_{i}$.
\end{proof}

\begin{hazard}[What the source does instead]\label{haz:maximal-mult}
\cite{ZhukSimplified} splits on whether $\bar B_{1}$ is $\TS$-free and, in the
negative case, argues that a BA-and-central $E\subt{\TS}\bar B_{1}$ meets both
$\bar C_{1}$ and $\bar B_{2}$, ``hence'' those two meet. They need not: two sets
meeting a third need not meet each other, and the configuration reproduces
itself on $E$, so no iteration of those facts can close it. Passing to a
minimal subfamily avoids the split altogether --- minimality is precisely the
leave-one-out hypothesis the stable-intersection theorem wants, and its
conclusion has no room for $\TS$.
\end{hazard}
\clearpage
%% ---------------------------------------------------------------- ch6-main.tex
\section{The main statements}\label{sec:main-statements}

This section is Zhuk's \S\S3.3--3.4. Everything converges here.
Theorem~\ref{thm:main-inductive} is the single induction that carries the whole
proof; the three theorems the algorithm consumes are corollaries of it.

\subsection{The three claims the algorithm needs}

The algorithm of \S\ref{sec:algorithm} is justified by exactly three facts.
Informally:

\begin{description}\tightlist
\item[\textup{(IC1)}]\label{ic:existence} every domain of size $\ge 2$ has a
  strong subset, or an equivalence $\sigma$ with
  $D_{x}/\sigma\cong\mathbb Z_{p}$;
\item[\textup{(IC2)}]\label{ic:safe} reducing a variable's domain to a strong
  subset of it does not destroy solvability, for a consistent enough instance;
\item[\textup{(IC3)}]\label{ic:codim} for a linked, consistent enough,
  strong-subset-free instance, the set of ``linear cosets'' containing a
  solution is empty, full, or a codimension-one affine subspace.
\end{description}

(IC1) is Lemma~\ref{lem:ubiquity} together with
Lemma~\ref{lem:linear-on-top}. (IC2) is
Theorem~\ref{thm:reductions-safe}. (IC3) is
Theorem~\ref{thm:codim-one}.

\subsection{The main induction}

\begin{longtheorem}[Main inductive claim]\label{thm:main-inductive}
Let $\mathcal I$ be an irreducible cycle-consistent instance and let $D^{(1)}$
be a $1$-consistent reduction for $\mathcal I$ with $D^{(1)}\lll D$.
\begin{enumerate}\tightlist
\item[\textup{(1)}] If $\mathcal I$ is crucial in $D^{(1)}$ then
  \textup{(1a)} holds, and \textup{(1b)} or \textup{(1c)} holds:
  \begin{enumerate}\tightlist
  \item[\textup{(1a)}] every constraint of $\mathcal I$ has the parallelogram
    property;
  \item[\textup{(1b)}] $\mathcal I$ is a connected linear-type instance whose
    solution set is subdirect;
  \item[\textup{(1c)}] there is an expanded covering $\mathcal J$ of
    $\mathcal I$, crucial in $D^{(1)}$, with a linked connected subinstance
    $\Upsilon$ whose solution set is not subdirect.
  \end{enumerate}
\item[\textup{(2)}] If $D^{(2)}\subte{T}D^{(1)}$ is a $1$-consistent reduction
  for $\mathcal I$ with $T\in\{\TBA,\TC\}$, and $\mathcal I^{(1)}$ has a
  solution, then $\mathcal I^{(2)}$ has a solution.
\end{enumerate}
\end{longtheorem}

\begin{hazard}[The induction, stated precisely]\label{haz:the-induction}
The proof reads ``we prove the claim by induction on the size of $D^{(1)}$'',
and this one line hides the entire well-foundedness argument. Making it precise
is the first thing a formalization must do, because \emph{the two halves of the
conclusion are proved by mutual induction and each invokes the other at a
smaller reduction}.

The measure is
\[
  \mu(D^{(1)})=\sum_{x\in\Var(\mathcal I)}\bigl|D^{(1)}_{x}\bigr|\in\mathbb N,
\]
and the statement being proved by strong induction on $\mu$ is the
conjunction of (1) and (2), \emph{universally quantified over $\mathcal I$}:
\[
  P(k):\quad
  \forall\,\mathcal I,\ \forall\,D^{(1)}\ \text{with}\ \mu(D^{(1)})\le k,\
  \textup{(1)}\ \text{and}\ \textup{(2)}\ \text{hold}.
\]
The quantifier over $\mathcal I$ must be \emph{inside}, because every use of the
inductive hypothesis changes the instance: (2) applies it to a weakening
$\mathcal I'$ of $\mathcal I^{(2)}$; Case~1 of (1) applies it to $\mathcal I$
with $D^{(2)}$; Case~2 applies it to weakenings and to expanded coverings. An
induction that fixed $\mathcal I$ and inducted on $D^{(1)}$ alone does not go
through.

Two further points that the one-line phrasing conceals.
\begin{enumerate}\alphlabels\tightlist
\item \emph{The instance grows.} Expanded coverings have \emph{more} variables
  than $\mathcal I$. So $\mu$ is not monotone in the instance, and the
  induction is genuinely on $\mu(D^{(1)})$ alone --- with the reduction of a
  covering being the extension of $D^{(1)}$ along the parent map
  (Proposition~\ref{prop:covering-basics}(h)), whose $\mu$ is \emph{larger}.
  Every appeal to the inductive hypothesis at a covering must therefore be
  checked to occur at a strictly smaller reduction, not merely at a different
  instance. In Case~2 of (1) the reductions used are the $D^{(B,\bot)}$, which
  are contained in $D^{(1)}$ restricted at $z$ to a proper subset $B$; that is
  where the decrease comes from.
\item \emph{Which appeals decrease.} (2) invokes the hypothesis at
  $D^{(2)}\subsetneq D^{(1)}$: strict because $D^{(2)}\subt{T}D^{(1)}$ has
  $C\lneq B$ at some variable. Case~1 of (1) likewise. Case~2 invokes it at
  $D^{(B,\bot)}\subseteq D^{(B,\top)}$, which is $D^{(1)}$ with the $z$-domain
  cut to $B\lneq D^{(1)}_{z}$. In every case the decrease is at a single
  variable and by at least one element. A formalization should extract this as
  a lemma: \emph{if $D'\subte{T}D$ with $T$ any type and $D'\neq D$, then
  $\mu(D')<\mu(D)$}, and use it uniformly.
\end{enumerate}
\end{hazard}

\begin{proof}[Proof of \textup{(2)}]
Suppose $\mathcal I^{(2)}$ has no solution. Weaken $\mathcal I^{(2)}$ to an
instance $\mathcal I'$ crucial in $D^{(2)}$ (Remark~\ref{rem:get-crucial}). By
the inductive hypothesis applied to $\mathcal I'$ and $D^{(2)}$, $\mathcal I'$
satisfies (1a) and (1b) or (1c).

\emph{Case \textup{(1c)}.} There is an expanded covering $\mathcal J$ of
$\mathcal I'$, crucial in $D^{(2)}$, with a linked connected subinstance
$\Upsilon$. Let $x$ be a variable of a constraint $C\in\Upsilon$. By
Lemma~\ref{lem:connected}(p), $\ConOne(C,x)$ is a perfect linear congruence;
choose $\zeta\subseteq\mathbf D_{x}\times\mathbf D_{x}\times\mathbf Z_{p}$ with
$(y_{1},y_{2},0)\in\zeta\iff(y_{1},y_{2})\in\ConOne(C,x)$ and
$\proj_{1,2}(\zeta)=\cover{\ConOne(C,x)}$.

Replace the variable $x$ of $C$ in $\mathcal J$ by a fresh $x'$ and add the
constraint $\zeta(x,x',z)$ with $z$ fresh of domain $\mathbf Z_{p}$; call the
result $\Theta$, and extend $D^{(1)},D^{(2)}$ by $D^{(1)}_{x'}=D^{(2)}_{x'}
=D_{x'}$. Let $E_{i}$ be the set of $b\in\mathbf Z_{p}$ such that $\Theta$ has
a solution in $D^{(i)}$ with $z=b$. By
Lemma~\ref{lem:ba-center-implies}, $E_{2}\dsubte{T}E_{1}$.

Since $\mathcal J$ is crucial in $D^{(2)}$, $\Theta^{(2)}$ has a solution but
none with $z=0$; since $\mathcal I^{(1)}$ has a solution, so does
$\mathcal J^{(1)}$, hence $\Theta^{(1)}$ has one with $z=0$. Thus
$0\in E_{1}$, $0\notin E_{2}$, and $E_{1}$ has at least two elements. As
$\mathbf Z_{p}$ has no proper subalgebra of size $>1$, $E_{1}=\mathbf Z_{p}$;
then $E_{2}\subt{T}\mathbf Z_{p}$ with $T\in\{\TBA,\TC\}$, contradicting
Lemma~\ref{lem:no-abs-in-linear}.

\emph{Case \textup{(1b)}.} Corollary~\ref{cor:same-type} gives a
contradiction directly.
\end{proof}

\begin{hazard}[Three gaps in the case \textup{(1c)} argument]\label{haz:case1c}
\begin{enumerate}\alphlabels\tightlist
\item ``$E_{1}$ contains at least two different elements, $0\in E_{1}$ and
  $0\notin E_{2}$, hence $E_{1}=\mathbf Z_{p}$'' uses that $E_{1}$ is a
  subuniverse of $\mathbf Z_{p}$ --- true because $E_{1}=\Theta^{(1)}(z)$ is a
  relation defined by an instance (Definition~\ref{def:instance-relation}) ---
  \emph{and} that the subuniverses of $\mathbf Z_{p}$ are the singletons and
  the whole set. The latter needs $w^{\mathbf Z_{p}}$ to generate: from $a\ne b$
  one reaches every element. Prove it; it is one line but it is the crux.
\item $E_{2}\dsubte{T}E_{1}$ is applied with the \emph{dotted} conclusion of
  Lemma~\ref{lem:ba-center-implies}, and the argument then treats $E_{2}$ as a
  proper subuniverse. The case $E_{2}=\varnothing$ must be handled: it is
  excluded because $\Theta^{(2)}$ has a solution. Say so.
\item The choice of $\zeta$ is by Definition~\ref{def:perfect-linear}, which
  supplies $\zeta$ with $\proj_{1,2}\zeta=\cover\sigma$ and
  $(a_{1},a_{2},b)\in\zeta\Rightarrow((a_{1},a_{2})\in\sigma\iff b=0)$. The
  proof uses the \emph{biconditional} $(y_{1},y_{2},0)\in\zeta\iff(y_{1},y_{2})
  \in\ConOne(C,x)$, which is stronger: it needs, for each $(y_1,y_2)\in\sigma$,
  that $(y_{1},y_{2},0)\in\zeta$. This follows from $\proj_{1,2}\zeta
  =\cover\sigma\supseteq\sigma$ plus the implication, but the step is
  suppressed.
\end{enumerate}
\end{hazard}

\begin{proof}[Proof of \textup{(1)}, sketch of the structure]
First, $|D^{(1)}_{x}|>1$ for some $x$: otherwise $\mathcal I^{(1)}$ has all
domains singletons and $1$-consistency makes it solvable, contradicting
cruciality.

\emph{Case 1: some $D^{(1)}_{x}$ has a nontrivial $\TBA$ or central
subuniverse.} Lemma~\ref{lem:find-consistent} yields a $1$-consistent reduction
$D^{(2)}\subte{T}D^{(1)}$, $T\in\{\TBA,\TC\}$. One checks $\mathcal I$ is
crucial in $D^{(2)}$: for a weakening $\mathcal J$ of a constraint,
$\mathcal J^{(1)}$ has a solution, so by the inductive hypothesis (part (2) at
$D^{(2)}$) $\mathcal J^{(2)}$ has one. Applying the inductive hypothesis
(part (1)) at $D^{(2)}$ gives the conclusion.

\emph{Case 2: otherwise.} Lemma~\ref{lem:ubiquity} gives
$E\subt[D_{z}]{T(\sigma)}D^{(1)}_{z}$ for some $z$ and $T\in\{\TPC,\TL\}$.
Property (1a) is obtained by choosing, for each $x$, a minimal
$D^{(2)}_{x}\subte{\MT T}D^{(1)}_{x}$ containing the value at $x$ of a solution
produced by weakening a fixed constraint $C$; Lemma~\ref{lem:multitype-chain}
gives $D^{(2)}\lll^{D}D^{(1)}$ and
Lemma~\ref{lem:minimal-consistent} splits into two subcases, in one of which the
inductive hypothesis applies and in the other of which
Lemma~\ref{lem:parallelogram-crucial} applies directly.

For (1b)/(1c) one forms, for each $B$ in the family
$\mathcal B=\{B:B\subt[D_{z}]{T(\sigma)}D^{(1)}_{z}\}$, the reduction
$D^{(B,\top)}$ that equals $D^{(1)}$ off $z$ and $B$ at $z$, and the maximal
$1$-consistent reduction $D^{(B,\bot)}\subseteq D^{(B,\top)}$ (possibly empty).
Lemma~\ref{lem:tree-coverings} supplies tree-coverings $\Upsilon_{B,x}$
computing $D^{(B,\bot)}_{x}$, and their conjunction $\Upsilon_{x}$ works
uniformly in $B$. Writing $\mathcal B_{0}$ for the $B$ with $D^{(B,\bot)}$
nonempty, the argument splits:
\begin{itemize}\tightlist
\item $\mathcal B_{0}=\varnothing$: the type must be $\TL$, and
  Lemma~\ref{lem:bridge-from-instance} applied to a suitably weakened
  tree-covering produces, for any two constraints sharing a variable, a
  reflexive bridge between their $\ConOne$'s. Hence $\mathcal I$ is connected,
  and (1b) or (1c) holds according as its solution set is subdirect.
\item $\mathcal B_{0}\neq\varnothing$: one constructs a family $\Omega$ of
  weakenings of $\mathcal I$ that is \emph{critically} unable to solve all of
  $\mathcal B_{0}$ (conditions 1--4 of
  Construction~\ref{constr:omega}) and argues by whether some member fails
  (1b).
\end{itemize}
\end{proof}

\begin{hazard}[Construction~\ref{constr:omega} is the hardest object]\label{haz:omega}
The set $\Omega$ is specified by four conditions, one of which (condition 3) is
a minimality condition ``if we replace any instance in $\Omega$ by all weaker
instances then 2 is not satisfied''. The paper obtains $\Omega$ by starting
from $\{\mathcal I\}$ and repeatedly weakening, justified by ``we cannot weaken
forever''. Condition 4 is then \emph{derived} from condition 3.

For a formalization this is the single most demanding step in
\S\ref{sec:main-statements}, for three reasons. (i) The iteration weakens a
\emph{set} of instances, so the termination measure is a measure on finite sets
of instances, not on instances --- see Hazard~\ref{haz:weakening}. (ii) The
derivation of condition 4 from condition 3 is a two-line argument that
instantiates condition 3 at
$\Omega\cup\{\mathcal J_{C}:C\in\mathcal J\}\setminus\{\mathcal J\}$ and reads
off an existential; the existential witness $B$ depends on $\mathcal J$, so
condition 4 is a $\Pi\Sigma$ statement and must be stated as such. (iii) The
subsequent Subcase~1 builds a \emph{second} descending sequence
$\mathcal B_{1}\supseteq\mathcal B_{2}\supseteq\dots$ inside the first, again
terminating by finiteness. Nested well-founded constructions of this shape are
where informal proofs most often hide a circularity, and this one should be
isolated and stated as a standalone lemma before Theorem~\ref{thm:main-inductive}
is attempted.
\end{hazard}

\begin{remark}[Construction of $\Omega$]\label{constr:omega}
$\Omega$ is a finite set of weakenings of $\mathcal I$ with:
\begin{enumerate}\tightlist
\item every member is a weakening of $\mathcal I$;
\item $\bigcap_{\mathcal J\in\Omega}\Sol(\mathcal J)=\varnothing$, where
  $\Sol(\mathcal J)=\{B\in\mathcal B_{0}:\mathcal J^{(B,\bot)}$ has a
  solution$\}$;
\item replacing any member by all weaker instances destroys 2;
\item for every $\mathcal J\in\Omega$ there is $B\in\mathcal B_{0}$ with
  (a) $\mathcal J$ crucial in $D^{(B,\bot)}$ and (b) $B\in\Sol(\mathcal J')$
  for every $\mathcal J'\in\Omega\setminus\{\mathcal J\}$.
\end{enumerate}
\end{remark}

\subsection{The three consumed theorems}\label{sec:codim-one}

\begin{theorem}[Reductions are safe]\label{thm:reductions-safe}
Let $\Theta$ be a cycle-consistent irreducible instance and
$B\subt[D_{x}]{T}D_{x}$ with $T\in\{\TBA,\TC,\TPC\}$. Then $\Theta$ has a
solution if and only if $\Theta$ has a solution with $x\in B$.
\end{theorem}

\begin{proof}
For $T=\TPC$ this is Theorem~\ref{thm:pc-safe}. For $T\in\{\TBA,\TC\}$,
Lemma~\ref{lem:find-consistent} gives a $1$-consistent reduction
$D^{(1)}\subte{T}D$ with $D^{(1)}_{x}\subseteq B$; by
Theorem~\ref{thm:main-inductive}(2), $\Theta^{(1)}$ has a solution, hence
$\Theta$ has one with $x\in B$.
\end{proof}

\begin{theorem}[PC reductions do not kill all solutions]\label{thm:pc-safe}
If $\mathcal I$ is cycle-consistent and irreducible, $B\subt[D_{y}]{\TPC(\sigma)}D_{y}$
for some $y\in\Var(\mathcal I)$, and $\mathcal I$ has a solution, then
$\mathcal I$ has a solution with $y\in B$.
\end{theorem}

\begin{longtheorem}[Codimension one]\label{thm:codim-one}
Suppose
\begin{enumerate}\romanlabels\tightlist
\item $\mathcal I$ is a linked cycle-consistent irreducible instance with
  $\Var(\mathcal I)=\{x_{1},\dots,x_{m}\}$;
\item $\mathbf D_{x_{i}}$ is $\TS$-free for every $i$;
\item weakening all constraints of $\mathcal I$ gives an instance with
  subdirect solution set;
\item $\sigma_{x_{i}}$ is the intersection of all linear congruences $\sigma$
  on $\mathbf D_{x_{i}}$ with $\cover{\sigma}=D_{x_{i}}^{2}$;
\item $L_{x_{i}}=\mathbf D_{x_{i}}/\sigma_{x_{i}}$;
\item $\phi\colon\mathbf Z_{q_{1}}\times\dots\times\mathbf Z_{q_{k}}\to
  L_{x_{1}}\times\dots\times L_{x_{m}}$ is a homomorphism, $q_{i}$ prime;
\item weakening any one constraint of $\mathcal I$ gives an instance with a
  solution in $\phi(\bar a)$ for every
  $\bar a\in\mathbf Z_{q_{1}}\times\dots\times\mathbf Z_{q_{k}}$.
\end{enumerate}
Then
$\Delta=\{\bar a:\mathcal I\text{ has a solution in }\phi(\bar a)\}$ is empty,
full, or an affine subspace of codimension $1$ --- the solution set of a single
linear equation.
\end{longtheorem}

\begin{proof}[Proof]
If $\Delta$ is full we are done. Otherwise pick $\bar b\notin\Delta$ and regard
$\phi(\bar b)$ as a reduction $D^{(1)}$ for $\mathcal I$; condition (vii) says
$\mathcal I$ is crucial in $D^{(1)}$.

\emph{Step 1: some $\ConOne(C,x)$ is a perfect linear congruence.} By
Lemma~\ref{lem:minimal-consistent}, either $C_{0}^{(1)}=\varnothing$ for some
constraint $C_{0}$, or $D^{(1)}$ is $1$-consistent.
\begin{itemize}\tightlist
\item If $C_{0}^{(1)}=\varnothing$: cruciality forces $\mathcal I=\{C_{0}\}$,
  say $C_{0}=R(y_{1},\dots,y_{t})$. By
  Lemma~\ref{lem:parallelogram-crucial}, $R$ has the parallelogram property and
  $\ConOne(R,1)$ is linear with $\cover{\ConOne(R,1)}=D_{y_{1}}^{2}$. By
  Lemma~\ref{lem:linear-on-top}, $\mathbf D_{y_{1}}/\delta\cong\mathbf Z_{p}$;
  with $\psi$ the quotient map, $\zeta=\{(a_{1},a_{2},b):\psi(a_{1})-\psi(a_{2})=b\}$
  witnesses perfection.
\item If $D^{(1)}$ is $1$-consistent: by Theorem~\ref{thm:main-inductive}(1)
  every constraint has the parallelogram property and (1b) or (1c) holds. In
  either case Lemma~\ref{lem:connected}(p) applies --- under (1c) to the
  linked connected subinstance, whose containing a constraint of $\mathcal I$
  is forced by condition (iii); under (1b) to $\mathcal I$ itself.
\end{itemize}

\emph{Step 2: add a $\mathbf Z_{p}$ coordinate.} Let $\zeta$ witness
perfection of $\ConOne(C,x)$. Add a variable $z$ with domain $\mathbf Z_{p}$,
replace $x$ in $C$ by a fresh $x'$, and add $\zeta(x,x',z)$, giving
$\mathcal I'$. Put
\[
  L=\{(\bar a,b):\mathcal I'\text{ has a solution with }z=b\text{ in }\phi(\bar a)\}
   \le\mathbf Z_{q_{1}}\times\dots\times\mathbf Z_{q_{k}}\times\mathbf Z_{p}.
\]
By condition (vii), $\proj_{1..k}(L)$ is full; and $(\bar b,0)\notin L$. So $L$
has dimension $k$ and is the solution set of one linear equation. If that
equation is $z=c$ with $c\neq 0$ then $\Delta=\varnothing$; otherwise setting
$z=0$ gives an equation defining $\Delta$, of dimension $k-1$.
\end{proof}

\begin{hazard}[The linear algebra of Step 2]\label{haz:step2-linear-algebra}
Step 2 uses four facts about subgroups of
$\mathbf Z_{q_{1}}\times\dots\times\mathbf Z_{q_{k}}\times\mathbf Z_{p}$ that
are stated as if obvious:
\begin{enumerate}\alphlabels\tightlist
\item $L$ is a subuniverse, hence (the $q_{i},p$ being prime and $w$ being
  a sum) a subgroup;
\item ``dimension'' is well defined for a subgroup of a product of
  $\mathbf Z_{q}$ with \emph{different} primes $q$. It is not a vector space;
  the correct statement is that such a group splits as a product over primes
  of $\mathbb F_{q}$-vector spaces, and dimension is the tuple of dimensions.
  The footnote in \cite{ZhukSimplified} about ``$J$ may contain only variables
  on the same field'' is the same issue surfacing in the algorithm;
\item a subgroup of codimension one is the kernel of a single homomorphism to
  $\mathbb F_{q}$, i.e.\ the solution set of one linear equation;
\item ``$\proj_{1..k}(L)$ full and $(\bar b,0)\notin L$ imply $\dim L=k$''.
\end{enumerate}
Item (b) is the one to be careful about: it is the reason
Theorem~\ref{thm:codim-one} is stated with mixed primes $q_{1},\dots,q_{k}$ and
the reason (c) must be proved prime by prime. In Mathlib the substrate exists
(\texttt{ZMod}, \texttt{Module} over \texttt{ZMod p}) but the mixed-prime
dimension notion and the codimension-one characterisation are not there; see \S\ref{sec:formalization}.
\end{hazard}

\subsection{The supporting lemmas}

The lemmas invoked above, in dependency order.

\begin{lemma}[{\cite[Lem.~6.1]{ZhukJACM}}]\label{lem:expanded-consistency}
An expanded covering of a cycle-consistent irreducible instance is
cycle-consistent and irreducible.
\end{lemma}

\begin{lemma}\label{lem:minimal-consistent}
Suppose $D^{(1)}$ is a $1$-consistent reduction for $\mathcal I$, every
$D^{(1)}_{x}$ is $\TS$-free, $T\in\{\TPC,\TL,\TD\}$, $D^{(1)}\lll D$, and
$D^{(2)}_{x}\subte{\MT T}D^{(1)}_{x}$ is a minimal $\MT T$ subuniverse for
every $x$. Then either $C^{(2)}=\varnothing$ for some constraint $C$, or
$\mathcal I^{(2)}$ is $1$-consistent.
\end{lemma}

\begin{lemma}\label{lem:crucial-irreducible}
A constraint crucial in a reduction of a cycle-consistent irreducible instance
is rectangular, and its $\ConOne$'s are irreducible congruences.
\end{lemma}

\begin{lemma}\label{lem:bridge-from-relation}
A relation whose restriction along strong subsets becomes empty yields a bridge
between the witnessing congruences.
\end{lemma}

\begin{lemma}[Connected instances]\label{lem:connected}
Let $\mathcal I$ be a cycle-consistent connected instance. Then
\begin{enumerate}\tightlist
\item[\textup{(a)}] any two constraints with a common variable are adjacent;
\item[\textup{(b)}] for constraints $C_{1},C_{2}$, variables
  $x_{i}\in\Var(C_{i})$, and any path from $x_{1}$ to $x_{2}$, there is a
  bridge $\delta$ from $\ConOne(C_{1},x_{1})$ to $\ConOne(C_{2},x_{2})$ with
  $\bcol{\delta}$ containing all pairs connected by that path;
\item[\textup{(p)}] if $\mathcal I$ is linked then $\ConOne(C,x)$ is a perfect
  linear congruence for every $C\in\mathcal I$, $x\in\Var(C)$.
\end{enumerate}
\end{lemma}

\begin{hazard}\label{haz:connected-p}
Clause (p) is the bridge between the bridge machinery and the linear algebra;
it is what makes Theorem~\ref{thm:codim-one} possible. Its proof is (b)
followed by Lemma~\ref{lem:perfect-from-linked}: linkedness of the instance
makes $\bcol{\delta}$ linked as a relation, which upgrades the bridge. The two
senses of ``linked'' (Hazard~\ref{haz:linked}) meet exactly here, and the step
where one becomes the other should be isolated.
\end{hazard}

\begin{lemma}[{\cite[Lem.~5.6]{ZhukStrong}}]\label{lem:tree-coverings}
For a $1$-consistent instance and a reduction, there is a tree-covering whose
projection onto a given variable computes the maximal $1$-consistent
subreduction at that variable.
\end{lemma}

\begin{lemma}\label{lem:find-consistent}
Let $D^{(1)}$ be a $1$-consistent reduction of a cycle-consistent instance
$\mathcal I$ with $D^{(1)}\lll D$, and $B\subt[D_{x}]{T}D^{(1)}_{x}$ with
$T\in\{\TBA,\TC,\TPC\}$. Then there is a nonempty $1$-consistent reduction
$D^{(2)}\lll D^{(1)}$ with $D^{(2)}_{x}\subseteq B$. Moreover, if
$T\in\{\TBA,\TC\}$ then $D^{(2)}\subte{T}D^{(1)}$; and if $T=\TPC$ and every
$D^{(1)}_{y}$ is $\TS$-free then $D^{(2)}\subte{\TMPC}D^{(1)}$.
\end{lemma}

\begin{lemma}\label{lem:parallelogram-crucial}
A constraint that becomes empty under a minimal multi-type reduction, in an
instance crucial there, has the parallelogram property; and in the linear case
its first induced congruence $\sigma$ satisfies $\cover\sigma=D^{2}$.
\end{lemma}

\begin{longlemma}[Bridges from a subdirect PC/linear instance]\label{lem:bridge-from-instance}
Suppose
\begin{enumerate}\romanlabels\tightlist
\item $\mathcal I$ has subdirect solution set;
\item $D^{(1)}$ is a reduction with $D^{(1)}_{x}\lll D_{x}$ for every $x$;
\item $C\in\mathcal I$ is a constraint of type $T\in\{\TPC,\TL\}$;
\item $B\subt[D_{z}]{\mathcal T(\xi)}D^{(1)}_{z}$ for some variable $z$, with
  $\mathcal T\in\{\TBA,\TC,\TS,\TPC,\TL\}$;
\item if $T=\TPC$ then $\mathcal T\in\{\TPC,\TL\}$;
\item $\mathcal I^{(1)}$ has a solution;
\item $\mathcal I^{(1)}$ has no solution with $z\in B$;
\item weakening $C$ in $\mathcal I$ gives an instance with a solution in
  $D^{(1)}$ having $z\in B$.
\end{enumerate}
Then $\mathcal T=T$, and for every variable $x$ of $C$ there is a bridge
$\delta$ from $\xi$ to $\ConOne(C,x)$ with
$\bcol{\delta}\supseteq\mathcal I(z,x)$.
\end{longlemma}

\begin{hazard}[The longest proof in the section]\label{haz:bridge-from-instance}
Lemma~\ref{lem:bridge-from-instance} carries eight numbered hypotheses and its
proof runs to roughly ten pages. It is the only source of bridges in
\S\ref{sec:main-statements}, and every use of connectedness routes through it.
Its conclusion has two independent parts --- ``$\mathcal T=T$'' and ``there is a
bridge with $\bcol\delta\supseteq\mathcal I(z,x)$'' --- and different call sites
use different parts; splitting it into two lemmas is tempting but the proof
produces them together.

A formalization should attack this lemma \emph{last} among the \S3.2 material
and \emph{first} among the things budgeted generously. It, together with
Theorem~\ref{thm:stable-intersection} (which it consumes) and
Theorem~\ref{thm:main-inductive} (which consumes it), is where the majority of
the effort will go.
\end{hazard}

\begin{corollary}\label{cor:same-type}
In the situation of Theorem~\ref{thm:main-inductive}(2) with $\mathcal I'$
satisfying (1b), the type $\mathcal T$ of the reduction and the type of the
constraints must agree, which contradicts $\mathcal T\in\{\TBA,\TC\}$.
\end{corollary}
\clearpage
%% ---------------------------------------------------------------- ch7-algorithm.tex
\section{The algorithm}\label{sec:algorithm}

\cite{ZhukSimplified} proves the correctness statements but does not contain the
algorithm; that is in \cite{ZhukJACM}. This section states it, because target
\textup{(T1)} of \S\ref{sec:target} is a proposition \emph{about} it and cannot
be written down otherwise.

\subsection{\textsc{Solve}}

\begin{verbatim}
function Solve(I):
    repeat
        I := ForceConsistency(I)
        if I = false: return "No"
        if some D_x has a strong subset B:
            I := ReduceDomain(I, x, B)
    until nothing changed
    return SolveLinear(I)
\end{verbatim}

\textsc{ForceConsistency} enforces cycle-consistency and irreducibility,
returning \texttt{false} if some domain empties. The reduction step is licensed
by Theorem~\ref{thm:reductions-safe}: restricting to a strong subset cannot
destroy solvability. When no strong subset remains, Lemma~\ref{lem:ubiquity}
together with Lemma~\ref{lem:linear-on-top} gives, for each domain of size at
least $2$, a congruence $\sigma_{x}$ with $D_{x}/\sigma_{x}\cong\mathbb Z_{q_{1}}
\times\dots\times\mathbb Z_{q_{n_{x}}}$; choose $\sigma_{x}$ minimal. That case
is \textsc{SolveLinear}.

\subsection{\textsc{SolveLinear}}

Write $\phi\colon\mathbb Z_{p_{1}}\times\dots\times\mathbb Z_{p_{k}}\to
D_{x_{1}}/\sigma_{x_{1}}\times\dots\times D_{x_{m}}/\sigma_{x_{m}}$ for a linear
map and
\[
  \phi^{-1}(\mathcal I)=\{\bar a:\mathcal I\text{ has a solution in }\phi(\bar a)\}.
\]
Computing $\phi^{-1}(\mathcal I)$ would solve $\mathcal I$; the algorithm cannot
do that, but it can do four things:

\begin{description}\tightlist
\item[\textup{(p0)}] decide $\bar a\in\phi^{-1}(\mathcal I)$, by reducing each
  domain to the corresponding block and recursing on a smaller domain;
\item[\textup{(p1)}] decide $\phi^{-1}(\mathcal I)=\mathbb Z_{p_{1}}\times\dots
  \times\mathbb Z_{p_{k}}$, by testing $\bar 0$ and the $k$ unit vectors with
  \textup{(p0)} --- correct because $\phi^{-1}(\mathcal I)$ is a subgroup;
\item[\textup{(p2)}] compute $\phi^{-1}(\mathcal I)$ when $\mathcal I$ is not
  linked, by splitting into linked pieces, recursing, and taking the union;
\item[\textup{(p3)}] compute $\phi^{-1}(\mathcal I)$ when it is empty or of
  codimension one, by finding $\bar a\notin\phi^{-1}(\mathcal I)$, then for each
  $i$ a $b_{i}$ with $\bar a[i\mapsto b_{i}]\in\phi^{-1}(\mathcal I)$ if one
  exists, and reading off the equation
  $\sum_{i\in J}(y_{i}-a_{i})/(b_{i}-a_{i})=1$ over the $i$ for which $b_{i}$
  exists.
\end{description}

\begin{verbatim}
function SolveLinear(I):
    m := dim(D_{x_1}/s_1 x ... x D_{x_n}/s_n)
    phi := a bijective linear map Z_{p_1}x...xZ_{p_m} -> that product
    while phi^-1(I) != Z_{p_1}x...xZ_{p_m}:              // (p1)
        I' := I
        for C in I':                                // drop what can be dropped
            if phi^-1(I' \ {C}) != Z_{p_1}x...xZ_{p_m}:   // (p1)
                I' := I' \ {C}
        F := phi^-1(I')                                // (p2) or (p3)
        if F = {}: return "No"
        m := dim F ;  psi := a bijective linear map onto F ;  phi := phi o psi
    return "Yes"
\end{verbatim}

The inner loop makes $\mathcal I'$ minimal with $\phi^{-1}(\mathcal I')$ not
full. At that point Theorem~\ref{thm:codim-one} applies --- its hypothesis (vii)
is exactly minimality --- so $\phi^{-1}(\mathcal I')$ is empty or of codimension
one, and \textup{(p3)} computes it; or $\mathcal I'$ is not linked and
\textup{(p2)} does. Since $\phi^{-1}(\mathcal I)\subseteq\phi^{-1}(\mathcal I')$,
the image of $\phi$ shrinks while still containing all solutions. The dimension
strictly decreases, so the loop terminates.

\subsection{What is honestly formalizable}

\begin{hazard}[Three obstacles between the pseudocode and a Lean function]
\label{haz:algorithm-formalization}
\begin{enumerate}\alphlabels\tightlist
\item \emph{Termination is not structural.} \textsc{Solve} recurses through
  \textup{(p0)} on strictly smaller domains, and \textsc{SolveLinear} loops on a
  strictly decreasing dimension. Neither is a structural recursion; both need an
  explicit well-founded measure, and the two interleave --- \textsc{SolveLinear}
  calls \textsc{Solve} on a smaller domain, which calls \textsc{SolveLinear}
  again. The measure is lexicographic:
  $\bigl(\sum_{x}|D_{x}|,\ \dim\phi\bigr)$.
\item \emph{Several steps quantify over objects that are only finitely many by
  accident.} ``If some $D_{x}$ has a strong subset $B$'' quantifies over subsets
  of $D_{x}$, over congruences on $D_{x}$, and --- through absorption --- over
  \emph{term operations}, of which there are infinitely many. Decidability comes
  from the fact that on a finite set there are finitely many operations of each
  arity and finitely many subsets, so "$B$ absorbs $A$ with some term'' is
  decidable by exhaustive search over $\Clo_{k}$ for $k\le|A|^{|A|}$; but that
  bound is a theorem, not a definition, and it is the only reason the algorithm
  is an algorithm.
\item \emph{Polynomial time is a different subject.} The bound in
  \cite{ZhukJACM} counts calls and bounds the recursion depth by $|A|$; making
  that formal needs a cost model. See \S\ref{sec:complexity-layer}.
\end{enumerate}
\end{hazard}

\begin{theorem}[Target \textup{(T1)}]\label{thm:t1}
There is a function $\textsc{Solve}$ from instances of $\CSP(\Gamma)$ to
Booleans such that for every instance $\mathcal I$,
\[
  \textsc{Solve}(\mathcal I)=\texttt{true}
  \iff \mathcal I\ \text{has a solution}.
\]
\end{theorem}

Theorem~\ref{thm:t1} mentions no machine, is provable from
\S\ref{sec:main-statements}, and is the honest formal content of ``the algorithm
works''. It is strictly stronger than decidability of $\CSP(\Gamma)$, which is
trivial for a finite template: it says that \emph{this} algorithm --- the one
whose existence is the dichotomy --- is correct.
\clearpage
%% ---------------------------------------------------------------- ch8-hardness.tex
\section{The hard half}\label{sec:hardness}

If no weak near-unanimity operation preserves $\Gamma$, then $\CSP(\Gamma)$ is
\textup{NP}-complete. This half is older, shorter, and better understood than
the tractable half, and it factors cleanly into algebra and complexity.

\subsection{The chain}

\begin{description}\tightlist
\item[\textup{(1)}] \emph{Reduction to a core, and to an idempotent algebra.}
  One may assume $\Gamma$ contains all singleton unary relations, so that
  $\Pol(\Gamma)$ is idempotent, at the cost of a polynomial-time reduction.
\item[\textup{(2)}] \emph{Maróti--McKenzie.} A finite idempotent algebra has a
  Taylor term if and only if it has a weak near-unanimity term
  \cite{MarotiMcKenzie}.
\item[\textup{(3)}] \emph{No Taylor term $\Rightarrow$ pp-interpretation.} If a
  finite idempotent algebra has no Taylor term, then the variety it generates
  contains a two-element algebra whose only operations are projections;
  equivalently, $\Gamma$ primitively positively interprets every finite
  constraint language, in particular a template for \textsc{$3$-Sat} or
  \textsc{Nae-$3$-Sat}.
\item[\textup{(4)}] \emph{Bulatov--Jeavons--Krokhin.} A primitive positive
  interpretation of $\Delta$ in $\Gamma$ yields a polynomial-time many-one
  reduction $\CSP(\Delta)\le_{p}\CSP(\Gamma)$ \cite{BJK}.
\item[\textup{(5)}] \emph{Cook--Levin.} $\textsc{Nae-$3$-Sat}$ is
  \textup{NP}-hard.
\end{description}

Steps (1)--(3) are algebra; step (4) is a syntactic construction plus a
complexity-theoretic wrapper; step (5) is pure complexity theory.

\subsection{The layered targets}

\begin{theorem}[Target \textup{(H0)}]\label{thm:h0}
Let $\mathbf A$ be a finite idempotent algebra with no weak near-unanimity term
operation. Then $\Inv(\mathbf A)$ primitively positively interprets the
constraint language $\{\textsc{Nae}\}$ on $\{0,1\}$, where
$\textsc{Nae}=\{0,1\}^{3}\setminus\{(0,0,0),(1,1,1)\}$.
\end{theorem}

\begin{theorem}[Target \textup{(H1)}]\label{thm:h1}
If $\Gamma$ primitively positively interprets $\Delta$ then
$\CSP(\Delta)\le_{p}\CSP(\Gamma)$. Consequently, under \textup{(H0)},
$\CSP(\Gamma)$ is \textup{NP}-hard.
\end{theorem}

Theorem~\ref{thm:h0} mentions no machine. Theorem~\ref{thm:h1} is where the
complexity layer enters, and it is small --- the reduction is a syntactic
substitution of formulas --- but it cannot be stated at all without a
formalized notion of polynomial-time many-one reduction.

\subsection{Definitions}

\begin{definition}[pp-formula, pp-definability]\label{def:pp}
A \emph{primitive positive formula} over $\Gamma$ is
$\exists y_{1}\dots\exists y_{s}\;\bigwedge_{j}R_{j}(\bar z_{j})$ with
$R_{j}\in\Gamma\cup\{=\}$. A relation is \emph{pp-definable} from $\Gamma$ if
some pp-formula defines it.
\end{definition}

\begin{definition}[pp-interpretation]\label{def:pp-interp}
$\Gamma$ on $A$ \emph{pp-interprets} $\Delta$ on $B$ if there are $d\ge 1$, a
pp-definable $F\subseteq A^{d}$, and a surjection $\pi\colon F\to B$ such that
the $\pi$-preimage of every relation of $\Delta$, and of the equality relation
on $B$, is pp-definable from $\Gamma$.
\end{definition}

\begin{theorem}[Geiger; Bodnarchuk--Kalu\v znin--Kotov--Romov]\label{thm:galois}
For a finite $A$, a relation is pp-definable from $\Gamma$ if and only if it is
preserved by every polymorphism of $\Gamma$.
\end{theorem}

\begin{hazard}[The Galois connection is the expensive import]\label{haz:galois}
Theorem~\ref{thm:galois} is what converts the algebraic hypothesis (``no WNU
polymorphism'') into a syntactic conclusion (``pp-defines a hard relation''),
and it is the only genuinely substantial import of \S\ref{sec:hardness}. The
easy direction is a one-line calculation. The hard direction --- every relation
preserved by all polymorphisms is pp-definable --- goes through the free algebra
on $|A|^{k}$ generators and an indicator-instance construction; it is roughly
four printed pages and is the single largest item in the \textup{(H0)} budget.
It is also of independent value and would be a reasonable Mathlib contribution
on its own.
\end{hazard}

\begin{remark}[Why the hard half is not the hard half]\label{rem:hardness-easier}
Despite its name, \textup{(H0)} is far smaller than \textup{(T0)}: perhaps
twenty printed pages against a hundred. The asymmetry is intrinsic. Showing a
problem is hard requires exhibiting one construction; showing it is easy
requires an algorithm that works on every instance, and the correctness proof
must survive every configuration the instance can present.
\end{remark}
\clearpage
%% ---------------------------------------------------------------- ch10-defects.tex
\section{Defects in the sources}\label{sec:defects}

A blueprint earns its keep by finding the places where the source cannot be
transcribed. This section is now a register rather than an argument: every item
has been resolved, and the reasoning lives beside the statement it repairs.
They are not criticisms of \cite{ZhukSimplified}, which is unusually careful
prose --- across nine thousand lines of \TeX{} it says ``obvious'' three times
and ``clearly'' never. They are the residue that any informal exposition
leaves.

Ten places where the source is consistent but leaves the reader to supply a
reading are legislated in \S\ref{sec:legislation}; two of the ten turned out to
be theorems rather than choices. Fourteen places where it is wrong, incomplete,
or states less than its own proof gives are listed below. All fourteen are
repaired, and none needed new mathematics: three needed a citation the later
paper had dropped, one was a typographical slip, one was a statement that is
false as written, and the rest were omitted steps.

\begin{center}\small
\begin{longtable}{@{}p{0.31\textwidth}p{0.60\textwidth}@{}}
\hline
\textbf{Where} & \textbf{What, and where the repair is} \\
\hline\endhead
Propagation modulo a congruence &
  Stated, used five times, never proved; it is
  Lemma~\ref{lem:propagation} at the canonical surjection.
  A mechanical audit of all $81$ labelled statements shows it is the
  \emph{only} statement in the source that is cited while being neither proved
  nor attributed. \\
Stable intersection, type $\TC$ &
  $n=2$ is argued only for $\TPC$; the argument has no analogue for $\TC$.
  Theorem~\ref{thm:stable-intersection} and Hazard~\ref{haz:d4}. \\
Bridges from relations &
  Applied without its third hypothesis, which can fail; cruciality supplies it.
  \S\ref{sec:main-statements}. \\
Main induction, (1c) &
  The endgame proves \emph{connected} where (1c) demands \emph{linked}, and
  Case~1 derives (1c) for the wrong reduction. Hazard~\ref{haz:case1c}. \\
Minimal $\MT T$ subuniverses &
  The bridging lemma was deleted from the source along with its citation.
  Hazard~\ref{haz:minimality}. \\
Symmetrisation of a bridge &
  The displayed formula defines a relation that is never a bridge.
  Hazard~\ref{haz:bridge-repairs}(a). \\
Bridges and abelian groups &
  \emph{False as stated}; a three-element counterexample.
  Hazard~\ref{haz:bridge-repairs}(b) and Remark~\ref{rem:phi-twist}. \\
Nontrivial reflexive bridges &
  Linkedness asserted from an inclusion bearing on another hypothesis.
  Hazard~\ref{haz:bridge-repairs}(c). \\
The PC inductive step &
  Cites the box form of absorption-preserves-linkedness for a restriction that
  is not a box. Lemma~\ref{lem:absorb-linked-sub},
  Hazard~\ref{haz:box-vs-sub}. \\
Existence of an irreducible congruence &
  Factors by the wrong congruence, hiding an unsupported properness step.
  Hazard~\ref{haz:main-existence}. \\
Linear congruences on top &
  The conclusion is asserted; proved here by a route that also removes an
  external import. Hazard~\ref{haz:on-top}. \\
Maximal multi-type extension &
  Concludes that two sets meet because each meets a third.
  Hazard~\ref{haz:maximal-mult}. \\
The central-relation citation &
  Drops a third alternative present in the source it cites, and needs an
  ``empty'' alternative under Convention~\ref{conv:empty}.
  Hazard~\ref{haz:central-relation}. \\
Full fibres from surjectivity &
  A non sequitur at three sites; true for a reason that makes the arity $|A|$
  load-bearing. Hazard~\ref{haz:full-fibre}. \\
\hline
\end{longtable}
\end{center}

Two further patterns are worth naming, because they recur rather than occurring
once. The first is a \emph{proper} containment asserted where the cited lemma
gives only $\subseteq$, usually after factoring by a congruence, which can
close the gap; it appears at least five times and is repaired the same way each
time, by producing the witness that the source has already displayed two lines
earlier. The second is an upgrade from ``the image has size one'' to ``the
image is \emph{this} block''; it appears four times and should become a named
lemma.

Finally, two statements are weaker than their own proofs and should be
restated: Lemma~\ref{lem:intersect-all}(i), whose proof gives
$B\cap D\dlll^{A}D$ (Hazard~\ref{haz:intersect-all-strength}), and
Lemma~\ref{lem:linked-implies-abs}, which needs the words \emph{proper
nonempty} (Hazard~\ref{haz:linked-abs-proper}).

\begin{remark}[A retracted item, and what replaced it]\label{rem:d10-retracted}
An earlier draft of this section carried an eleventh blocking item: that the
recursion-depth bound $|A|+|\Gamma|$ of \cite[Lem.~5.2]{ZhukJACM} was ``not
established for the \textsc{CheckTuple} $\to$ \textsc{SolveNonlinked} path'',
and that ``without the repair the algorithm is not polynomial''. That was
wrong, and it is withdrawn.

Lemma~5.2 bounds the depth by arguing that every recursion site except
\textsc{CheckWeakerInstance} ``reduce[s] all domains of size greater than 1''.
At the \textsc{SolveNonlinked} site this needs the instance to be neither
linked nor fragmented --- otherwise the linked-component congruence can be
trivial on some domain, that domain does not shrink, and the bound degrades
from the constant $|A|$ to $O(|A|\cdot n)$, which is fatal, since the total work
is $\mathrm{poly}(N)^{\text{depth}}$. The source does state the hypothesis and
does discharge it. It is simply not in one place:
\begin{itemize}\tightlist
\item the precondition ``an instance that is not linked and not fragmented''
  is in the prose introducing the function, and \emph{not} in the
  \texttt{Input:} line of its pseudocode;
\item it is discharged at the call site by ``$\Theta_{0}$ is not fragmented
  because it is crucial'' --- itself a one-line argument the source does not
  give: a fragmented unsatisfiable instance has an unsatisfiable part, and no
  constraint of the \emph{other} part is then crucial;
\item the step from ``not fragmented'' to ``\emph{every} domain splits'' ---
  propagate the partition along a path, which exists by non-fragmentedness and
  is value-preserving by $1$-consistency --- appears once, inside the proof of a
  \emph{different} lemma, and is never cited here.
\end{itemize}
So the mathematics is present and the theorem is not in doubt. What is absent is
any single place where the six recursion sites of Lemma~5.2 are discharged, and
that is a blueprint's job rather than a defect of the source. The same pattern
holds for the $\Gamma$-order argument, whose case~(3) needs the projections to
be unchanged under weakening, which holds because the instance is
$1$-consistent when \textsc{CheckWeakerInstance} runs --- true, and unstated.

\emph{The one step that really is missing} is supplied by
Lemma~\ref{lem:solvenonlinked-callsite} below. It is two lines, it is not in
the source, and without it the $|A|$ bound at this site is unjustified.
\end{remark}

\begin{lemma}[The \textsc{SolveNonlinked} call site]\label{lem:solvenonlinked-callsite}
Let $\Theta_{0}$ be cycle-consistent, non-fragmented and not linked, and let
$\Theta_{0}'$ be $\Theta_{0}$ with each domain restricted to a subset
$D_{x}'\subseteq D_{x}$. Then every domain of $\Theta_{0}$ of size greater than
one is strictly shrunk by the time \textsc{SolveNonlinked}$(\Theta_{0}')$
recurses into \textsc{Solve}.
\end{lemma}

\begin{proof}
Write $\tau_{x}=\mathrm{LinkedCon}(\Theta_{0},x)$, a congruence by
\cite[Lem.~6.2]{ZhukJACM}. First, $\tau_{x}$ is proper at \emph{every} variable.
Since $\Theta_{0}$ is not linked there are at least two connected components of
the graph on $\bigsqcup_{z}\{z\}\times D_{z}$; and every component meets every
variable, because from any $(x,a)$ non-fragmentedness supplies a path to $y$ and
$1$-consistency extends $a$ along it. So each $D_{y}$ is partitioned into at
least two nonempty $\tau_{y}$-blocks, each therefore proper.

Now split on whether $\Theta_{0}'$ is linked.

\emph{Not linked.} Then \textsc{SolveNonlinked} recurses on a component, whose
trace on each $D_{y}$ is a $\tau_{y}$-block intersected with $D_{y}'$ --- a
proper subset of $D_{y}$ by the previous paragraph.

\emph{Linked.} Then for every $z$ and all $a,b\in D_{z}'$ some path of
$\Theta_{0}'$ connects them; every such path is a path of $\Theta_{0}$, so
$D_{z}'\times D_{z}'\subseteq\tau_{z}$, i.e.\ $D_{z}'$ lies inside a single
$\tau_{z}$-block. That block is proper, so $D_{z}'\subsetneq D_{z}$ --- every
domain has \emph{already} been shrunk, at the \textsc{CheckTuple} step, before
\textsc{SolveNonlinked} was called at all.
\end{proof}

\begin{remark}[Why this is the right argument]\label{rem:callsite-moral}
Two things are worth extracting.

First, the question this replaces was the wrong one. The natural worry is
whether the restriction preserves \textsc{SolveNonlinked}'s stated precondition
``not linked'', and one then starts asking how the minimal linear congruence
$\mathrm{ConLin}$ relates to $\mathrm{LinkedCon}$ --- a relation the source
never states, and which is not needed. The dichotomy above uses no relation
between them: \emph{either} the precondition survives and the component split
does the shrinking, \emph{or} it fails and the restriction has already done it.

Second, it follows that \textsc{SolveNonlinked}'s precondition is not required
for the depth bound at all. The function is correct on a linked instance too ---
there is then one component and it recurses on the whole of it --- and
Lemma~\ref{lem:solvenonlinked-callsite} covers that case. The precondition
should be read as documentation of the intended use, not as an obligation the
caller must discharge.
\end{remark}

\begin{remark}[What this list is evidence of]\label{rem:defects-meaning}
Nine blocking items in a fifty-page paper is not a high rate; it is roughly what
one should expect, and it is the reason blueprints exist. Note the
distribution: one is a genuine misstatement (D1), one is an omission (D2), two
are gaps in the mathematics that need new arguments (D3, D4), three are
hypotheses that are too weak or missing (D5, D6, D8), and two are artefacts of
the interface with the older paper (D7, D9). None of them threatens the
theorem. All of them cost time.

Remark~\ref{rem:d10-retracted} is the counterweight, and the more instructive
entry. A tenth item was listed, survived a review round, and was published ---
and was wrong, because it was taken from a reader who worked from the pseudocode
while the hypothesis it needed was in the surrounding prose. A list like this
one is exactly as reliable as the reading behind each entry, and the failure
mode is not carelessness but reading the formal-looking part of a paper and not
the sentence before it.
\end{remark}
\clearpage
%% ---------------------------------------------------------------- ch9-formalization.tex
\section{The complexity layer}\label{sec:complexity-layer}

Targets \textup{(T2)} and \textup{(H1)} are stated here and not proved. The
point of stating them is to be precise about what is \emph{not} being claimed.

\begin{definition}[What a cost model would have to supply]\label{def:cost-model}
To state \textup{(T2)} one needs: an encoding of instances of $\CSP(\Gamma)$ as
strings; a machine model with a step count; a predicate ``$f$ is computed in
time bounded by a polynomial in the input length''; and closure of that
predicate under composition. To state \textup{(H1)} one needs, in addition, a
polynomial-time many-one reduction relation $\le_{p}$ and the class
\textup{NP}.
\end{definition}

Mathlib supplies the machine (\texttt{Turing.TM2}), an encoding
(\texttt{Computability.Encoding}), and a time-bounded computation predicate
(\texttt{Turing.TM2ComputableInPolyTime}). It does not supply a complexity
\emph{class}, \textup{P}, \textup{NP}, \textup{NP}-completeness, $\le_{p}$, or
\textsc{Sat}; and the composition lemma
\texttt{TM2ComputableInPolyTime.comp} is an open \texttt{proof\_wanted} in
Mathlib itself. Moreover \texttt{TM2ComputableInPolyTime} applies to total
functions $\alpha\to\beta$, never to languages, so even the existing predicate
would need reshaping.

\begin{remark}[This layer now exists, and it was cheap]\label{rem:complexity-scope}
An earlier draft of this remark said that building the complexity layer was
``a substantial project --- comparable in size to the algebraic core'' and that
attempting it as part of this project ``would be a scoping error''. That was
wrong, and it has been tested rather than argued: the vocabulary --- words, a
cost model, $\mathrm P$, $\mathrm{NP}$ by verifiers, $\le_{p}$,
\textup{NP}-hardness, \textup{NP}-completeness, and a generic
\textup{NP}-hard problem --- is about $850$ lines of Lean, \texttt{sorry}-free,
and took an afternoon.

That figure is what the literature should have led one to expect. The class
layer is $702$ lines in Balbach's Isabelle Cook--Levin development and $992$ in
G\"aher--Kunze's Coq one, against totals of $56\,514$ and $\approx 16\,500$
respectively. Defining the classes has never been the expensive part; the
expensive part is connecting them to a concrete machine, which is a different
activity.

So the reason to keep \textup{(T2)} out of scope is \emph{not} that complexity
theory is unavailable. It is Hazard~\ref{haz:t2-is-not-small}: \textup{(T2)}
requires exhibiting $\textsc{Solve}$ as a program with a proved step bound, and
that is the largest program-construction task in the project. \textup{(H1)},
by contrast, needs only a syntactic substitution and should be built.
\end{remark}

\begin{remark}[What the exercise cost, and what it caught]
Two defects surfaced in those $850$ lines, both in \emph{statements} rather
than proofs, and neither of a kind a type-checker reports. The generic problem
was first defined with a single budget bounding both the certificate length and
the running time, which makes the theorem false --- the two directions pull the
bound opposite ways. And the language was first given only the tests ``register
is empty'' and ``register's first bit is $1$'', with which no loop that consumes
its input can be written at all; that was found by trying to write the first
program.

The ratio is worth noting for the rest of this document. Two statement-level
errors in $850$ lines, against ten found in Zhuk's fifty pages by readers who
did not formalize (\S\ref{sec:defects}). The defect list should be expected to
grow under formalization, not to be complete.
\end{remark}

\section{Formalization}\label{sec:formalization}

\subsection{What is reused}\label{sec:reuse}

The predecessor development \texttt{ZhukLean} is a complete, \texttt{sorry}-free
Lean~4 formalization of Zhuk's centre theorem, built on
\texttt{Mathlib.ModelTheory}. It is consumed here as a \texttt{lake} dependency,
not copied. It supplies:

\begin{itemize}\tightlist
\item products of structures (\texttt{piStructure}, \texttt{prodStructure},
  coordinatewise realization, reindexing and evaluation homomorphisms) --- the
  one part of the universal-algebra vocabulary Mathlib lacks;
\item absorption in the tuple form (\texttt{Witnesses}, \texttt{Absorbs},
  \texttt{BinAbsorbs}), Taylor terms (\texttt{IsTaylorOn}), and central
  absorption (\texttt{CentrallyAbsorbs});
\item the relational description of absorption by essential relations
  (Barto--Kazda), which is what converts absorption at some arity into
  absorption at arity~$3$;
\item Zhuk's centre theorem itself (\texttt{zhuk\_center}) and the ternary
  collapse (\texttt{exists\_ternary\_witnesses}) --- which is precisely
  Lemma~\ref{lem:central-ternary}, cited without proof in
  \cite{ZhukSimplified}.
\end{itemize}

Three of \cite{ZhukSimplified}'s sixteen black-box imports are therefore already
discharged.

\subsection{What is built here}\label{sec:built}

\begin{center}
\begin{tabular}{@{}llr@{}}
\hline
Module & Contents & Status \\
\hline
\texttt{CSP/Product} & products of structures & complete \\
\texttt{CSP/Congruence} & \texttt{Congruence L M}, order, \texttt{sInf}, quotient & complete \\
\texttt{CSP/Irreducible} & irreducibility, existence and uniqueness of $\cover\sigma$ & complete \\
\texttt{CSP/Bridge} & bridges, collapse, composition, linear/PC & 2 imports open \\
\texttt{CSP/Types} & the six types, multi-types, $\lll$ & complete \\
\texttt{CSP/Instance} & instances, reductions, consistency, cruciality & complete \\
\texttt{CSP/Main} & the main statements & targets \\
\hline
\end{tabular}
\end{center}

Everything below \texttt{CSP/Main} is \texttt{sorry}-free. The two open items in
\texttt{CSP/Bridge} are Lemma~\ref{lem:bridge-comp} and its collapse identity,
both imported from \cite{ZhukJACM}.

\subsection{Design decisions}\label{sec:design}

\begin{description}
\item[Build on \texttt{Mathlib.ModelTheory}, keep the language generic.]
  A bespoke ``finite algebra with one WNU operation'' type would bundle the
  standing hypotheses as fields and avoid threading them, which is a real cost
  in the predecessor development. But it would discard
  \texttt{Substructures.lean} --- 985 lines of exactly the $\Sg$ API needed,
  including \texttt{mem\_closure\_iff\_exists\_term} with the variable type
  equal to the generating set, an interface no hand-rolled clone matches --- and
  it would discard the predecessor's 1600 lines wholesale. Congruences and
  quotients are also \emph{cheaper} in \texttt{ModelTheory}, not dearer.
  Instantiating to $\Vn$ is a thin layer on top.

\item[$\lll$ is \texttt{Relation.ReflTransGen}.]
  Its induction principle is the one the proofs use.

\item[Dotted relations are defined as disjunctions.]
  \texttt{DotLll C B := C = }$\varnothing$\texttt{ }$\vee$\texttt{ Lll C B}, so the case split is forced at every
  use site rather than resolved by the reader.

\item[The bridge criterion is the definition of linearity.]
  Zhuk defines a linear congruence by a per-block isomorphism to
  $\mathbb Z_{p}^{m}$ and derives the bridge criterion. We reverse this: the
  criterion is first-order over finite data and is what every consumer uses; the
  block description becomes a structure theorem. See
  Hazard~\ref{haz:linear-def}.

\item[$\cover\sigma$ is data attached to irreducibility.]
  \texttt{Irreducible.exists\_isCover} produces it; there is no standalone
  function, because there is no cover without irreducibility. A pleasant
  by-product: with the definition taken literally, irreducibility \emph{implies}
  $\sigma\neq A^{2}$, since the empty family intersects to the universe. The
  source never says this.

\item[Instances are lists; scopes are tuples; the variable type is a parameter.]
  See Hazards~\ref{haz:instance-list} and~\ref{haz:renaming}. The last of these
  is the decision that has not yet been cashed out, and it should be made before
  anything about coverings is written.

\item[Nothing is stubbed with \texttt{sorry} at the level of a definition.]
  \texttt{IsConnectedInstance} and \texttt{IsExpandedCovering} are
  \emph{absent} rather than defined as \texttt{sorry}, because a
  \texttt{def~:~Prop~:=~sorry} is a junk constant that downstream proofs can
  silently lean on. Absence forces the design decision; a stub hides it.
\end{description}

\subsection{Order of attack}\label{sec:order}

\begin{enumerate}\tightlist
\item Lemma~\ref{lem:ba-center-implies}. Small, self-contained, used everywhere,
  and adjacent to what is already formalized.
\item The mixed-prime linear algebra of Hazard~\ref{haz:step2-linear-algebra}.
  Independent of everything else; can be done in parallel.
\item Lemma~\ref{lem:bridge-comp}. Twelve source lines of relational algebra;
  closes the two open items in \texttt{CSP/Bridge}.
\item The coproduct design for coverings, then
  Proposition~\ref{prop:covering-basics}.
\item Theorem~\ref{thm:stable-intersection}. The one hard statement in the
  interface, and the sole source of bridges.
\item Lemma~\ref{lem:bridge-from-instance}. Eight hypotheses, ten pages.
\item Theorem~\ref{thm:main-inductive}. Last.
\end{enumerate}

\subsection{Scope}\label{sec:scope}

The calibration point is that the predecessor formalized about $3.5$ printed
pages in about $1670$ lines of Lean --- roughly $480$ lines per page. Excluding
\S\ref{sec:complexity-layer} and \cite{ZhukSimplified}'s \S4, which is an
independent second result and not needed for the dichotomy, the transitive
closure of what must be proved is on the order of $100$--$130$ printed pages,
so $35\,000$--$60\,000$ lines of Lean.

Line count, however, is not the binding constraint, and an earlier draft of this
section was wrong to convert it into person-years at conventional rates. The
pipeline this project belongs to has since produced roughly $33\,000$ verified,
\texttt{sorry}-free Lean lines in about ten active hours on the
Jordan--Sch\"onflies formalization, by fanning independent modules into separate
git worktrees and landing them in waves; the rate held on the hard content, not
merely on the foundation.

What binds is instead the \emph{critical path}: the depth of the serial chain,
and the number of places where the source is wrong and new mathematics is
needed. Here the serial chain is
Lemma~\ref{lem:bridge-from-instance} $\to$
Theorem~\ref{thm:stable-intersection} $\to$
Theorem~\ref{thm:main-inductive} --- each a single large proof that cannot be
split across workers --- and the new mathematics is the four blocking items
D3, D4, D6 and D7 of \S\ref{sec:defects}. Those are not throughput-limited.
Schedule around them, not around the line count.

What is finished here is the design: the route, the statements, the vocabulary
at formalization granularity, the hazard list, the defect list, the
missing-layer inventory, and a building development whose foundations are proved
and whose summit is stated.

\clearpage
\begin{thebibliography}{99}

\bibitem{BJK} A.~Bulatov, P.~Jeavons, A.~Krokhin,
  \emph{Classifying the complexity of constraints using finite algebras},
  SIAM J. Comput.\ \textbf{34} (2005), 720--742.

\bibitem{BradyNotes} Z.~Brady, \emph{Notes on CSPs and polymorphisms},
  arXiv:2210.07383.

\bibitem{BartoKozik} L.~Barto, M.~Kozik, \emph{Absorbing subalgebras, cyclic
  terms, and the constraint satisfaction problem}, Log. Methods Comput. Sci.
  \textbf{8} (2012), 1:07.

\bibitem{DecidingAbsorption} L.~Barto, A.~Kazda, \emph{Deciding absorption},
  Internat. J. Algebra Comput. \textbf{26} (2016), 1033--1060.

\bibitem{Bulatov} A.~Bulatov, \emph{A dichotomy theorem for nonuniform CSPs},
  arXiv:1703.03021; FOCS 2017.

\bibitem{FederVardi} T.~Feder, M.~Y.~Vardi, \emph{The computational structure of
  monotone monadic SNP and constraint satisfaction}, SIAM J. Comput.\
  \textbf{28} (1999), 57--104.

\bibitem{HobbyMcKenzie} D.~Hobby, R.~McKenzie, \emph{The structure of finite
  algebras}, Contemp. Math.~76, Amer. Math. Soc., 1988.

\bibitem{IstingerKaiser} M.~Istinger, H.~K.~Kaiser, \emph{A characterization of
  polynomially complete algebras}, J. Algebra \textbf{56} (1979), 103--110.

\bibitem{MarotiMcKenzie} M.~Mar\'oti, R.~McKenzie, \emph{Existence theorems for
  weakly symmetric operations}, Algebra Universalis \textbf{59} (2008),
  463--489.

\bibitem{MinimalTaylor} L.~Barto, Z.~Brady, A.~Bulatov, M.~Kozik, D.~Zhuk,
  \emph{Minimal Taylor algebras as a common framework for the three algebraic
  approaches to the CSP}, LICS 2021.

\bibitem{RossSlides} R.~Willard, slides on bridges, similarity and
  centralizers.

\bibitem{ZhukJACM} D.~Zhuk, \emph{A proof of the CSP dichotomy conjecture},
  J. ACM \textbf{67} (2020), no.~5, 1--78; arXiv:1704.01914.

\bibitem{ZhukSimplified} D.~Zhuk, \emph{A simplified proof of the CSP dichotomy
  conjecture and XY-symmetric operations}, arXiv:2404.01080v2.

\bibitem{ZhukStrong} D.~Zhuk, \emph{Strong subalgebras and the constraint
  satisfaction problem}, J. Mult.-Valued Logic Soft Comput.\ \textbf{36} (2021);
  arXiv:2005.00593.

\end{thebibliography}

\end{document}
